% La cybercriminalité 2 : le harcèlement numérique — Allemand Tle A — Cours signé OnBuch+
% Programme officiel MINESEC. Compiler avec : tectonic AL18-la-cybercriminalite-2-le-harcelement-numerique.tex
\newcommand{\DOCMATIERE}{Allemand}
\newcommand{\DOCNIVEAU}{Tle A}
\newcommand{\DOCMODULE}{Chapitre 4 : Média et communication}
\newcommand{\DOCLECON}{AL18}
\newcommand{\DOCTITRE}{La cybercriminalité 2 : le harcèlement numérique}
% OnBuch+ — préambule « cours signés » ENRICHI (compilé avec tectonic / XeLaTeX)
% Système visuel « Pop » d'OnBuch+ : fond crème (jamais blanc), bordures encre
% épaisses, ombres « dures » décalées, titres Archivo Black en capitales.
% Version MATHÉMATIQUES : graphes (pgfplots/tikz), boîtes pédagogiques
% (définition, propriété, méthode, attention, le savais-tu, exercice résolu,
% exercice, TP, bilan), page de garde et signature OnBuch+ ancrée en bas.
%
% Macros attendues dans main.tex AVANT \input{preamble} :
%   \DOCMATIERE  \DOCNIVEAU  \DOCMODULE  \DOCLECON (numéro)  \DOCTITRE
\documentclass[11pt]{article}

\usepackage{fontspec}
\usepackage[ngerman,french]{babel} % français = langue principale ; l'allemand via \all{..} / textebox
\usepackage[a4paper,margin=2cm,top=3.4cm,bottom=2.8cm,headheight=1.4cm,footskip=1.3cm]{geometry}
\usepackage{amsmath,amssymb,amsfonts}
\usepackage{mathtools} % \xmapsto, \coloneqq... souvent utilisés par les modèles
\usepackage{cancel} % simplifications barrées (bilans de forces)
% \abs{z} (module, valeur absolue) et \norm{u} (norme), utilisés par les modèles
\providecommand{\abs}{}\let\abs\relax\DeclarePairedDelimiter{\abs}{\lvert}{\rvert}
\providecommand{\norm}{}\let\norm\relax\DeclarePairedDelimiter{\norm}{\lVert}{\rVert}
\usepackage[table]{xcolor} % [table] : fournit aussi \rowcolors (alternance de lignes)
\usepackage{pagecolor}
\usepackage{tikz}
\usetikzlibrary{arrows.meta,positioning,calc,shapes.geometric,shapes.misc,shapes.arrows,decorations.pathmorphing,decorations.pathreplacing,decorations.markings,patterns,shadows,mindmap,trees,fit,backgrounds,matrix,intersections,angles,quotes}
\usepackage{pgfplots}
\usepackage[version=4]{mhchem} % équations nucléaires \ce{^{235}_{92}U}
\usepackage[european,siunitx]{circuitikz} % schémas électriques
\pgfplotsset{compat=1.18}
\usepgfplotslibrary{fillbetween,groupplots} % aires entre courbes (calcul intégral)
\usepackage{enumitem}
\usepackage{fancyhdr}
% [most] charge toutes les bibliothèques tcolorbox (listings, minted, raster,
% documentation...) et gonfle inutilement la mémoire TeX sur un document long
% (~300 boîtes) : on ne charge que ce qui est réellement utilisé.
\usepackage[skins,breakable]{tcolorbox}
\usepackage{graphicx}
\usepackage{colortbl}
\usepackage{booktabs}
\usepackage{tabularx}
\usepackage{diagbox} % case d'en-tête diagonale des tableaux à double entrée
% Colonnes extensibles centrée / gauche / droite (C, L, R), souvent utilisées par les modèles
\newcolumntype{C}{>{\centering\arraybackslash}X}
\newcolumntype{L}{>{\raggedright\arraybackslash}X}
\newcolumntype{R}{>{\raggedleft\arraybackslash}X}
\usepackage{array}
\usepackage{multicol}
\usepackage{multirow}
\usepackage{siunitx}
% parse-numbers=false : accepte aussi \SI{2,5 \times 10^{-3}}{...}
\sisetup{locale=FR,output-decimal-marker={,},group-minimum-digits=4,parse-numbers=false}
\DeclareSIUnit\ohm{\ensuremath{\symup{\Omega}}} % U+2126 absent de la police
\DeclareSIUnit\division{div} % réglages d'oscilloscope (ms/div, V/div)
\DeclareSIUnit\tr{tr} % tours (vitesse de rotation, ex. \SI{1500}{\tr\per\minute})
\DeclareSIUnit\molar{M} % concentration molaire (ex. \SI{3}{\molar}, \SI{1}{\micro\molar})
\DeclareSIUnit\an{an} % année (le modèle confond parfois avec \year, un compteur TeX) — ex. \SI{5}{\kilo\meter\per\an}
\DeclareSIUnit\FCFA{FCFA} % franc CFA (ex. \SI{500}{\FCFA\per\kilo\gram})
\DeclareSIUnit\degreeBrix{\textdegree Brix} % teneur en sucre d'un sirop/jus (réfractométrie)
\DeclareSIUnit\month{mois} % le modèle utilise parfois \month, un compteur TeX, comme unité de durée
\usepackage{qrcode}
% \degree : commande fréquemment utilisée par les modèles pour un angle ou une
% température en dehors de \SI{}{} (non fournie par les paquets chargés).
\providecommand{\degree}{^{\circ}}
% \qed (fin de démonstration), utilisé par les modèles sans amsthm
\providecommand{\qed}{\hfill$\square$}
% \barre : double barre des valeurs interdites dans un tableau de variations (tkz-tab)
\providecommand{\barre}{\ensuremath{\|}}
% environnement proof (amsthm non chargé) utilisé par les modèles
\ifdefined\proof\else\newenvironment{proof}[1][Démonstration]{\par\noindent\textit{#1.}\ }{\qed\par}\fi
\usepackage{lastpage}
\usepackage{hyperref}
\hypersetup{hidelinks}

% ============================================================
% Palette « Pop » OnBuch+ — mêmes hex que la classe P (plus_ui.dart)
% ============================================================
\definecolor{poporange}{HTML}{F59321}
\definecolor{popdark}{HTML}{E07A0C}
\definecolor{popcream}{HTML}{FFF6E8}
\definecolor{popink}{HTML}{1C1714}
\definecolor{poppurple}{HTML}{7A5AE0}
\definecolor{popgreen}{HTML}{2E9E6A}
\definecolor{poppink}{HTML}{E0457B}
\definecolor{popblue}{HTML}{2F7DE1}
\definecolor{popgold}{HTML}{C98A12}
\definecolor{popmuted}{HTML}{8A7F76}
\definecolor{popline}{HTML}{EADFD2}
\definecolor{popgray}{HTML}{8A7F76}
\colorlet{popgrey}{popgray}
% Alias tolérants (le modèle nomme parfois les couleurs différemment)
\colorlet{popwhite}{white}
\colorlet{popblack}{popink}
\colorlet{popred}{poppink}
\colorlet{popyellow}{popgold}
% Teintes claires (fonds de tableaux/encadrés) — le modèle nomme parfois
% les couleurs avec un suffixe L (ex. popblueL) sans les avoir définies.
\colorlet{poporangeL}{poporange!20}
\colorlet{popdarkL}{popdark!20}
\colorlet{poppurpleL}{poppurple!20}
\colorlet{popgreenL}{popgreen!20}
\colorlet{poppinkL}{poppink!20}
\colorlet{popblueL}{popblue!20}
\colorlet{popgoldL}{popgold!20}
\colorlet{popredL}{popred!20}
\colorlet{popgrayL}{popgray!20}
% Teintes claires pour fonds de tableaux / zones
\colorlet{popblueL}{popblue!10!white}
\colorlet{poporangeL}{poporange!14!white}
\colorlet{popgreenL}{popgreen!12!white}
\colorlet{poppurpleL}{poppurple!12!white}
\colorlet{poppinkL}{poppink!12!white}
\colorlet{popgoldL}{popgold!14!white}

\pagecolor{popcream}

% ============================================================
% Typographie
% ============================================================
\defaultfontfeatures{Ligatures=TeX}
\setmainfont{PlusJakartaSans-Medium.ttf}[Path=../../../fonts/,BoldFont=PlusJakartaSans-Bold.ttf]
% Maths en sans-serif (Fira Math) avec chiffres et lettres droites de Plus
% Jakarta, pour que les formules se fondent dans le texte.
\usepackage[math-style=ISO,bold-style=ISO]{unicode-math}
\setmathfont{FiraMath-Regular.otf}[Path=../../../fonts/]
\setmathfont{PlusJakartaSans-Medium.ttf}[Path=../../../fonts/,range={up/num,up/{Latin,latin},\mathup}]
\usepackage{icomma} % virgule décimale sans espace en mode math
% Caractères mathématiques Unicode absents des polices de texte (Plus Jakarta,
% Archivo Black) : rendus par leur équivalent mathématique, en texte comme en
% formule — sinon ils s'affichent en carré vide (ex. « Arithmétique dans ℤ »).
\usepackage{newunicodechar}
\newunicodechar{ℝ}{\ensuremath{\mathbb{R}}}
\newunicodechar{ℤ}{\ensuremath{\mathbb{Z}}}
\newunicodechar{ℂ}{\ensuremath{\mathbb{C}}}
\newunicodechar{ℕ}{\ensuremath{\mathbb{N}}}
\newunicodechar{ℚ}{\ensuremath{\mathbb{Q}}}
\newunicodechar{𝔻}{\ensuremath{\mathbb{D}}}
\newunicodechar{α}{\ensuremath{\alpha}}
\newunicodechar{β}{\ensuremath{\beta}}
\newunicodechar{γ}{\ensuremath{\gamma}}
\newunicodechar{δ}{\ensuremath{\delta}}
\newunicodechar{ε}{\ensuremath{\varepsilon}}
\newunicodechar{θ}{\ensuremath{\theta}}
\newunicodechar{λ}{\ensuremath{\lambda}}
\newunicodechar{μ}{\ensuremath{\mu}}
\newunicodechar{σ}{\ensuremath{\sigma}}
\newunicodechar{τ}{\ensuremath{\tau}}
\newunicodechar{φ}{\ensuremath{\varphi}}
\newunicodechar{ψ}{\ensuremath{\psi}}
\newunicodechar{ω}{\ensuremath{\omega}}
\newunicodechar{Σ}{\ensuremath{\Sigma}}
% majuscules produites par \MakeUppercase dans les titres (α-aminés -> Α)
\newunicodechar{Α}{\ensuremath{\alpha}}
\newunicodechar{Β}{\ensuremath{\beta}}
\newunicodechar{Γ}{\ensuremath{\Gamma}}
\newunicodechar{Δ}{\ensuremath{\Delta}}
\newunicodechar{Ε}{\ensuremath{\varepsilon}}
\newunicodechar{Θ}{\ensuremath{\Theta}}
\newunicodechar{Λ}{\ensuremath{\Lambda}}
\newunicodechar{Μ}{\ensuremath{\mu}}
\newunicodechar{Τ}{\ensuremath{\tau}}
\newunicodechar{Φ}{\ensuremath{\Phi}}
\newunicodechar{Ψ}{\ensuremath{\Psi}}
\newunicodechar{Ω}{\ensuremath{\Omega}}
\newunicodechar{↦}{\ensuremath{\mapsto}}
\newunicodechar{∈}{\ensuremath{\in}}
\newunicodechar{∉}{\ensuremath{\notin}}
\newunicodechar{∧}{\ensuremath{\wedge}}
\newunicodechar{⇔}{\ensuremath{\Leftrightarrow}}
\newunicodechar{⟺}{\ensuremath{\Longleftrightarrow}}
\newunicodechar{⇒}{\ensuremath{\Rightarrow}}
\newunicodechar{⟹}{\ensuremath{\Longrightarrow}}
\newunicodechar{∘}{\ensuremath{\circ}}
\newunicodechar{∪}{\ensuremath{\cup}}
\newunicodechar{∩}{\ensuremath{\cap}}
\newunicodechar{≡}{\ensuremath{\equiv}}
\newunicodechar{⊂}{\ensuremath{\subset}}
\newunicodechar{⊥}{\ensuremath{\perp}}
\newunicodechar{∃}{\ensuremath{\exists}}
\newunicodechar{∀}{\ensuremath{\forall}}
\newunicodechar{∛}{\ensuremath{\sqrt[3]{\,}}}
\newunicodechar{∜}{\ensuremath{\sqrt[4]{\,}}}
\newunicodechar{⃗}{\ensuremath{{}^{\to}}}
\newunicodechar{ⁿ}{\ifmmode^{n}\else\textsuperscript{\ensuremath{n}}\fi}
\newunicodechar{ˣ}{\ifmmode^{x}\else\textsuperscript{\ensuremath{x}}\fi}
\newunicodechar{ᵇ}{\ifmmode^{b}\else\textsuperscript{\ensuremath{b}}\fi}
\newunicodechar{ʳ}{\ifmmode^{r}\else\textsuperscript{\ensuremath{r}}\fi}
\newunicodechar{ᵘ}{\ifmmode^{u}\else\textsuperscript{\ensuremath{u}}\fi}
\newunicodechar{ʸ}{\ifmmode^{y}\else\textsuperscript{\ensuremath{y}}\fi}
\newunicodechar{ᵃ}{\ifmmode^{a}\else\textsuperscript{\ensuremath{a}}\fi}
\newunicodechar{ᵉ}{\ifmmode^{e}\else\textsuperscript{\ensuremath{e}}\fi}
\newunicodechar{ᵏ}{\ifmmode^{k}\else\textsuperscript{\ensuremath{k}}\fi}
\newunicodechar{ᵖ}{\ifmmode^{p}\else\textsuperscript{\ensuremath{p}}\fi}
\newunicodechar{ᴺ}{\ifmmode^{N}\else\textsuperscript{\ensuremath{N}}\fi}
\newunicodechar{ᶜ}{\ifmmode^{c}\else\textsuperscript{\ensuremath{c}}\fi}
\newunicodechar{ʰ}{\ifmmode^{h}\else\textsuperscript{\ensuremath{h}}\fi}
\newunicodechar{ᵠ}{\ifmmode^{q}\else\textsuperscript{\ensuremath{q}}\fi}
\newunicodechar{ₙ}{\ifmmode_{n}\else\textsubscript{\ensuremath{n}}\fi}
\newunicodechar{ₐ}{\ifmmode_{a}\else\textsubscript{\ensuremath{a}}\fi}
\newunicodechar{ᵢ}{\ifmmode_{i}\else\textsubscript{\ensuremath{i}}\fi}
\newunicodechar{ᵣ}{\ifmmode_{r}\else\textsubscript{\ensuremath{r}}\fi}
\newunicodechar{ₖ}{\ifmmode_{k}\else\textsubscript{\ensuremath{k}}\fi}
\newunicodechar{ᵦ}{\ifmmode_{\beta}\else\textsubscript{\ensuremath{\beta}}\fi}
\newunicodechar{ₓ}{\ifmmode_{x}\else\textsubscript{\ensuremath{x}}\fi}
\newfontfamily\popdisplay{ArchivoBlack-Regular.ttf}[Path=../../../fonts/]
\newfontfamily\poplogo{SpaceGrotesk-Bold.ttf}[Path=../../../fonts/]
\newfontfamily\popsemi{PlusJakartaSans-SemiBold.ttf}[Path=../../../fonts/]
\newfontfamily\popxbold{PlusJakartaSans-ExtraBold.ttf}[Path=../../../fonts/]
\newfontfamily\popxboldls{PlusJakartaSans-ExtraBold.ttf}[Path=../../../fonts/,LetterSpace=3.0]

\setlist[enumerate]{leftmargin=1.7em,itemsep=5pt,topsep=4pt}
\setlist[itemize]{leftmargin=1.7em,itemsep=4pt,topsep=4pt,label={\color{poporange}\textbullet}}
\setlength{\parindent}{0pt}
\setlength{\parskip}{5pt}
\renewcommand{\arraystretch}{1.25}

% pgfplots : style Pop par défaut
\pgfplotsset{
  every axis/.append style={axis line style={popink,line width=1pt},tick style={popink},
    grid style={popline,line width=0.5pt},label style={font=\small},tick label style={font=\footnotesize}},
  popcurve/.style={poporange,line width=1.8pt,no markers,smooth},
}
% Même style disponible dans un \draw TikZ simple (hors pgfplots)
\tikzset{popcurve/.style={poporange,line width=1.8pt}}
\tikzset{popgrid/.style={popline,very thin}}
\colorlet{popcurve}{poporange} % le nom du style est parfois utilisé comme couleur

% ============================================================
% Pill / squiggle
% ============================================================
\newcommand{\poppill}[3]{%
  \tikz[baseline=-0.55ex]\node[draw=popink,line width=1.2pt,rounded corners=2.6pt,%
    fill=#2,inner xsep=6.5pt,inner ysep=3pt]{%
    \popxboldls\fontsize{7}{7}\selectfont\color{#3}\MakeUppercase{#1}};%
}
\newcommand{\popsquiggle}[1]{%
  \tikz[baseline=-0.4ex]{
    \draw[#1,line width=2.6pt,line cap=round]
      plot [smooth] coordinates {(0,0.09) (0.34,0.01) (0.68,0.21) (1.02,0.01) (1.36,0.10)};
  }%
}
% Mot-clé surligné (marqueur orange sous le texte)
\newcommand{\cle}[1]{{\popxbold\color{popdark}#1}}

% ============================================================
% En-tête / pied
% ============================================================
\pagestyle{fancy}
\fancyhf{}
\renewcommand{\headrulewidth}{0pt}
\renewcommand{\footrulewidth}{0pt}
\lhead{\raisebox{-0.15ex}{\poplogo\fontsize{13}{13}\selectfont\color{popink}OnBuch\color{poporange}+}}
\rhead{\poppill{\DOCMATIERE\ \textbullet\ \DOCNIVEAU\ \textbullet\ Leçon \DOCLECON}{white}{popink}}
\lfoot{\popxbold\fontsize{7.6}{7.6}\selectfont\color{popmuted}\MakeUppercase{Cours signé OnBuch+}\ \textbullet\ {\popsemi\DOCTITRE}}
\rfoot{\tikz[baseline=-0.6ex]\node[rounded corners=8.5pt,fill=popink,minimum height=17pt,minimum width=17pt,inner xsep=5pt,inner sep=0]{\popxbold\fontsize{8}{8}\selectfont\color{popcream}\thepage\,/\,\pageref*{LastPage}};}

% ============================================================
% Titres
% ============================================================
\newcommand{\chaptitre}[2]{%
  \begin{tcolorbox}[enhanced,colback=poporange,colframe=popink,boxrule=2.6pt,arc=11pt,
      drop shadow=popink,left=15pt,right=15pt,top=12pt,bottom=11pt,before skip=3pt,after skip=18pt]
    \poppill{#2}{popink}{popcream}\par\vspace{9pt}
    {\raggedright\hyphenpenalty=10000\exhyphenpenalty=10000\popdisplay\fontsize{22}{24}\selectfont\color{popcream}\MakeUppercase{#1}\par}
  \end{tcolorbox}
}
% Section numérotée : pastille orange avec le numéro + titre Archivo
\newcounter{popsec}
\newcommand{\coursec}[1]{%
  \stepcounter{popsec}\setcounter{popsub}{0}%
  \par\vspace{16pt}\noindent
  \begin{minipage}{\linewidth}
  \tikz[baseline=-0.7ex]\node[draw=popink,line width=1.6pt,fill=poporange,rounded corners=4pt,minimum size=22pt,inner sep=0]{\popdisplay\fontsize{12}{12}\selectfont\color{popcream}\thepopsec};%
  \hspace{8pt}{\popdisplay\fontsize{15}{17}\selectfont\color{popink}\MakeUppercase{#1}}\par
  \vspace{4pt}\noindent{\color{poporange}\rule{36pt}{3.6pt}}
  \end{minipage}\par\nopagebreak\vspace{8pt}
}
\newcounter{popsub}
\newcommand{\courssub}[1]{%
  \stepcounter{popsub}%
  \par\vspace{9pt}\noindent{\popxbold\fontsize{11.5}{13}\selectfont\color{popink}{\color{poporange}\thepopsec.\thepopsub}\hspace{6pt}#1}\par\nopagebreak\vspace{4pt}
}

% ============================================================
% Boîtes pédagogiques — même recette : fond blanc, bordure épaisse colorée,
% ombre dure encre, badge plein. Titre optionnel [#1].
% ============================================================
\tcbset{popbase/.style={breakable,enhanced,colback=white,boxrule=2.3pt,arc=9pt,
  shadow={3pt}{-3pt}{0pt}{popink},left=14pt,right=14pt,top=11pt,bottom=11pt,before skip=11pt,after skip=15pt,
  fontupper=\popsemi\fontsize{10.6}{14.5}\selectfont\color{popink},
  fontlower=\popsemi\fontsize{10.6}{14.5}\selectfont\color{popink}}}
% Coche dessinée (le glyphe ✓ n'existe pas dans les polices du document)
\newcommand{\popcheck}[1]{\tikz[baseline=-0.5ex]\draw[#1,line width=1.6pt,line cap=round,line join=round](-0.13,0.0)--(-0.04,-0.09)--(0.13,0.1);}
\providecommand{\coche}{\popcheck{popgreen}}
\providecommand{\resultat}[1]{\par\smallskip\noindent\fcolorbox{popgreen}{popgreenL}{\parbox{\dimexpr\linewidth-2\fboxsep-2\fboxrule\relax}{\textbf{Résultat.} #1}}\par\smallskip}
\newcommand{\popboxhead}[3]{\poppill{#1}{#2}{white}\ifx\relax#3\relax\else\hspace{6pt}{\popxbold\fontsize{10.6}{12}\selectfont\color{popink}#3}\fi\par\vspace{7pt}}

\newtcolorbox{aretenir}[1][]{popbase,colframe=popgreen,before upper={\popboxhead{À retenir}{popgreen}{#1}}}
% Texte en allemand (document de lecture, dialogue, extrait) : cadre
% d'étude. Un titre optionnel (source, auteur) s'affiche dans l'en-tête.
\newtcolorbox{textebox}[1][]{popbase,colframe=popblue,before upper={\popboxhead{Text}{popblue}{#1}\begin{otherlanguage}{ngerman}},after upper={\end{otherlanguage}}}
% Marques ✓ / ✗ (forme correcte / fautive) souvent utilisées par les modèles
\providecommand{\cross}{\textcolor{poppink}{\ensuremath{\times}}}
\providecommand{\xmark}{\cross}\providecommand{\crossmark}{\cross}\providecommand{\cmark}{\ensuremath{\checkmark}}
% Ligne de réponse / justification à compléter (inventée par les modèles)
\providecommand{\justification}[1]{\par\noindent\textit{Justification~:} \dotfill\par}
% \textipa (package tipa non chargé) : la police ne couvre pas l'API -> simple passage
\providecommand{\textipa}[1]{#1}
% Soulignement (ulem non chargé) : \uline, \ul
\providecommand{\uline}[1]{\underline{#1}}\providecommand{\ul}[1]{\underline{#1}}
% \glossaire{\item ...} inventée par les modèles : liste de glossaire
\providecommand{\glossaire}[1]{\begin{itemize}#1\end{itemize}}
% \alert (beamer) inventée par les modèles : texte mis en évidence
\providecommand{\alert}[1]{\textbf{\textcolor{poppink}{#1}}}
% variantes ASCII d'ordinaux écrites en macro
\providecommand{\iere}{\textsuperscript{re}}\providecommand{\ieme}{\textsuperscript{e}}\providecommand{\ere}{\textsuperscript{re}}\providecommand{\eme}{\textsuperscript{e}}\providecommand{\er}{\textsuperscript{er}}\providecommand{\ie}{\textsuperscript{e}}
% Ordinaux français écrits en macro par les modèles : 1\ière, 3\ième
\newcommand{\ière}{\textsuperscript{re}}\newcommand{\ième}{\textsuperscript{e}}
% \exemple (inventée par les modèles) : étiquette « Exemple : »
\providecommand{\exemple}{\textbf{Exemple~:}\ }
% \square utilisable aussi hors mode math (cases à cocher)
\AtBeginDocument{\let\squareMath\square\renewcommand{\square}{\ensuremath{\squareMath}}}
% Mot ou phrase en allemand à l'intérieur d'un paragraphe français
\newcommand{\all}[1]{\ifmmode\text{\textit{#1}}\else\begingroup\selectlanguage{ngerman}\itshape #1\endgroup\fi}
\let\esp\all % alias toléré
\newtcolorbox{exemplebox}[1][]{popbase,colframe=poporange,before upper={\popboxhead{Exemple}{poporange}{#1}}}
\newtcolorbox{definition}[1][]{popbase,colframe=popblue,before upper={\popboxhead{Définition}{popblue}{#1}}}
\newtcolorbox{propriete}[1][]{popbase,colframe=poppurple,before upper={\popboxhead{Propriété}{poppurple}{#1}}}
\newtcolorbox{methode}[1][]{popbase,colframe=popgold,before upper={\popboxhead{Méthode}{popgold}{#1}}}
\newtcolorbox{attention}[1][]{popbase,colframe=poppink,before upper={\popboxhead{Attention}{poppink}{#1}}}
\newtcolorbox{experience}[1][]{popbase,colframe=popblue,colback=popblueL,before upper={\popboxhead{Expérience}{popblue}{#1}}}
% « Le savais-tu ? » : carte violette pleine (contexte, culture, Cameroun)
\newtcolorbox{savaistu}[1][]{popbase,colback=poppurple,colframe=popink,
  fontupper=\popsemi\fontsize{10.4}{14.2}\selectfont\color{white},
  before upper={\poppill{Le savais-tu\,?}{white}{poppurple}\ifx\relax#1\relax\else\hspace{6pt}{\popxbold\color{white}#1}\fi\par\vspace{7pt}}}
% Exercice résolu : énoncé en haut, \tcblower puis la solution sur fond crème
\newcounter{popexo}\newcounter{popexr}
\newtcolorbox{exoresolu}[1][]{popbase,colframe=poporange,colbacklower=popcream,
  segmentation style={popink,line width=1.2pt,dashed},
  before upper={\stepcounter{popexr}\popboxhead{Exercice résolu \thepopexr}{poporange}{#1}},
  before lower={\poppill{Solution}{popink}{popcream}\par\vspace{6pt}}}
% Exercice (non résolu en place) — la correction est dans la partie Corrigés
\newtcolorbox{exercice}[1][]{popbase,colframe=popink,
  before upper={\stepcounter{popexo}\popboxhead{Exercice \thepopexo}{popink}{#1}}}
% Corrigé
\newtcolorbox{corrige}[1][]{popbase,colframe=popgreen,colback=popgreenL,
  before upper={\popboxhead{Corrigé}{popgreen}{#1}}}
% Objectifs / prérequis / bilan
\newtcolorbox{objectifs}{popbase,colframe=popink,colback=poporangeL,
  before upper={\popboxhead{Objectifs de la leçon}{popink}{}}}
\newtcolorbox{prerequis}{popbase,colframe=popink,colback=popblueL,
  before upper={\popboxhead{Prérequis}{popblue}{}}}
\newtcolorbox{bilan}{popbase,colframe=popink,colback=white,boxrule=2.8pt,
  before upper={\popboxhead{Fiche bilan}{poporange}{L'essentiel à connaître pour l'examen}}}
% Figure encadrée avec légende
\newtcolorbox{popfigure}[1][]{enhanced,breakable=false,colback=white,colframe=popline,boxrule=1.4pt,arc=8pt,
  left=10pt,right=10pt,top=10pt,bottom=8pt,before skip=10pt,after skip=12pt,halign=center,
  lower separated=false,halign lower=center,
  fontlower=\popsemi\fontsize{9}{11.5}\selectfont\color{popmuted},
  #1}
\newcommand{\legende}[1]{\par\vspace{6pt}{\popsemi\fontsize{9}{11.5}\selectfont\color{popmuted}#1}}

% ============================================================
% Signature OnBuch+
% ============================================================
\newcommand{\poplogoinline}[1]{{\poplogo\fontsize{#1}{#1}\selectfont\color{popink}OnBuch\color{poporange}+}}
\newcommand{\signatureonbuch}{%
  \begin{tcolorbox}[enhanced,colback=popink,colframe=popink,boxrule=2.2pt,arc=10pt,
      drop shadow=poporange,left=14pt,right=14pt,top=10pt,bottom=10pt,before skip=0pt,after skip=0pt]
    \tikz[baseline=-0.7ex]\node[circle,fill=popgreen,draw=popcream,line width=1.3pt,minimum size=22pt,inner sep=0]{\popcheck{white}};%
    \hspace{9pt}\begin{minipage}[c]{0.62\linewidth}
      {\popxbold\fontsize{12}{13}\selectfont\color{popcream}Signé\ \textbullet\ {\poplogo OnBuch\color{poporange}+}}\par\vspace{2pt}
      {\raggedright\popsemi\fontsize{8}{10.5}\selectfont\color{popline}\MakeUppercase{Équipe pédagogique OnBuch+}\\\MakeUppercase{Conforme au programme officiel MINESEC}\par}
    \end{minipage}\hfill
    {\poplogo\fontsize{22}{22}\selectfont\color{popcream}OnBuch\color{poporange}+}
  \end{tcolorbox}%
}

% ============================================================
% Page de garde
% #1 titre · #2 matière · #3 niveau · #4 module · #5 n° leçon · #6 durée ·
% #7 description / objectifs (texte ou itemize)
% ============================================================
\newcommand{\pagegarde}[7]{%
  \thispagestyle{empty}
  \newgeometry{a4paper,margin=1.7cm,top=1.6cm,bottom=1.4cm}%
  \begin{tikzpicture}[remember picture,overlay]
    \fill[poporange!18] ([xshift=-3.2cm,yshift=-3.6cm]current page.north east) circle (4.6cm);
    \fill[poppurple!14] ([xshift=2.4cm,yshift=7.6cm]current page.south west) circle (3.8cm);
  \end{tikzpicture}%
  {\poplogo\fontsize{30}{30}\selectfont\color{popink}OnBuch\color{poporange}+}\hfill
  \raisebox{0.6ex}{\poppill{Cours signé \textbullet\ Premium}{popink}{popcream}}\par
  \vspace{1pt}\popsquiggle{poppurple}\par
  \vspace{8pt}
  \begin{tcolorbox}[enhanced,colback=poporange,colframe=popink,boxrule=2.8pt,arc=13pt,
      drop shadow=popink,left=18pt,right=18pt,top=12pt,bottom=13pt,after skip=10pt]
    \poppill{#2 \textbullet\ #3}{popink}{popcream}\hspace{5pt}\poppill{#4}{white}{popink}\par\vspace{8pt}
    {\popdisplay\fontsize{14}{14}\selectfont\color{popink}LEÇON #5}\par\vspace{4pt}
    {\raggedright\hyphenpenalty=10000\exhyphenpenalty=10000\popdisplay\fontsize{25}{27}\selectfont\color{popcream}\MakeUppercase{#1}\par}
    \vspace{8pt}
    {\popxbold\fontsize{9.5}{9.5}\selectfont\color{popink}\MakeUppercase{Durée conseillée : #6}}
  \end{tcolorbox}
  \begin{tcolorbox}[enhanced,colback=white,colframe=popink,boxrule=2.2pt,arc=10pt,
      drop shadow=popink,left=16pt,right=16pt,top=10pt,bottom=10pt,after skip=8pt]
    \poppill{Dans cette leçon}{poppurple}{white}\par\vspace{5pt}
    \popsemi\fontsize{9.3}{12.6}\selectfont\color{popink}#7
  \end{tcolorbox}
  \noindent\begin{minipage}[t]{0.487\linewidth}
    \begin{tcolorbox}[enhanced,colback=popgreen,colframe=popink,boxrule=2.2pt,arc=10pt,
        drop shadow=popink,left=11pt,right=11pt,top=10pt,bottom=10pt,height=3.1cm,valign=center]
      \tcbox[colback=white,colframe=popink,boxrule=1.3pt,arc=5pt,boxsep=0pt,nobeforeafter,
        left=3pt,right=3pt,top=3pt,bottom=3pt,valign=center]{\qrcode[height=1.9cm]{https://play.google.com/store/apps/details?id=com.luvvix.onbuch&hl=fr}}%
      \hspace{8pt}\begin{minipage}[c]{0.5\linewidth}
        {\popdisplay\fontsize{11}{12.5}\selectfont\color{white}\MakeUppercase{Télécharge OnBuch}}\par\vspace{4pt}
        {\raggedright\popsemi\fontsize{8.4}{11}\selectfont\color{white}Gratuit \textbullet\ Google Play\\Tuteur IA, annales, concours, résultats\par}
      \end{minipage}
    \end{tcolorbox}
  \end{minipage}\hfill
  \begin{minipage}[t]{0.487\linewidth}
    \begin{tcolorbox}[enhanced,colback=poppurple,colframe=popink,boxrule=2.2pt,arc=10pt,
        drop shadow=popink,left=11pt,right=11pt,top=10pt,bottom=10pt,height=3.1cm,valign=center]
      \tikz[baseline=-0.7ex]\node[circle,fill=white,draw=popink,line width=1.3pt,minimum size=1.6cm,inner sep=0]{\popdisplay\fontsize{24}{24}\selectfont\color{poppurple}+};%
      \hspace{8pt}\begin{minipage}[c]{0.6\linewidth}
        {\popdisplay\fontsize{11}{12.5}\selectfont\color{white}\MakeUppercase{Passe à OnBuch+}}\par\vspace{4pt}
        {\raggedright\popsemi\fontsize{8.4}{11}\selectfont\color{white}Tous les cours signés, Tuteur IA illimité, fiches PDF — onglet OnBuch+ de l'app\par}
      \end{minipage}
    \end{tcolorbox}
  \end{minipage}
  \vspace{8pt}
  \popcontenu
  \vfill
  \signatureonbuch
  \restoregeometry
  \newpage
}

% Rangée « Ce cours contient » (page de garde)
\newcommand{\poptile}[3]{% #1 couleur #2 gros texte #3 légende
  \begin{tcolorbox}[enhanced,colback=#1,colframe=popink,boxrule=2pt,arc=9pt,shadow={2.5pt}{-2.5pt}{0pt}{popink},
      left=6pt,right=6pt,top=9pt,bottom=9pt,halign=center,valign=center,height=1.75cm,width=\linewidth,nobeforeafter]
    {\popdisplay\fontsize{12.5}{13}\selectfont\color{white}\MakeUppercase{#2}}\par\vspace{3pt}
    {\centering\popsemi\fontsize{7.8}{9.5}\selectfont\color{white}#3\par}
  \end{tcolorbox}}
\newcommand{\popcontenu}{%
  \par\noindent{\popxboldls\fontsize{7.5}{7.5}\selectfont\color{popmuted}CE COURS SIGNÉ CONTIENT}\par\vspace{7pt}
  \noindent\begin{minipage}[t]{0.235\linewidth}\poptile{popblue}{Cours}{complet, illustré, pas à pas}\end{minipage}\hfill
  \begin{minipage}[t]{0.235\linewidth}\poptile{poporange}{Méthodes}{et exercices résolus}\end{minipage}\hfill
  \begin{minipage}[t]{0.235\linewidth}\poptile{poppink}{Exercices}{type examen + corrigés}\end{minipage}\hfill
  \begin{minipage}[t]{0.235\linewidth}\poptile{popgreen}{Fiche bilan}{l'essentiel à retenir}\end{minipage}\par}

% Sommaire
\newtcolorbox{sommaire}{enhanced,colback=white,colframe=popink,boxrule=2.2pt,arc=10pt,
  drop shadow=popink,left=18pt,right=18pt,top=13pt,bottom=13pt,before skip=6pt,after skip=20pt,
  before upper={\poppill{Sommaire}{poppurple}{white}\par\vspace{9pt}\popsemi\fontsize{10.6}{14.5}\selectfont\color{popink}}}

% Signature de fin de cours — poussée tout en bas de la dernière page
\newcommand{\findecours}{%
  \par\vfill\nopagebreak
  \begin{center}\popsquiggle{poporange}\end{center}\vspace{4pt}
  \signatureonbuch
}
\AtBeginDocument{\def\llbracket{\mathopen{[\mkern-3.5mu[}}\def\rrbracket{\mathclose{]\mkern-3.5mu]}}} % intervalles d'entiers (glyphe ⟦ absent de la police)
% Glyphes absents de Fira Math : redéfinis à partir de glyphes disponibles
\AtBeginDocument{%
  \def\setminus{\mathbin{\backslash}}\let\smallsetminus\setminus
  \def\vdots{\mathord{\vbox{\baselineskip3.2pt\lineskiplimit0pt\kern4pt\hbox{.}\hbox{.}\hbox{.}}}}%
  \def\ddots{\mathinner{\mkern1mu\raise6.5pt\hbox{.}\mkern2mu\raise3.6pt\hbox{.}\mkern2mu\raise0.7pt\hbox{.}\mkern1mu}}%
  \def\triangle{\mathord{\scalebox{0.9}{$\symup{\Delta}$}}}%
  \def\nabla{\mathord{\raisebox{\depth}{\scalebox{1}[-1]{$\symup{\Delta}$}}}}%
  \def\checkmark{\popcheck{popdark}}%
  \def\star{\mathbin{\ast}}\def\bigstar{\mathord{\ast}}%
  \def\diamond{\mathbin{\scalebox{0.6}{$\Diamond$}}}%
  \def\bigcup{\mathop{\mathchoice{\raisebox{-0.3em}{\scalebox{1.7}{$\cup$}}}{\scalebox{1.25}{$\cup$}}{\cup}{\cup}}\displaylimits}%
  \def\bigcap{\mathop{\mathchoice{\raisebox{-0.3em}{\scalebox{1.7}{$\cap$}}}{\scalebox{1.25}{$\cap$}}{\cap}{\cap}}\displaylimits}%
  \def\lesseqqgtr{\mathrel{\lessgtr}}\def\bigoplus{\mathop{\oplus}\displaylimits}%
}
\providecommand{\ssi}{si et seulement si} % abréviation fréquente
\AtBeginDocument{\providecommand{\wideparen}{\overparen}} % arc de cercle

%% FIGFIX-BEGIN v3 — correctifs globaux des figures (docs/onbuch-generation/tools/figfix)
\usepackage{adjustbox}\usepackage{etoolbox}
% toute figure TikZ/pgfplots plus large que la ligne est réduite à la largeur disponible
\BeforeBeginEnvironment{tikzpicture}{\begin{adjustbox}{max width=\linewidth}}
\AfterEndEnvironment{tikzpicture}{\end{adjustbox}}
% légendes pgfplots lisibles par-dessus les courbes
\pgfplotsset{every axis/.append style={legend style={fill=white,fill opacity=0.92,text opacity=1,draw=black!15,inner sep=3pt}}}
% étiquettes d'arêtes (syntaxe edge["..."]) avec fond blanc
\tikzset{every edge quotes/.append style={fill=white,inner sep=1.2pt}}
% bulles de cartes mentales : texte à la ligne dans la bulle, bulle un peu plus grande
\tikzset{every concept/.append style={text width=2.0cm,align=center,minimum size=2.3cm,inner sep=2pt}}
%% FIGFIX-END
\begin{document}

\pagegarde{La cybercriminalité 2 : le harcèlement numérique}{Allemand}{Tle A}{Chapitre 4 : Média et communication}{AL18}{2 h}{Leçon de lecture-commentaire sur le harcèlement numérique (Cybermobbing) : compréhension d'un texte allemand, vocabulaire thématique, passif (présent et prétérit) et connecteurs de concession/opposition. Elle prépare à la compréhension écrite et à l'expression guidée du Bac A.
\vspace{4pt}
\begin{multicols}{2}\small
\begin{itemize}
\item À la fin de cette leçon, je sais comprendre globalement puis en détail un texte allemand sur le harcèlement numérique.
\item À la fin de cette leçon, je sais employer le vocabulaire du thème (das Cybermobbing, das Opfer, der Täter, die Anzeige, schützen, das Passwort) avec articles et pluriels.
\item À la fin de cette leçon, je sais former et utiliser le passif au présent et au prétérit (Vorgangspassiv).
\item À la fin de cette leçon, je sais transformer une phrase active en phrase passive et inversement.
\item À la fin de cette leçon, je sais exprimer la concession et l'opposition avec obwohl, trotzdem, aber, während, trotz + génitif.
\item À la fin de cette leçon, je sais répondre par écrit à des questions de compréhension en phrases complètes.
\end{itemize}\end{multicols}}

\begin{sommaire}
\begin{multicols}{2}
\begin{enumerate}
\item Texte support : Cybermobbing – ein Fall aus Berlin
\item Compréhension du texte
\item Lexique thématique : das Cybermobbing
\item Grammaire 1 : le passif (Vorgangspassiv)
\item Grammaire 2 : concession et opposition
\item Méthode du commentaire et commentaire modèle
\item Production guidée et entraînement
\item Activité d'intégration
\item Méthodes et exercices résolus
\item Exercices
\item Corrigés des exercices
\item Fiche bilan
\end{enumerate}
\end{multicols}
\end{sommaire}

\begin{objectifs}
\begin{itemize}
\item À la fin de cette leçon, je sais comprendre globalement puis en détail un texte allemand sur le harcèlement numérique.
\item À la fin de cette leçon, je sais employer le vocabulaire du thème (das Cybermobbing, das Opfer, der Täter, die Anzeige, schützen, das Passwort) avec articles et pluriels.
\item À la fin de cette leçon, je sais former et utiliser le passif au présent et au prétérit (Vorgangspassiv).
\item À la fin de cette leçon, je sais transformer une phrase active en phrase passive et inversement.
\item À la fin de cette leçon, je sais exprimer la concession et l'opposition avec obwohl, trotzdem, aber, während, trotz + génitif.
\item À la fin de cette leçon, je sais répondre par écrit à des questions de compréhension en phrases complètes.
\item À la fin de cette leçon, je sais commenter un texte (situation, structure, thème, avis argumenté).
\item À la fin de cette leçon, je sais produire un court paragraphe guidé (conseils contre le Cybermobbing) en réutilisant la grammaire de la leçon.
\end{itemize}
\end{objectifs}

\begin{prerequis}
\begin{itemize}
\item Les verbes à accusatif et à datif (Première) : jemanden beleidigen, jemandem helfen.
\item Les auxiliaires haben/sein et le prétérit des verbes faibles et forts.
\item La place du verbe : proposition principale (V2) et subordonnée (verbe en fin).
\item Les connecteurs de base : und, aber, denn, weil, dass.
\item Le vocabulaire des médias vu au chapitre 4 (das Handy, das Internet, die sozialen Netzwerke).
\end{itemize}
\end{prerequis}

\newpage

\coursec{Texte support : Cybermobbing – ein Fall aus Berlin}

\courssub{Situation d'accroche et problématique}

Imaginez la scène : un élève de Terminale à Douala reçoit, soir après soir, des messages humiliants sur WhatsApp. Des photos retouchées circulent dans le groupe de la classe. Il n'ose plus aller au lycée. Ce scénario, hélas banal, se joue aussi à Berlin, Munich ou Vienne. En Allemagne, ce phénomène a un nom précis : \all{das Cybermobbing}. Il ne s'agit pas d'une simple « blague » : l'insulte publique (\all{die Beleidigung}) est un délit pénal poursuivi d'office selon le \all{\S 185 StGB} (Code pénal), passible d'amende ou de peine de prison.

\textbf{Problématique de la leçon :} Comment comprendre, analyser et discuter en allemand d'un cas concret de harcèlement numérique ? Pour y répondre, nous allons exploiter un texte authentique, maîtriser le vocabulaire thématique, et découvrir deux outils grammaticaux indispensables pour rapporter des faits : \textbf{le passif} (qui met l'accent sur la victime ou l'action) et \textbf{les connecteurs de concession/opposition} (qui nuancent le propos).

\begin{methode}[Stratégie de lecture en trois temps]
\begin{enumerate}
    \item \textbf{Première lecture (globale) :} Lisez le texte une fois sans dictionnaire. Cherchez l'idée générale : \emph{De qui parle-t-on ? Qu'est-il arrivé ? Quelle est l'issue ?} Ne traduisez pas mot à mot.
    \item \textbf{Deuxième lecture (détaillée) :} Soulignez les \textbf{mots-clés} (noms, verbes, connecteurs), les \textbf{dates/chiffres} et les \textbf{formes verbales passives} (auxiliaire \all{werden} + participe passé). Repérez les connecteurs de concession (\all{obwohl}, \all{trotz}).
    \item \textbf{Troisième lecture (exploitation) :} Répondez aux questions de compréhension précises. Reformulez l'histoire en 3 phrases à l'oral.
\end{enumerate}
\end{methode}

\courssub{Texte support : Cybermobbing – ein Fall aus Berlin}

\begin{textebox}[Cybermobbing – ein Fall aus Berlin]
In Berlin \textbf{wurde} ein 16-jähriger Schüler monatelang im Internet \textbf{beleidigt}. Jeden Abend \textbf{wurden} hässliche Fotos von ihm in einer WhatsApp-Gruppe \textbf{verbreitet}. \textbf{Obwohl} er die Nachrichten \textbf{löschte}, \textbf{wurden} sie immer wieder \textbf{gesendet}. Das Opfer hatte Angst und wollte nicht mehr zur Schule gehen. \textbf{Trotz} der vielen Beweise wagte er lange keine Anzeige. Schließlich \textbf{wurde} der Täter von der Polizei \textbf{gefunden}: es war ein Mitschüler. Er \textbf{wurde} von der Schule \textbf{suspendiert} und musste eine Strafe zahlen. Heute \textbf{wird} der Junge von einer Psychologin \textbf{betreut}. Experten raten: Passwörter \textbf{müssen} geheim \textbf{bleiben}, und Opfer \textbf{müssen} \textbf{geschützt} \textbf{werden}. Wer \textbf{wird} \textbf{gemobbt}, \textbf{soll} Beweise speichern und sich an Lehrer oder Eltern wenden.
\end{textebox}

\noindent\textbf{Observation guidée :} Dans le texte ci-dessus, les formes du \textbf{passif (Vorgangspassiv)} sont en \textbf{gras} (ex. \textbf{wurde ... beleidigt}, \textbf{wurden ... verbreitet}, \textbf{wurden ... gesendet}, \textbf{wurde ... gefunden}, \textbf{wurde ... suspendiert}, \textbf{wird ... betreut}, \textbf{müssen ... geschützt werden}, \textbf{wird ... gemobbt}). Les connecteurs de concession/opposition sont également en \textbf{gras} (\textbf{Obwohl}, \textbf{Trotz}).

\courssub{Glossaire en marge (10 mots-clés)}

\begin{center}
\begin{tabularx}{\linewidth}{|>{\bfseries}l|X|X|}
\hline
\rowcolor{popblueL}
\textbf{Allemand (Article + Pluriel)} & \textbf{Français} & \textbf{Contexte / Famille} \\
\hline
\all{die Beleidigung, -en} & l'insulte, l'offense & \all{beleidigen} (v.) insulter ; \all{\S 185 StGB} \\
\hline
\all{verbreiten} (regelm.) & propager, diffuser, répandre & \all{die Verbreitung} (propagation) ; \all{verbreitet} (part. passé) \\
\hline
\all{melden} (regelm.) & signaler, déclarer, annoncer & \all{die Meldung} (signalement) ; \all{sich melden} (se manifester) \\
\hline
\all{die Anzeige, -n} & la plainte (police), la dénonciation & \all{Anzeige erstatten} (déposer plainte) ; \all{anzeigen} (v., non séparable) \\
\hline
\all{der Täter, -} & l'auteur (d'un délit), l'agresseur & fém. \all{die Täterin} ; \all{die Tat} (le fait, l'acte) \\
\hline
\all{das Opfer, -} & la victime & \all{opfern} (sacrifier) ; \all{Opferschutz} (protection victimes) \\
\hline
\all{schützen} (regelm.) & protéger & \all{der Schutz} (protection) ; \all{Schutzmaßnahme} (mesure de protection) \\
\hline
\all{das Passwort, -wörter} & le mot de passe & \all{geheim bleiben} (rester secret) ; \all{Passwortsicherheit} \\
\hline
\all{suspendieren} (regelm.) & suspendre, exclure temporairement & \all{die Suspendierung} ; scolaire : \all{von der Schule suspendieren} \\
\hline
\all{betreuen} (regelm.) & prendre en charge, encadrer, suivre & \all{die Betreuung} (prise en charge) ; \all{der Betreuer} (tuteur, encadrant) \\
\hline
\end{tabularx}
\end{center}

\courssub{Compréhension globale : Repérage des informations clés}

\begin{exoresolu}[Repérage « Qui ? Où ? Quoi ? Conséquences ? »]
\textbf{Énoncé :} Lisez le texte et complétez le tableau de repérage ci-dessous en allemand (citation du texte) et en français.
\tcblower
\textbf{Solution détaillée :}
\begin{center}
\begin{tabularx}{\linewidth}{|>{\bfseries}l|X|X|}
\hline
\rowcolor{popblueL}
\textbf{Question} & \textbf{Extrait du texte (Allemand)} & \textbf{Traduction / Réponse (Français)} \\
\hline
\textbf{Wer? (Qui ?)} & \all{ein 16-jähriger Schüler} / \all{der Täter ... ein Mitschüler} & Un élève de 16 ans / L'auteur : un camarade de classe \\
\hline
\textbf{Wo? (Où ?)} & \all{In Berlin} / \all{im Internet} / \all{in einer WhatsApp-Gruppe} / \all{zur Schule} & À Berlin / Sur Internet / Dans un groupe WhatsApp / À l'école \\
\hline
\textbf{Was? (Quoi ?)} & \all{beleidigt} / \all{hässliche Fotos ... verbreitet} / \all{Nachrichten löschte} / \all{Anzeige} / \all{Polizei gefunden} & Insulté / Photos moches diffusées / Messages effacés / Plainte / Police a trouvé \\
\hline
\textbf{Folgen? (Conséquences)} & \all{Angst, nicht mehr zur Schule gehen} / \all{suspendiert, Strafe zahlen} / \all{von Psychologin betreut} & Peur, refus d'aller à l'école / Suspension, amende / Suivi psychologique \\
\hline
\end{tabularx}
\end{center}
\end{exoresolu}

\courssub{Analyse linguistique : Le passif (Vorgangspassiv) et la Concession}

\begin{definition}[Le Vorgangspassiv (Passif de procès / action)]
Le \textbf{Vorgangspassiv} décrit une \textbf{action}, un \textbf{processus} ou un \textbf{événement}. Il met l'accent sur \textbf{ce qui arrive} au sujet (qui subit l'action), et non sur l'auteur de l'action.
\par\medskip
\textbf{Forme :} \all{werden} (conjugué) + \textbf{Participe II} (ge-mach-t, ge-seh-en, be-lei-dig-t).
\par\medskip
\textbf{Auteur (Agent) :} Souvent omis. Si nécessaire : préposition \all{von} (+ datif) pour personne/organisme, \all{durch} (+ accusatif) pour moyen/instrument.
\end{definition}

\begin{center}
\begin{tabularx}{\linewidth}{|X|X|X|}
\hline
\rowcolor{popblueL}
\textbf{Temps} & \textbf{Allemand (Exemple du texte / construit)} & \textbf{Français} \\
\hline
\bfseries Präsens & \all{Heute \textbf{wird} der Junge \textbf{betreut}.} & Aujourd'hui, le garçon \textbf{est suivi} / \textbf{est pris en charge}. \\
\hline
\bfseries Präteritum & \all{\textbf{Wurde} ein Schüler \textbf{beleidigt}.} / \all{\textbf{Wurden} Fotos \textbf{verbreitet}.} & Un élève \textbf{a été insulté}. / Des photos \textbf{ont été diffusées}. \\
\hline
\bfseries Perfekt & \all{Der Täter \textbf{ist} ... \textbf{gefunden} \textbf{worden}.} & L'auteur \textbf{a été retrouvé}. (Remarque : \all{worden} sans \emph{ge-}) \\
\hline
\bfseries Futur I & \all{Opfer \textbf{werden} \textbf{geschützt} \textbf{werden}.} & Les victimes \textbf{seront protégées}. \\
\hline
\bfseries Modal (Präsens) & \all{Opfer \textbf{müssen} \textbf{geschützt} \textbf{werden}.} & Les victimes \textbf{doivent être protégées}. (Passif modal : modal + Part. II + \all{werden}) \\
\hline
\end{tabularx}
\end{center}

\begin{attention}[Piège majeur pour francophones : \all{werden} vs \all{sein} / \all{worden} vs \all{geworden}]
\begin{itemize}
    \item \all{werden} = auxiliaire du \textbf{passif de procès} (action en cours / accomplie). \all{Er \textbf{wird} geschlagen.} (Il est en train d'être battu / Il se fait battre).
    \item \all{sein} = auxiliaire de l'\textbf{état passif (Zustandspassiv)} ou du perfect actif. \all{Er \textbf{ist} geschlagen.} (Il est battu = il est dans l'état « battu »).
    \item \textbf{Perfekt passif :} \all{ist ... \textbf{worden}} (sans \emph{ge-} !). \textbf{Jamais} \all{ist ... geworden} (qui signifie « est devenu »).
    \item \textbf{Français « être + participe »} $\neq$ toujours passif allemand. \all{Er ist interessiert} (Il est intéressé = état) vs \all{Er wird interessiert} (Il est intéressé par qqn = action rare).
    \item \all{geheim bleiben} (rester secret) = voix active (verbe intransitif), \textbf{pas} un passif. Le passif modal correct est \all{geheim gehalten werden} (être gardés secrets).
\end{itemize}
\end{attention}

\begin{definition}[Concession et Opposition : \all{obwohl} vs \all{trotz}]
\begin{itemize}
    \item \textbf{\all{obwohl}} (conjonction de subordination) + \textbf{phrase complète} (verbe à la fin). Exprime une concession : « bien que, quoique ».
    \item \textbf{\all{trotz}} (préposition) + \textbf{nom/groupe nominal au GÉNITIF} (formel, écrit) ou DATIF (courant, oral). Exprime une opposition : « malgré, en dépit de ».
\end{itemize}
\end{definition}

\begin{exemplebox}[Exemples tirés du texte et variations]
\begin{itemize}
    \item \all{\textbf{Obwohl} er die Nachrichten \textbf{löschte}, wurden sie gesendet.} (Bien qu'il \textbf{effaçât} les messages, ils furent envoyés.) $\rightarrow$ Verbe \all{löschte} à la fin.
    \item \all{\textbf{Trotz} \textbf{der} vielen \textbf{Beweise} wagte er keine Anzeige.} (Malgré \textbf{les} nombreuses \textbf{preuves}, il n'osa pas porter plainte.) $\rightarrow$ \all{der Beweise} = Génitif pluriel (article \all{der} audible).
    \item \textit{Variante courante (Datif) :} \all{Trotz \textbf{dem} Regen gingen wir spazieren.} (Malgré la pluie, nous sommes sortis.)
\end{itemize}
\end{exemplebox}

\courssub{Carte mentale lexicale : L'univers du Cybermobbing}

\begin{popfigure}
\begin{tikzpicture}[
    mindmap,
    grow cyclic,
    every node/.style={concept, execute at begin node=\hskip0pt},
    concept color=popblue,
    level 1/.style={level distance=4.5cm, sibling angle=90, concept color=popblue!70},
    level 2/.style={level distance=3cm, sibling angle=45, concept color=poporange!70},
    text=white,
    font=\small\bfseries
]
\node {Cybermobbing}
    child[concept color=poppink] { node {Täter}
        child { node {Mitschüler} }
        child { node {Fremde} }
        child { node {Anonym} }
    }
    child[concept color=popgreen] { node {Opfer}
        child { node {Schüler/in} }
        child { node {Angst} }
        child { node {Schule verweigern} }
    }
    child[concept color=poppurple] { node {Folgen}
        child { node {Psychische Belastung} }
        child { node {Strafanzeige (\S 185)} }
        child { node {Suspendierung} }
        child { node {Strafe / Geldbuße} }
    }
    child[concept color=popgold] { node {Hilfe \& Schutz}
        child { node {Beweise speichern} }
        child { node {Eltern / Lehrer} }
        child { node {Polizei / Anzeige} }
        child { node {Psychologin / Beratung} }
        child { node {Passwörter schützen} }
    };
\end{tikzpicture}
\legende{Champ lexical structuré : Acteurs, Victimes, Conséquences juridiques/psychologiques, Pistes de solution.}
\end{popfigure}

\courssub{Focus culturel : Le cadre légal en Allemagne}

\begin{savaistu}[Cybermobbing et droit allemand (\S 185 StGB)]
En Allemagne, l'insulte (\all{die Beleidigung}) est un \textbf{délit poursuivi d'office} (\all{Offizialdelikt}) selon le \all{\S 185 StGB} (Strafgesetzbuch). La peine encourt une amende (\all{Geldstrafe}) ou une peine de prison jusqu'à un an (\all{Freiheitsstrafe bis zu einem Jahr}), deux ans si l'insulte est publique ou par écrit (\all{\S 186, 187 StGB : üble Nachrede, Verleumdung}).
\par\medskip
Depuis 2021, la loi \all{NetzDG} (Netzwerkdurchsetzungsgesetz) oblige les grands réseaux sociaux (Facebook, X, TikTok, YouTube) à supprimer les contenus manifestement illégaux (haine, harcèlement) sous 24h sous peine de lourdes amendes (jusqu'à 50 millions €).
\par\medskip
\textbf{Contexte scolaire :} La plupart des \all{Bundesländer} imposent des \all{Anti-Mobbing-Konzepte} dans les écoles. Des \all{Anti-Mobbing-Tage} (journées de prévention), des \all{Streitschlichter} (médiateurs élèves) et des \all{Vertrauenslehrer} (enseignants de confiance) sont mis en place. Au Cameroun, la loi n°2010/012 sur la cybercriminalité sanctionne aussi le harcèlement électronique (art. 78), mais la prévention scolaire reste à renforcer.
\end{savaistu}

\courssub{Méthode du commentaire écrit (Bac)}

\begin{methode}[Commentaire structuré d'un texte d'actualité]
\begin{enumerate}
    \item \textbf{Introduction (3--4 lignes) :} Présenter le document (type, source supposée, thème, date). Annoncer le problème (\all{das Cybermobbing}) et le plan (1. Les faits 2. La réaction juridique 3. La prévention).
    \item \textbf{Développement (2--3 paragraphes) :} 
    \begin{itemize}
        \item \textit{Paragraphe 1 (Les faits) :} Résumer l'affaire (victime, auteur, moyens, durée). Utiliser le \textbf{passif} (\all{werden} + Part. II) et le \textbf{Präteritum}.
        \item \textit{Paragraphe 2 (Réaction \& Droit) :} Expliquer l'issue (police, sanction, loi \S 185). Utiliser la \textbf{concession} (\all{obwohl}, \all{trotz}) pour montrer la difficulté de la victime à agir.
        \item \textit{Paragraphe 3 (Prévention) :} Citer les conseils d'experts (mots de passe, preuves, adultes de confiance). Utiliser le \textbf{passif modal} (\all{müssen} + Part. II + \all{werden}) et le \textbf{Futur I}.
    \end{itemize}
    \item \textbf{Conclusion (2--3 lignes) :} Bilan : le harcèlement numérique est un délit grave. Ouverture : rôle de l'école (Cameroun / Allemagne), responsabilité des plateformes (\all{NetzDG}).
\end{enumerate}
\end{methode}

\begin{exoresolu}[Commentaire modèle (environ 120 mots)]
\textbf{Énoncé :} Rédigez un commentaire en allemand sur le texte « Cybermobbing – ein Fall aus Berlin ».
\tcblower
\textbf{Proposition de rédaction :}
\par\medskip
\all{Der Text „Cybermobbing – ein Fall aus Berlin“ thematisiert einen Fall von digitaler Gewalt an Schulen. Ein 16-jähriger Schüler wurde monatelang über WhatsApp beleidigt und mit bearbeiteten Fotos diffamiert. \textbf{Obwohl} er die Nachrichten löschte, hörte das Mobbing nicht auf. \textbf{Trotz} vorhandener Beweise zögerte das Opfer, Anzeige zu erstatten. Schließlich wurde der Täter – ein Mitschüler – von der Polizei ermittelt. Er wurde von der Schule suspendiert und bestraft.}
\par\medskip
\all{Der Fall zeigt die rechtliche Dimension: Nach \S 185 StGB ist Beleidigung ein Offizialdelikt. Das NetzDG zwingt Plattformen zum schnellen Löschen illegaler Inhalte. Experten betonen Prävention: Passwörter müssen geheim gehalten werden, Opfer müssen geschützt werden. Wer gemobbt wird, soll Beweise sichern und sich an Vertrauenspersonen wenden.}
\par\medskip
\all{Zusammenfassend ist Cybermobbing keine Bagatelle, sondern ein Strafdelikt mit schweren psychischen Folgen. Schulen in Deutschland und Kamerun müssen Präventionskonzepte stärken.}
\end{exoresolu}

\courssub{Entraînement guidé : Reformulation, Transformation, Thème \& Production}

\begin{exercice}[Entraînement complet (Grammaire, Thème, Version, Expression)]
\textbf{A. Reformulation orale (3 phrases) :} Résumez l'histoire en \textbf{3 phrases allemandes} (Präteritum/Perfekt), en utilisant le vocabulaire du glossaire et le passif.
\par\textbf{B. Transformation Actif $\rightarrow$ Passif (Präteritum) :} Transformez les phrases actives en \textbf{Passif (Präteritum)}. L'agent (\all{von} + datif) est optionnel.
\begin{enumerate}
    \item Ein Mitschüler beleidigte den Jungen monatelang.
    \item Die Polizei fand den Täter schließlich.
    \item Die Schule suspendierte den Täter.
    \item Eine Psychologin betreut das Opfer jetzt. (Présent passif)
    \item Man muss die Passwörter geheim halten. (Passif modal présent)
\end{enumerate}
\par\textbf{C. Thème (Français $\rightarrow$ Allemand) :} Traduisez en allemand.
\begin{enumerate}
    \item Malgré les insultes, le garçon est allé à l'école. (\all{trotz} + génitif)
    \item Bien qu'il ait peur, il a déposé plainte. (\all{obwohl} + Präteritum)
    \item Les mots de passe doivent être gardés secrets. (Passif modal)
    \item La victime est soutenue par un psychologue. (Présent passif)
\end{enumerate}
\par\textbf{D. Version (Allemand $\rightarrow$ Français) :} Traduisez en français.
\all{„Obwohl die Täter oft anonym bleiben, werden sie durch digitale Spuren identifiziert. Trotz strenger Gesetze ist die Dunkelziffer hoch. Deshalb müssen Schulen und Eltern zusammenarbeiten.“}
\par\textbf{E. Production écrite guidée (50--80 mots) :} Vous êtes délégué(e) de classe. Écrivez un courriel formel au directeur (\all{Schulleiter}) pour signaler un cas de cyberharcèlement observé et proposer une action de prévention (\all{Anti-Mobbing-Tag}, \all{Workshop}). Utilisez le passif, \all{obwohl}/\all{trotz}, le vocabulaire du thème.
\end{exercice}

\begin{corrige}[Exercice Entraînement complet]
\textbf{A. Proposition de reformulation orale (Modèle) :}
\begin{enumerate}
    \item \all{In Berlin wurde ein 16-jähriger Schüler monatelang per WhatsApp gemobbt.}
    \item \all{Obwohl er Beweise hatte, erstattete er lange keine Anzeige.}
    \item \all{Der Täter wurde bestraft und das Opfer wird jetzt psychologisch betreut.}
\end{enumerate}

\textbf{B. Transformations (Passif) :}
\begin{enumerate}
    \item \all{Der Junge \textbf{wurde} monatelang \textbf{beleidigt} (von einem Mitschüler).}
    \item \all{Der Täter \textbf{wurde} schließlich \textbf{gefunden} (von der Polizei).}
    \item \all{Der Täter \textbf{wurde} \textbf{suspendiert} (von der Schule).}
    \item \all{Das Opfer \textbf{wird} jetzt \textbf{betreut} (von einer Psychologin).} (Présent passif car « jetzt »)
    \item \all{Die Passwörter \textbf{müssen} geheim \textbf{gehalten} \textbf{werden}.} (Passif modal présent : \all{müssen} + Part. II + \all{werden})
\end{enumerate}
\par\medskip
\textbf{Note grammaticale :} À la phrase 4, le texte original utilise le \textbf{Präsens} (\all{wird betreut}) car c'est une situation actuelle. À la phrase 5, \all{man} disparaît ; le verbe modal \all{müssen} reste à la 2e position, l'infinitif passif \all{gehalten werden} va à la fin.

\textbf{C. Thème (Français $\rightarrow$ Allemand) :}
\begin{enumerate}
    \item \all{Trotz der Beleidigungen ging der Junge zur Schule.} (ou \all{Trotz der Beleidigungen ist der Junge zur Schule gegangen.})
    \item \all{Obwohl er Angst hatte, erstattete er Anzeige.}
    \item \all{Die Passwörter müssen geheim gehalten werden.}
    \item \all{Das Opfer wird von einem Psychologen betreut.}
\end{enumerate}

\textbf{D. Version (Allemand $\rightarrow$ Français) :}
« Bien que les auteurs restent souvent anonymes, ils sont identifiés par des traces numériques. Malgré des lois strictes, le chiffre noir (nombre d'affaires non signalées) est élevé. C'est pourquoi les écoles et les parents doivent collaborer. »

\textbf{E. Production écrite guidée (Modèle) :}
\par\medskip
\all{Betreff: Meldung eines Cybermobbing-Falls und Präventionsvorschlag}
\par\medskip
\all{Sehr geehrter Herr Schulleiter,}
\par\medskip
\all{ich möchte Sie über einen Fall von Cybermobbing in der 12. Klasse informieren. Ein Schüler wurde seit Wochen in einer Chat-Gruppe beleidigt und mit gefälschten Fotos diffamiert. \textbf{Obwohl} das Opfer die Nachrichten löschte, wurden neue Inhalte gepostet. \textbf{Trotz} der Beweise wagte der Betroffene lange nicht, sich zu melden.}
\par\medskip
\all{Das Opfer wird nun von der Schulpsychologin betreut. Der Täter wurde identifiziert und \textbf{von der Schule suspendiert}. Nach \S 185 StGB ist dies eine Straftat.}
\par\medskip
\all{Ich schlage vor, einen \textbf{Anti-Mobbing-Tag} mit Workshops zu organisieren. Passwörter \textbf{müssen geschützt werden}, und Opfer \textbf{müssen ermutigt werden}, sich an \textbf{Vertrauenslehrer} zu wenden.}
\par\medskip
\all{Mit freundlichen Grüßen}
\par\medskip
\all{Max Mustermann, Klassensprecher 12A}
\end{corrige}

\courssub{Synthèse vocabulaire : Familles de mots (Wortfamilien)}

\begin{center}
\begin{tabularx}{\linewidth}{|>{\bfseries}l|X|X|}
\hline
\rowcolor{popblueL}
\textbf{Verbe (Infinitif)} & \textbf{Nom (Article + Pluriel)} & \textbf{Adjectif / Participe / Composés} \\
\hline
\all{beleidigen} & \all{die Beleidigung, -en} & \all{beleidigt} (part.), \all{beleidend} (adj.), \all{die Beleidigungsanzeige} \\
\hline
\all{verbreiten} & \all{die Verbreitung} & \all{verbreitet} (part./adj.), \all{weit verbreitet} \\
\hline
\all{anzeigen / Anzeige erstatten} & \all{die Anzeige, -n} / \all{der Anzeigenerstatter} & \all{anzeigepflichtig}, \all{erstatten} (v.) \\
\hline
\all{schützen} & \all{der Schutz} / \all{die Schutzmaßnahme, -n} & \all{geschützt} (part.), \all{schützenswert}, \all{der Opferschutz} \\
\hline
\all{mobben} (anglicisme, rég.) & \all{das Mobbing} / \all{der Mobber, -} & \all{gemobbt} (part.), \all{mobbingfrei}, \all{Cybermobbing} \\
\hline
\all{suspendieren} & \all{die Suspendierung} & \all{suspendiert} (part.), \all{vorläufig suspendiert} \\
\hline
\all{betreuen} & \all{die Betreuung} / \all{der Betreuer, - / die Betreuerin, -nen} & \all{betreuungsbedürftig}, \all{die Betreuungsstelle} \\
\hline
\end{tabularx}
\end{center}

\begin{attention}[Faux amis et pièges de traduction]
\begin{itemize}
    \item \all{die Anzeige} $\neq$ « l'annonce » (publicité = \all{die Werbung}). Contexte juridique : \all{Anzeige erstatten} = déposer plainte. Contexte commercial : \all{eine Anzeige aufgeben} = passer une petite annonce.
    \item \all{der Täter} $\neq$ « le tuteur » (\all{der Betreuer} / \all{der Vormund}). \all{Täter} = auteur d'un crime/délit (connotation négative).
    \item \all{das Opfer} $\neq$ « l'ouvrier » (\all{der Arbeiter}). \all{Opfer} = victime (aussi : sacrifice religieux \all{Opfer bringen}).
    \item \all{suspendieren} $\neq$ « suspendre » (pendre = \all{aufhängen}). Sens scolaire/admin : exclure temporairement, mettre à pied (\all{vorläufig suspendieren}).
    \item \all{melden} : « signaler » (à la police, plateforme : \all{etwas melden}) OU « s'annoncer/inscrire » (\all{sich melden}, \all{zum Dienst melden}).
    \item \all{anzeigen} (non séparable) = dénoncer, signaler (\all{er zeigte ihn an} $\rightarrow$ Präteritum \all{anzeigte}). \textbf{Ne pas confondre avec \all{an|zeigen} (séparable) = afficher, indiquer.}
\end{itemize}
\end{attention}

\coursec{Compréhension du texte}

Après une première lecture silencieuse puis une lecture à voix haute du texte \all{„Cybermobbing – ein Fall aus Berlin“}, nous passons maintenant à l'étape décisive de toute leçon de texte : la \cle{compréhension}. L'objectif est double : vérifier que le sens global est acquis, puis aller chercher dans le texte des informations précises, en apprenant à formuler des réponses complètes et correctes en allemand. C'est exactement la démarche exigée au baccalauréat : lire, repérer, reformuler.

Pour faciliter le travail, rappelons le contenu du texte support avec ses numéros de ligne, indispensables pour les justifications.

\begin{textebox}[Rappel : „Cybermobbing – ein Fall aus Berlin“ (texte support avec numéros de ligne)]
1. In Berlin wurde ein 15-jähriger Schüler wochenlang im Internet gemobbt.
2. In einer WhatsApp-Gruppe seiner Klasse wurden hässliche Fotos von ihm verbreitet
3. und beleidigende Kommentare geschrieben.
4. Der Junge wurde ausgelacht und ausgegrenzt;
5. er wollte nicht mehr zur Schule gehen und konnte nachts nicht schlafen.
6. Schließlich wurde die Schulleitung informiert, und die Eltern erstatteten Anzeige bei der Polizei.
7. Der Täter, ein Mitschüler, wurde schnell gefunden,
8. weil die Nachrichten gespeichert worden waren.
9. Er wurde von der Schule suspendiert und musste sich entschuldigen.
10. Experten raten: Beweise sichern, das Passwort schützen,
11. die Täter blockieren und sich an eine Vertrauensperson wenden.
12. Niemand soll allein schweigen.
\end{textebox}

\textbf{Glossaire :} \all{mobben} (harceler) ; \all{die WhatsApp-Gruppe, -n} (le groupe WhatsApp) ; \all{verbreiten} (diffuser, répandre) ; \all{ausgrenzen} (exclure, mettre à l'écart) ; \all{Anzeige erstatten} (porter plainte) ; \all{der Täter, -} (l'auteur des faits, le coupable) ; \all{der Beweis, -e} (la preuve) ; \all{sich wenden an + Akk.} (se tourner vers, s'adresser à) ; \all{die Vertrauensperson, -en} (la personne de confiance) ; \all{schweigen} (se taire).

\courssub{Compréhension globale : de quoi parle le texte ?}

Avant les questions de détail, on pose toujours les questions générales. Elles portent sur le \cle{thème}, le \cle{lieu}, les \cle{personnages} et le \cle{type de texte}.

\begin{exemplebox}[Questions globales et réponses modèles]
\begin{itemize}
\item \all{Worum geht es in dem Text?} (De quoi parle le texte ?)\\
\all{Der Text handelt von einem Fall von Cybermobbing an einer Schule in Berlin.}\\
« Le texte traite d'un cas de harcèlement numérique dans une école à Berlin. »
\item \all{Wo spielt die Handlung?} (Où se passe l'action ?)\\
\all{Die Handlung spielt in Berlin, an einer Schule und in einer WhatsApp-Gruppe.}\\
« L'action se déroule à Berlin, dans une école et dans un groupe WhatsApp. »
\item \all{Wer sind die Personen?} (Qui sont les personnages ?)\\
\all{Die Personen sind ein 15-jähriger Schüler (das Opfer), ein Mitschüler (der Täter), die Eltern, die Schulleitung und die Polizei.}\\
« Les personnages sont un élève de 15 ans (la victime), un camarade de classe (l'auteur des faits), les parents, la direction de l'école et la police. »
\end{itemize}
\end{exemplebox}

\begin{methode}[Comment répondre à une question de compréhension]
\begin{enumerate}
\item \textbf{Lire la question} et souligner le mot interrogatif : \all{wer} (qui), \all{was} (quoi), \all{wo} (où), \all{warum} (pourquoi), \all{welche} (quel/le).
\item \textbf{Repérer la ligne} du texte qui contient la réponse.
\item \textbf{Reprendre les mots de la question} dans la réponse, puis compléter avec l'information du texte, sans recopier toute la phrase.
\item \textbf{Rédiger une phrase complète} : sujet + verbe conjugué en 2\textsuperscript{e} position + compléments.
\item \textbf{Citer la ligne} pour justifier : \all{(vgl. Zeile 5)} = « cf. ligne 5 ».
\end{enumerate}
\end{methode}

\begin{attention}[Erreurs fréquentes des francophones]
\begin{itemize}
\item Ne répondez \textbf{jamais} par \all{ja} ou \all{nein} seul à une question ouverte (\all{was, wer, warum...}) : il faut une phrase complète.
\item Respectez la \textbf{place du verbe conjugué en 2\textsuperscript{e} position} : \all{Hässliche Fotos wurden verbreitet} et non \all{Hässliche Fotos verbreitet wurden}.
\item Ne recopiez pas intégralement le texte : reformulez en gardant les mots-clés.
\item Attention au passif : \all{wurden verbreitet} = « ont été diffusées » (et non « diffusaient »).
\end{itemize}
\end{attention}

\courssub{Compréhension détaillée : questions et réponses}

Voici les questions détaillées, avec des réponses rédigées en phrases complètes. Les numéros de lignes renvoient au texte support ci-dessus.

\begin{exemplebox}[Questions détaillées et réponses modèles]
\begin{itemize}
\item \all{Was wurde in der WhatsApp-Gruppe verbreitet?} (Qu'est-ce qui a été diffusé dans le groupe WhatsApp ?)\\
\all{In der WhatsApp-Gruppe wurden hässliche Fotos des Schülers verbreitet (vgl. Zeile 2).}\\
« Dans le groupe WhatsApp, de vilaines photos de l'élève ont été diffusées. »
\item \all{Was wurde neben den Fotos noch geschrieben?} (Qu'a-t-on écrit d'autre que les photos ?)\\
\all{Es wurden auch beleidigende Kommentare geschrieben (vgl. Zeile 3).}\\
« Des commentaires insultants ont également été écrits. »
\item \all{Warum wollte der Schüler nicht mehr zur Schule gehen?} (Pourquoi l'élève ne voulait-il plus aller à l'école ?)\\
\all{Der Schüler wollte nicht mehr zur Schule gehen, weil er ausgelacht und ausgegrenzt wurde (vgl. Zeile 4--5).}\\
« L'élève ne voulait plus aller à l'école parce qu'on se moquait de lui et qu'il était mis à l'écart. »
\item \all{Wer war der Täter?} (Qui était l'auteur des faits ?)\\
\all{Der Täter war ein Mitschüler des Opfers (vgl. Zeile 7).}\\
« L'auteur des faits était un camarade de classe de la victime. »
\item \all{Warum konnte der Täter gefunden werden?} (Pourquoi l'auteur a-t-il pu être retrouvé ?)\\
\all{Der Täter konnte gefunden werden, weil die Nachrichten gespeichert worden waren (vgl. Zeile 8).}\\
« L'auteur a pu être retrouvé parce que les messages avaient été conservés. »
\item \all{Welche Strafe bekam er?} (Quelle punition a-t-il reçue ?)\\
\all{Er wurde von der Schule suspendiert und musste sich entschuldigen (vgl. Zeile 9).}\\
« Il a été suspendu de l'école et a dû présenter des excuses. »
\item \all{Was raten die Experten?} (Que conseillent les experts ?)\\
\all{Die Experten raten, Beweise zu sichern, das Passwort zu schützen, die Täter zu blockieren und sich an eine Vertrauensperson zu wenden (vgl. Zeile 10--11).}\\
« Les experts conseillent de conserver les preuves, de protéger son mot de passe, de bloquer les harceleurs et de s'adresser à une personne de confiance. »
\end{itemize}
\end{exemplebox}

\begin{center}
\begin{tabularx}{\linewidth}{|X|X|}
\hline
\rowcolor{popblueL} \textbf{Question} & \textbf{Élément de réponse (ligne du texte)} \\
\hline
\all{Was wurde verbreitet?} & \all{hässliche Fotos} (Zeile 2) \\
\hline
\all{Was wurde geschrieben?} & \all{beleidigende Kommentare} (Zeile 3) \\
\hline
\all{Warum wollte er nicht mehr zur Schule gehen?} & \all{Er wurde ausgelacht und ausgegrenzt.} (Zeile 4) \\
\hline
\all{Wer war der Täter?} & \all{ein Mitschüler} (Zeile 7) \\
\hline
\all{Warum wurde er gefunden?} & \all{Nachrichten waren gespeichert worden} (Zeile 8) \\
\hline
\all{Welche Strafe bekam er?} & \all{Suspendierung + Entschuldigung} (Zeile 9) \\
\hline
\all{Was raten die Experten?} & \all{Beweise sichern, Passwort schützen, blockieren, sich wenden an} (Zeile 10--11) \\
\hline
\end{tabularx}
\end{center}

\begin{popfigure}
\begin{center}
\begin{tikzpicture}[mindmap, grow cyclic, every node/.style={concept, align=center, font=\small},
root concept/.append style={concept color=popblue, font=\bfseries},
level 1/.append style={level distance=3.2cm, sibling angle=90},
level 1 concept/.append style={concept color=poporangeL}]
\node[root concept, align=center]{Der Text\\ \all{„Cybermobbing“}}
  child { node[align=center]{Wer?\\ das Opfer (15)\\ der Täter (Mitschüler)} }
  child { node[align=center]{Was?\\ hässliche Fotos\\ Beleidigungen} }
  child { node[align=center]{Wo?\\ Berlin\\ WhatsApp-Gruppe} }
  child { node[align=center]{Folgen?\\ Anzeige\\ Suspendierung\\ Ratschläge} };
\end{tikzpicture}
\end{center}
\legende{Carte mentale de compréhension : qui, quoi, où, conséquences.}
\end{popfigure}

\courssub{Vrai ou faux : justifier par une citation du texte}

L'exercice \all{Richtig oder Falsch} est classique au baccalauréat. La règle d'or : \textbf{toute réponse doit être justifiée par une citation exacte du texte} (entre guillemets allemands „…“), avec le numéro de ligne.

\begin{exoresolu}[Richtig oder Falsch ?]
\textbf{Énoncé.} Cochez \all{richtig} (R) ou \all{falsch} (F) et justifiez par une citation du texte.

1. \all{Der Schüler wurde in der Schule mit dem Stock geschlagen.}\\
2. \all{Die Eltern erstatteten Anzeige bei der Polizei.}\\
3. \all{Der Täter wurde nie gefunden.}\\
4. \all{Die Experten raten, allein zu schweigen.}
\tcblower
\textbf{Solution détaillée.}

1. \textbf{Falsch.} Le texte ne parle pas de coups physiques : \all{„In einer WhatsApp-Gruppe wurden hässliche Fotos von ihm verbreitet“} (Zeile 2). Le harcèlement était numérique, pas physique.

2. \textbf{Richtig.} \all{„Die Eltern erstatteten Anzeige bei der Polizei“} (Zeile 6). C'est une citation presque littérale du texte.

3. \textbf{Falsch.} \all{„Der Täter wurde schnell gefunden, weil die Nachrichten gespeichert worden waren“} (Zeile 7--8). Le coupable a été identifié grâce aux messages conservés.

4. \textbf{Falsch.} Le texte dit exactement le contraire : \all{„Niemand soll allein schweigen“} (Zeile 12). Les experts conseillent au contraire de parler et de se tourner vers une personne de confiance.
\end{exoresolu}

\begin{propriete}[Règle de justification]
Pour justifier une réponse \all{richtig} ou \all{falsch} :
\begin{itemize}
\item citez la phrase du texte entre guillemets allemands : \all{„…“} ;
\item indiquez la ligne : \all{(Zeile 6)} ou \all{(vgl. Zeile 6)} ;
\item si la réponse est \all{falsch}, la citation doit montrer la \textbf{contradiction} avec l'affirmation.
\end{itemize}
\end{propriete}

\courssub{Recherche d'équivalents : traduire à partir du texte}

On demande souvent de retrouver dans le texte l'équivalent allemand d'une expression française. Méthode : repérer d'abord l'idée, puis vérifier la construction exacte (verbe + cas + préposition).

\begin{center}
\begin{tabularx}{\linewidth}{|X|X|X|}
\hline
\rowcolor{popblueL} \textbf{Français} & \textbf{Équivalent dans le texte} & \textbf{Remarque} \\
\hline
porter plainte & \all{Anzeige erstatten (bei der Polizei)} & \all{die Anzeige, -n} ; préposition \all{bei} \\
\hline
les preuves & \all{die Beweise} & \all{der Beweis, -e} ; pluriel en \all{-e} \\
\hline
se tourner vers & \all{sich wenden an + Akkusativ} & verbe réfléchi ; \all{an eine Vertrauensperson} \\
\hline
diffuser, répandre & \all{verbreiten} & verbe régulier : \all{verbreitete, hat verbreitet} \\
\hline
protéger & \all{schützen} & \all{das Passwort schützen} \\
\hline
mettre à l'écart & \all{ausgrenzen} & particule séparable : \all{er grenzt aus} \\
\hline
\end{tabularx}
\end{center}

\begin{attention}[Faux amis et pièges de traduction]
\begin{itemize}
\item « Porter plainte » ne se traduit \textbf{pas} par \all{eine Klage tragen} : on dit \all{Anzeige erstatten} ou \all{eine Anzeige machen}.
\item \all{sich wenden an} exige l'\textbf{accusatif} : \all{sich an eine Vertrauensperson wenden} (et non \all{einer Vertrauensperson}).
\item \all{die Anzeige} signifie aussi « la petite annonce » (journal) : attention au contexte !
\item \all{der Täter} = l'auteur d'une infraction ; ne pas confondre avec \all{die Tat} (l'acte, l'action).
\end{itemize}
\end{attention}

\courssub{Entraînement : à vous de répondre}

\begin{exercice}[Fragen zum Text]
Répondez par des phrases complètes en allemand, sans recopier le texte, et citez la ligne.

1. \all{Wie alt ist das Opfer?}\\
2. \all{Was wurde neben den Fotos noch geschrieben?}\\
3. \all{Warum konnte der Täter gefunden werden?}\\
4. \all{An wen soll man sich wenden?}\\
5. Traduisez en allemand : « Les parents ont porté plainte parce que leur fils était harcelé. »
\end{exercice}

\begin{corrige}[Exercice : Fragen zum Text]
1. \all{Das Opfer ist 15 Jahre alt (vgl. Zeile 1).} — Reprise des mots de la question + information du texte.\\
2. \all{Es wurden auch beleidigende Kommentare geschrieben (vgl. Zeile 3).} — Passif au prétérit : \all{wurden geschrieben}.\\
3. \all{Der Täter konnte gefunden werden, weil die Nachrichten gespeichert worden waren (vgl. Zeile 8).} — Subordonnée de cause avec \all{weil} : verbe conjugué en fin de phrase.\\
4. \all{Man soll sich an eine Vertrauensperson wenden (vgl. Zeile 11).} — \all{sich wenden an + Akk.}\\
5. \all{Die Eltern erstatteten Anzeige, weil ihr Sohn gemobbt wurde.} — « porter plainte » = \all{Anzeige erstatten} ; subordonnée \all{weil} + verbe en position finale ; passif prétérit \all{wurde gemobbt}.
\end{corrige}

\begin{aretenir}[Les cinq réflexes de la compréhension de texte]
\begin{enumerate}
\item Identifier le thème, le lieu et les personnages avant le détail.
\item Souligner le mot interrogatif de chaque question.
\item Reprendre les mots de la question dans la réponse.
\item Rédiger une phrase complète, verbe conjugué en 2\textsuperscript{e} position.
\item Justifier par une citation entre „…“ avec le numéro de ligne.
\end{enumerate}
\end{aretenir}

Cette étape de compréhension étant acquise, nous pouvons maintenant nous approprier les mots du texte : la section suivante sera consacrée au \cle{lexique thématique} du harcèlement numérique (\all{das Cybermobbing, das Opfer, der Täter, die Anzeige, schützen, das Passwort}), avec ses familles de mots, ses articles et ses pluriels.

\coursec{Lexique thématique : das Cybermobbing}

Après avoir compris le texte \all{Cybermobbing – ein Fall aus Berlin} dans sa globalité puis dans le détail, arrêtons-nous maintenant sur ses mots-clés. Un texte ne se possède vraiment que si l'on maîtrise son vocabulaire : c'est lui qui permet de répondre aux questions, de résumer et de produire à son tour. Nous allons donc organiser le lexique du harcèlement numérique en quatre champs sémantiques (les personnes, les actes, la protection, les lieux et outils), puis étudier les familles de mots et les expressions figées. Pour chaque nom, mémorisez toujours trois choses ensemble : \cle{l'article}, \cle{le pluriel} et \cle{un exemple en contexte}.

\courssub{Observation : le vocabulaire dans un texte}

Lisons d'abord ce court texte original, qui reprend la situation de notre texte support sous un autre angle : un article de journal scolaire.

\begin{textebox}[Aus der Schülerzeitung „Die Pausenklingel“, Berlin]
\textbf{Cybermobbing an der Kant-Schule: Eine Mitschülerin wehrt sich}

Berlin. Seit Wochen wurde eine 16-jährige Schülerin der Kant-Schule in einer WhatsApp-Gruppe beleidigt und bedroht. Die Täter, zwei Mitschüler aus ihrer Klasse, verbreiteten gefälschte Fotos von ihr im Internet und in sozialen Netzwerken. Das Opfer schwieg zuerst, denn es hatte Angst vor den Reaktionen der anderen.

Ein Zeuge, ein Freund aus der Parallelklasse, beobachtete die Angriffe und speicherte alle Nachrichten als Beweise auf seinem Handy. Dann informierte er die Klassenlehrerin. „Man darf nicht wegsehen“, sagte er. „Wer schweigt, hilft den Tätern.“

Die Schule reagierte sofort: Die beleidigenden Inhalte wurden gelöscht, und die Eltern der Täter wurden informiert. Das Opfer wandte sich an eine Beratungsstelle und wurde dort geschützt und unterstützt. Schließlich erstattete die Familie Anzeige bei der Polizei. Die Täter wurden angezeigt und müssen jetzt mit einer Strafe rechnen.

Die Direktorin erinnert alle Schüler an wichtige Regeln: „Ändert regelmäßig euer Passwort, es muss geheim bleiben. Meldet verdächtige Nachrichten sofort. Und denkt daran: Mobbing ist kein Spaß, sondern eine Tat.“

Seit diesem Fall organisiert die Schule jedes Jahr eine Projektwoche gegen Cybermobbing. Viele Schüler sagen: „Jetzt wissen wir, wie wir uns und unsere Freunde schützen können.“
\end{textebox}

\textbf{Glossaire :} \all{sich wehren} (se défendre) ; \all{gefälscht} (falsifié, truqué) ; \all{wegsehen} (détourner le regard) ; \all{die Beratungsstelle (-n)} (le centre d'aide/d'écoute) ; \all{unterstützen} (soutenir) ; \all{mit einer Strafe rechnen} (s'attendre à une sanction) ; \all{verdächtig} (suspect) ; \all{die Projektwoche (-n)} (la semaine de projet).

\textbf{Questions de compréhension :}
\begin{enumerate}
\item \all{Wer wurde in der WhatsApp-Gruppe beleidigt?} (Qui a été insulté dans le groupe WhatsApp ?)
\item \all{Was hat der Zeuge gespeichert?} (Qu'est-ce que le témoin a enregistré ?)
\item \all{An wen hat sich das Opfer gewandt?} (À qui la victime s'est-elle adressée ?)
\item \all{Was empfiehlt die Direktorin?} (Que recommande la directrice ?)
\end{enumerate}

\courssub{Champ 1 : les personnes (die Personen)}

\begin{definition}[Les personnes autour du harcèlement]
\begin{itemize}
\item \all{das Opfer (-)} : la victime. \emph{Attention : pluriel identique au singulier.}
\item \all{der Täter (-)} : l'auteur (des faits), le coupable. Féminin : \all{die Täterin (-nen)}.
\item \all{der Mitschüler (-)} : le camarade de classe (masculin) ; \all{die Mitschülerin (-nen)} : la camarade.
\item \all{der Zeuge (-n)} : le témoin. \emph{Nom faible :} \all{der Zeuge, den Zeugen, dem Zeugen, des Zeugen}.
\item \all{die Polizei (nur Singular)} : la police (toujours au singulier en allemand).
\end{itemize}
\end{definition}

\begin{exemplebox}[Les personnes en contexte]
\all{Das Opfer wurde von zwei Mitschülern bedroht.} « La victime a été menacée par deux camarades de classe. »

\all{Ein Zeuge hat alles gesehen und die Polizei informiert.} « Un témoin a tout vu et a informé la police. »

\all{Die Täter müssen mit einer Strafe rechnen.} « Les auteurs doivent s'attendre à une sanction. »
\end{exemplebox}

\begin{attention}[der Täter n'est pas « le traître »]
\all{der Täter} signifie \cle{l'auteur d'un acte, le coupable} (celui qui \emph{fait} la chose : \all{tun} → \all{die Tat} → \all{der Täter}). Le « traître » se dit \all{der Verräter}. De même, \all{die Polizei} est singulier : on dit \all{Die Polizei \emph{wurde} informiert} (la police a été informée), jamais de pluriel.
\end{attention}

\courssub{Champ 2 : les actes (die Handlungen)}

Voici les verbes essentiels, avec leurs formes principales (infinitif – prétérit – participe II) :

\begin{center}
\begin{tabularx}{\linewidth}{|l|l|X|}
\hline
\rowcolor{popblueL} \textbf{Deutsch} & \textbf{Français} & \textbf{Beispiel aus dem Text} \\
\hline
\all{beleidigen (beleidigte, beleidigt)} & insulter & \all{Die Schülerin wurde beleidigt.} \\
\hline
\all{bedrohen (bedrohte, bedroht)} & menacer & \all{Die Täter bedrohten das Opfer.} \\
\hline
\all{verbreiten (verbreitete, verbreitet)} & diffuser, répandre & \all{Sie verbreiteten Fotos im Internet.} \\
\hline
\all{melden (meldete, gemeldet)} & signaler & \all{Meldet verdächtige Nachrichten!} \\
\hline
\all{löschen (löschte, gelöscht)} & effacer, supprimer & \all{Die Inhalte wurden gelöscht.} \\
\hline
\all{speichern (speicherte, gespeichert)} & enregistrer, sauvegarder & \all{Der Zeuge speicherte die Beweise.} \\
\hline
\all{anzeigen (zeigte an, angezeigt)} & dénoncer (à la police) & \all{Die Täter wurden angezeigt.} \\
\hline
\all{mobben (mobbte, gemobbt)} & harceler & \all{Er mobbt seine Mitschülerin.} \\
\hline
\end{tabularx}
\end{center}

\begin{propriete}[Le verbe séparable anzeigen]
\all{anzeigen} est un \cle{verbe à particule séparable} : la particule \all{an-} va en fin de phrase au présent et au prétérit, et le participe II s'écrit \all{an\emph{ge}zeigt} (le \all{ge-} s'insère entre la particule et le radical).

\all{Er \textbf{zeigt} den Täter \textbf{an}.} « Il dénonce l'auteur. » (présent)

\all{Er \textbf{zeigte} den Täter \textbf{an}.} « Il a dénoncé l'auteur. » (prétérit)

\all{Er hat den Täter \textbf{angezeigt}.} « Il a dénoncé l'auteur. » (parfait)

Au passif : \all{Der Täter wurde angezeigt.} « L'auteur a été dénoncé. »
\end{propriete}

\begin{attention}[Pièges de conjugaison pour les francophones]
\begin{itemize}
\item \all{beleidigen} et \all{bedrohen} sont des \cle{verbes faibles réguliers} : pas de changement de voyelle. Ne les conjuguez pas comme des verbes forts (pas d'umlaut) : la forme est \all{beleidigte, hat beleidigt}.
\item \all{mobben} double le \all{b} au participe : \all{gemobbt}.
\item En français on « signale quelque chose » ; en allemand : \all{etwas (Akk.) melden}, sans préposition.
\end{itemize}
\end{attention}

\courssub{Champ 3 : la protection (der Schutz)}

\begin{definition}[Se protéger et réagir]
\begin{itemize}
\item \all{schützen (schützte, geschützt)} : protéger ; \all{sich schützen} : se protéger.
\item \all{das Passwort (¨-er)} : le mot de passe. Pluriel : \all{die Passwörter}.
\item \all{geheim} (adjectif) : secret. \all{geheim bleiben} : rester secret.
\item \all{die Anzeige (-n)} : la plainte, la dénonciation.
\item \all{der Beweis (-e)} : la preuve.
\item \all{die Strafe (-n)} : la sanction, la peine.
\item \all{sich wenden an + Akkusativ} : s'adresser à quelqu'un.
\end{itemize}
\end{definition}

\begin{exemplebox}[Exemples traduits à retenir]
\all{Das Opfer wurde geschützt.} « La victime a été protégée. »

\all{Er erstattete Anzeige.} « Il a porté plainte. »

\all{Das Passwort muss geheim bleiben.} « Le mot de passe doit rester secret. »

\all{Der Zeuge speicherte die Beweise auf seinem Handy.} « Le témoin a enregistré les preuves sur son portable. »

\all{Sie wandte sich an eine Beratungsstelle.} « Elle s'est adressée à un centre d'aide. »
\end{exemplebox}

\begin{attention}[die Anzeige n'est pas « l'annonce »]
\all{die Anzeige} signifie ici \cle{la plainte} (déposée auprès de la police). L'« annonce publicitaire » se dit \all{die Werbung (-en)}. Expression figée à mémoriser : \all{\textbf{eine Anzeige erstatten}} = porter plainte (et non \all{\emph{eine Anzeige machen}}, même si l'on entend parfois cette forme familière). Autres expressions figées : \all{\textbf{jemanden mobben}} (harceler quelqu'un), \all{\textbf{Beweise speichern}} (conserver des preuves), \all{\textbf{jemanden bedrohen}} (menacer quelqu'un).
\end{attention}

\courssub{Champ 4 : les lieux et outils (Orte und Werkzeuge)}

\begin{definition}[Où se passe le Cybermobbing ?]
\begin{itemize}
\item \all{das Handy (-s)} : le téléphone portable. \emph{Mot allemand inventé : les Anglophones disent « mobile phone » !}
\item \all{das soziale Netzwerk (-e)} : le réseau social.
\item \all{die WhatsApp-Gruppe (-n)} : le groupe WhatsApp.
\item \all{das Internet (-s, meist ohne Artikel)} : Internet. \all{im Internet} : sur Internet.
\end{itemize}
\end{definition}

\begin{exemplebox}[Les outils en contexte]
\all{Die Fotos wurden in sozialen Netzwerken verbreitet.} « Les photos ont été diffusées sur les réseaux sociaux. »

\all{In der WhatsApp-Gruppe wurde sie jeden Tag beleidigt.} « Dans le groupe WhatsApp, elle était insultée chaque jour. »

\all{Viele Jugendliche in Yaoundé und in Berlin benutzen jeden Tag ihr Handy.} « Beaucoup de jeunes à Yaoundé et à Berlin utilisent leur portable chaque jour. »
\end{exemplebox}

\courssub{Les familles de mots (Wortfamilien)}

Apprendre les mots par familles multiplie votre vocabulaire sans effort : un seul radical, plusieurs mots.

\begin{center}
\begin{tabularx}{\linewidth}{|l|l|X|}
\hline
\rowcolor{popgreenL} \textbf{Famille} & \textbf{Membres} & \textbf{Français} \\
\hline
TUN & \all{tun – die Tat – der Täter} & faire – l'acte – l'auteur \\
\hline
ANZEIGEN & \all{anzeigen – die Anzeige} & dénoncer – la plainte \\
\hline
SCHÜTZEN & \all{schützen – der Schutz} & protéger – la protection \\
\hline
BELEIDIGEN & \all{beleidigen – die Beleidigung} & insulter – l'insulte \\
\hline
BEDROHEN & \all{bedrohen – die Bedrohung} & menacer – la menace \\
\hline
BEWEISEN & \all{beweisen – der Beweis} & prouver – la preuve \\
\hline
\end{tabularx}
\end{center}

\begin{propriete}[Verbe → nom d'action en -ung]
De nombreux verbes allemands forment un nom d'action féminin en \all{-ung} (toujours \all{die}, pluriel \all{-en}) : \all{beleidigen → die Beleidigung}, \all{bedrohen → die Bedrohung}, \all{melden → die Meldung}. C'est l'équivalent des noms français en « -tion/-ment » (insulte, menace, signalement). Attention aux exceptions : \all{schützen → der Schutz} (masculin, sans \all{-ung}) et \all{tun → die Tat}.
\end{propriete}

\begin{popfigure}
\begin{tikzpicture}[mindmap, grow cyclic, every node/.style={concept, align=center, font=\small}, root concept/.append style={concept color=popblue, text=white, font=\bfseries}, level 1/.append style={level distance=3.4cm, sibling angle=90}, level 2/.append style={level distance=2.2cm, sibling angle=50}]
\node{CYBERMOBBING}
  child[concept color=poporange] { node {Personen}
    child { node[concept color=poporangeL, text=popdark, align=center]{das Opfer\\der Täter} }
    child { node[concept color=poporangeL, text=popdark, align=center]{der Zeuge\\die Polizei} }
  }
  child[concept color=poppink] { node {Handlungen}
    child { node[concept color=poppinkL, text=popdark, align=center]{beleidigen\\bedrohen} }
    child { node[concept color=poppinkL, text=popdark, align=center]{verbreiten\\mobben} }
  }
  child[concept color=popgreen] { node {Schutz}
    child { node[concept color=popgreenL, text=popdark, align=center]{schützen\\das Passwort} }
    child { node[concept color=popgreenL, text=popdark, align=center]{die Anzeige\\der Beweis} }
  }
  child[concept color=popgold] { node {Orte}
    child { node[concept color=popgoldL, text=popdark, align=center]{das Handy\\das Internet} }
    child { node[concept color=popgoldL, text=popdark, align=center]{die WhatsApp-\\Gruppe} }
  };
\end{tikzpicture}
\legende{Carte mentale du champ lexical de \all{das Cybermobbing} : quatre champs à mémoriser.}
\end{popfigure}

\courssub{Tableau récapitulatif général}

\begin{center}
\begin{tabularx}{\linewidth}{|l|l|X|}
\hline
\rowcolor{popblueL} \textbf{Deutsch (Artikel + Plural)} & \textbf{Français} & \textbf{Beispiel} \\
\hline
\all{das Opfer (-)} & la victime & \all{Das Opfer schwieg zuerst.} \\
\hline
\all{der Täter (-)} & l'auteur, le coupable & \all{Die Täter wurden angezeigt.} \\
\hline
\all{der Zeuge (-n)} & le témoin & \all{Der Zeuge speicherte Beweise.} \\
\hline
\all{die Polizei (Sg.)} & la police & \all{Die Polizei wurde informiert.} \\
\hline
\all{die Anzeige (-n)} & la plainte & \all{Er erstattete Anzeige.} \\
\hline
\all{der Beweis (-e)} & la preuve & \all{Beweise speichern!} \\
\hline
\all{die Strafe (-n)} & la sanction & \all{Sie müssen mit einer Strafe rechnen.} \\
\hline
\all{das Passwort (¨-er)} & le mot de passe & \all{Das Passwort bleibt geheim.} \\
\hline
\all{das Handy (-s)} & le portable & \all{auf dem Handy speichern} \\
\hline
\all{das soziale Netzwerk (-e)} & le réseau social & \all{Fotos im Netzwerk verbreiten} \\
\hline
\all{die WhatsApp-Gruppe (-n)} & le groupe WhatsApp & \all{in der Gruppe beleidigt} \\
\hline
\end{tabularx}
\end{center}

\begin{aretenir}[Les expressions figées à connaître par cœur]
\begin{itemize}
\item \all{eine Anzeige erstatten} : porter plainte
\item \all{jemanden mobben / beleidigen / bedrohen} : harceler / insulter / menacer quelqu'un
\item \all{Beweise speichern} : conserver des preuves
\item \all{sich wenden an + Akk.} : s'adresser à
\item \all{mit einer Strafe rechnen} : s'attendre à une sanction
\item \all{geheim bleiben} : rester secret
\end{itemize}
\end{aretenir}

\begin{exoresolu}[Exercice de réemploi lexical]
\textbf{Énoncé.} Complétez les phrases avec le mot qui convient : \all{Opfer – Täter – Anzeige – Beweise – Passwort – gelöscht – wenden}.

\begin{enumerate}
\item \all{Die Familie erstattete \dots\ bei der Polizei.}
\item \all{Der Zeuge speicherte alle Nachrichten als \dots .}
\item \all{Das \dots\ wurde von zwei Mitschülern bedroht.}
\item \all{Die beleidigenden Fotos wurden sofort \dots .}
\item \all{Dein \dots\ muss geheim bleiben!}
\item \all{Die \dots\ müssen mit einer Strafe rechnen.}
\item \all{Sie können sich an eine Beratungsstelle \dots .}
\end{enumerate}

\tcblower
\textbf{Solution détaillée.}
\begin{enumerate}
\item \all{Die Familie erstattete \textbf{Anzeige} bei der Polizei.} (expression figée : \all{eine Anzeige erstatten} = porter plainte)
\item \all{Der Zeuge speicherte alle Nachrichten als \textbf{Beweise}.} (pluriel de \all{der Beweis} : \all{die Beweise})
\item \all{Das \textbf{Opfer} wurde von zwei Mitschülern bedroht.} (c'est la victime qui subit la menace ; neutre : \all{das Opfer})
\item \all{Die beleidigenden Fotos wurden sofort \textbf{gelöscht}.} (passif prétérit : \all{wurden + Partizip II} ; \all{löschen → gelöscht})
\item \all{Dein \textbf{Passwort} muss geheim bleiben!} (\all{das Passwort} ; possessif \all{dein} sans terminaison devant un neutre)
\item \all{Die \textbf{Täter} müssen mit einer Strafe rechnen.} (ce sont les auteurs qui risquent une sanction ; pluriel identique : \all{die Täter})
\item \all{Sie können sich an eine Beratungsstelle \textbf{wenden}.} (\all{sich wenden an + Akk.} : l'infinitif se place en fin de phrase après le modal \all{können})
\end{enumerate}
\end{exoresolu}

\begin{savaistu}[D'où vient le mot « mobben » ?]
Le verbe \all{mobben} vient de l'anglais \emph{mob}, « la foule » (en latin \emph{mobile vulgus}, « la foule changeante ») : le harcèlement, c'est la violence d'un groupe contre une personne seule. Le mot composé \all{das Cybermobbing} est utilisé dans toute l'aire germanophone : en Allemagne, en Autriche et en Suisse. En Suisse alémanique, on entend aussi \all{Cyber-Bullying}, emprunt direct à l'anglais. Et petit clin d'œil linguistique : \all{das Handy} est un « faux anglicisme » inventé en Allemagne — un Anglais dira \emph{mobile phone} ou \emph{smartphone} !
\end{savaistu}

\begin{attention}[Dernières erreurs à bannir]
\begin{itemize}
\item Ne traduisez pas « la victime » par \all{*die Opfer*} : c'est \all{das Opfer} (neutre), même pour une fille : \all{Das Opfer, eine 16-jährige Schülerin, \dots}.
\item \all{der Zeuge} est un nom faible : \all{Ich habe den Zeuge\textbf{n} gesehen.}
\item Pluriel de \all{das Passwort} : \all{die Passw\textbf{ö}rter} (avec umlaut), comme \all{das Wort → die Wörter}.
\item Ordre des mots avec les séparables : \all{Er zeigt den Täter \textbf{an}} — la particule en toute fin, jamais collée au verbe conjugué.
\end{itemize}
\end{attention}

\begin{exercice}[À vous de jouer]
Traduisez en allemand, en utilisant le lexique de la leçon :
\begin{enumerate}
\item La victime a porté plainte auprès de la police.
\item Les preuves ont été enregistrées sur le portable du témoin.
\item Ton mot de passe doit rester secret.
\item Les auteurs ont diffusé les photos sur les réseaux sociaux.
\item Adresse-toi à un centre d'aide !
\end{enumerate}
\end{exercice}

\begin{corrige}[Exercice « À vous de jouer »]
\begin{enumerate}
\item \all{Das Opfer erstattete Anzeige bei der Polizei.}
\item \all{Die Beweise wurden auf dem Handy des Zeugen gespeichert.}
\item \all{Dein Passwort muss geheim bleiben.}
\item \all{Die Täter verbreiteten die Fotos in sozialen Netzwerken.} (ou au passif : \all{Die Fotos wurden in sozialen Netzwerken verbreitet.})
\item \all{Wende dich an eine Beratungsstelle!} (impératif de \all{sich wenden}, pronom réfléchi \all{dich})
\end{enumerate}
\end{corrige}

Ce lexique est désormais votre boîte à outils : il vous servira dans la section suivante pour comprendre le passif (\all{Das Opfer \emph{wurde geschützt}}), puis pour produire votre propre résumé et votre courte production écrite sur le harcèlement numérique.

\coursec{Grammaire 1 : le passif (Vorgangspassiv)}

Après avoir enrichi notre vocabulaire autour du \all{Cybermobbing} — \all{das Opfer}, \all{der Täter}, \all{die Anzeige}, \all{schützen}, \all{das Passwort} — nous allons maintenant observer comment ces mots se comportent dans les phrases du texte. En effet, lorsqu'on parle de harcèlement numérique, on s'intéresse souvent plus à ce qui \emph{arrive} à la victime qu'à celui qui agit : on dit alors « les photos ont été diffusées », « l'auteur a été retrouvé ». Cette construction s'appelle le \cle{passif}. C'est l'une des structures les plus fréquentes de la presse germanophone et l'un des points de grammaire les plus rentables au baccalauréat.

\courssub{Observation : que devient le sujet ?}

Reprenons deux phrases tirées de notre texte support :

\begin{exemplebox}[Deux phrases du texte]
\all{Fotos wurden verbreitet.} « Des photos ont été diffusées. »

\all{Der Täter wurde gefunden.} « L'auteur a été retrouvé. »
\end{exemplebox}

Observons attentivement :

\begin{itemize}
\item Dans la première phrase, ce ne sont pas les photos qui agissent : elles \emph{subissent} l'action. Pourtant, elles occupent la place du \cle{sujet} grammatical (nominatif, avant le verbe conjugué).
\item Dans la deuxième phrase, l'auteur des faits est bien « retrouvé » : il ne fait rien, c'est la police qui agit, mais elle n'est même pas nommée.
\item Le verbe se présente en deux morceaux : une forme de \all{werden} conjuguée (\all{wurden}) en deuxième position, et un \cle{participe II} (\all{verbreitet}, \all{gefunden}) renvoyé \textbf{à la fin} de la phrase.
\end{itemize}

\begin{aretenir}[L'essentiel de l'observation]
Au passif, le sujet grammatical n'est \textbf{pas} celui qui agit : c'est celui qui \emph{subit} l'action. Celui qui agit (l'\cle{agent}) disparaît souvent, ou bien est introduit par \all{von} + datif. La forme verbale = \all{werden} conjugué + participe II en fin de phrase.
\end{aretenir}

\courssub{Formation du passif : werden + Partizip II}

\begin{propriete}[Règle de formation du Vorgangspassiv]
Le passif d'action (\all{Vorgangspassiv}) se forme avec :
\begin{enumerate}
\item l'auxiliaire \all{werden}, \textbf{conjugué} au temps voulu, en deuxième position ;
\item le \textbf{participe II} du verbe principal, placé \textbf{à la fin} de la phrase.
\end{enumerate}
L'agent (celui qui agit) est introduit par \all{von} + \textbf{datif} :

\all{Der Täter wurde von der Polizei gefunden.} « L'auteur a été retrouvé par la police. »

(\all{von der Polizei} = von + datif féminin ; \all{von dem Opfer} se contracte en \all{vom Opfer}.)
\end{propriete}

Comparons actif et passif sur un même événement :

\begin{center}
\begin{tabularx}{\linewidth}{|X|X|}
\hline
\rowcolor{popblueL} \textbf{Actif} & \textbf{Passif} \\
\hline
\all{Die Polizei findet den Täter.} (La police retrouve l'auteur.) & \all{Der Täter wird (von der Polizei) gefunden.} (L'auteur est retrouvé par la police.) \\
\hline
\all{Er verbreitet die Fotos.} (Il diffuse les photos.) & \all{Die Fotos werden (von ihm) verbreitet.} (Les photos sont diffusées par lui.) \\
\hline
\all{Man schützt das Opfer.} (On protège la victime.) & \all{Das Opfer wird geschützt.} (La victime est protégée.) \\
\hline
\end{tabularx}
\end{center}

\begin{popfigure}
\begin{tikzpicture}[font=\small]
\node[draw=popblue, fill=popblueL, rounded corners, align=center, minimum width=4.2cm, minimum height=1.6cm] (actif) at (0,0) {\textbf{ACTIF}\\ \all{Der Täter} $\rightarrow$ \all{verbreitet} $\rightarrow$ \all{die Fotos}\\ (sujet = agent)};
\node[draw=popgreen, fill=popgreenL, rounded corners, align=center, minimum width=4.2cm, minimum height=1.6cm] (passif) at (9,0) {\textbf{PASSIF}\\ \all{Die Fotos werden (von ihm) verbreitet.}\\ (sujet = celui qui subit)};
\draw[-{Stealth[length=4mm]}, line width=1.5pt, poporange] (actif) -- (passif) node[midway, above, align=center, text=popdark]{COD $\rightarrow$ sujet\\ agent $\rightarrow$ \all{von} + Dativ};
\end{tikzpicture}
\legende{Schéma de transformation actif $\Rightarrow$ passif : le complément d'objet direct devient sujet au nominatif ; l'agent est introduit par \all{von} + datif.}
\end{popfigure}

\courssub{Le passif aux temps : présent, prétérit, passé composé}

\begin{propriete}[Passif au présent et au prétérit]
\textbf{Passif présent} : \all{werden} au présent + participe II.

\all{Das Opfer wird geschützt.} « La victime est protégée. »

\textbf{Passif prétérit} : \all{werden} au prétérit + participe II.

\all{Das Opfer wurde geschützt.} « La victime a été / fut protégée. »
\end{propriete}

Le prétérit passif (\all{wurde} + participe II) est le temps normal du récit écrit : c'est celui du texte journalistique, comme dans notre texte support. Au passé composé (langue parlée), le passif se construit avec \all{sein} + participe II + \all{worden} :

\all{Er ist gefunden worden.} « Il a été retrouvé. »

Remarquez la forme spéciale \all{worden} (et non \all{geworden}) : c'est le participe II de \all{werden} \emph{uniquement} dans la construction passive.

\begin{center}
\begin{tabular}{|l|l|l|}
\hline
\rowcolor{popblueL} \textbf{Personne} & \textbf{Présent} & \textbf{Prétérit} \\
\hline
ich & werde & wurde \\
\hline
du & wirst & wurdest \\
\hline
er / sie / es & wird & wurde \\
\hline
wir & werden & wurden \\
\hline
ihr & werdet & wurdet \\
\hline
sie / Sie & werden & wurden \\
\hline
\end{tabular}
\end{center}

\begin{exemplebox}[Exemples traduits]
\all{Die Nachrichten werden gesendet.} « Les messages sont envoyés. » (présent)

\all{Die Nachrichten wurden gesendet.} « Les messages ont été envoyés. » (prétérit)

\all{Das Passwort wird geändert.} « Le mot de passe est changé. »

\all{Die Anzeige wurde von den Eltern gemacht.} « La plainte a été déposée par les parents. »

\all{Ihm wurde geholfen.} « On l'a aidé. » (verbe à datif : pas de sujet nominatif !)
\end{exemplebox}

\courssub{Cas particuliers : verbes à datif et passif impersonnel}

\begin{propriete}[Verbes à datif et passif impersonnel]
\textbf{1. Verbes à datif} (\all{helfen}, \all{danken}, \all{folgen}, \all{antworten}\ldots) : ils n'ont pas de COD, donc \textbf{rien} ne peut devenir sujet nominatif. Le complément \textbf{reste au datif} et le verbe se met à la 3\textsuperscript{e} personne du singulier :

\all{Man hilft ihm.} $\Rightarrow$ \all{Ihm wird geholfen.} « On l'aide. »

\textbf{2. Passif impersonnel} : sans aucun complément, le passif exprime une activité générale, avec \all{es} en tête (ou un complément de lieu) :

\all{Es wurde gelacht.} « On a ri. »

\all{Hier wird getanzt.} « Ici on danse. »
\end{propriete}

\begin{attention}[Piège : \all{werden} (action) ou \all{sein} (état) ?]
Ne confondez pas les deux « passifs » allemands :
\begin{itemize}
\item \all{Die Tür wird geschlossen.} = « On ferme la porte » : une \textbf{action} en train de se dérouler (\all{Vorgangspassiv}, avec \all{werden}).
\item \all{Die Tür ist geschlossen.} = « La porte est fermée » : un \textbf{état}, un résultat (\all{Zustandspassiv}, avec \all{sein}).
\end{itemize}
Au baccalauréat, c'est le \all{Vorgangspassiv} (avec \all{werden}) qui est demandé et valorisé. Le français « est + participe » est ambigu ; demandez-vous toujours : action ou état ?
\end{attention}

\begin{attention}[Piège de conjugaison : \all{wird}, pas *würd !]
Le prétérit \all{wurde} contient un \emph{u}, mais au présent la 3\textsuperscript{e} personne est \all{er wird} avec un \emph{i} — jamais d'umlaut : on écrit \all{er wird}, \all{du wirst}, et non *würd, *würdest. La voyelle change (e $\rightarrow$ i) au présent, mais sans tréma.
\end{attention}

\courssub{Méthode : transformer une phrase active en phrase passive}

\begin{methode}[De l'actif au passif en 4 étapes]
\begin{enumerate}
\item \textbf{Repérer le COD} de la phrase active (accusatif) : \all{Die Polizei findet \textbf{den Täter}.}
\item \textbf{Le mettre au nominatif} en tête de phrase : \all{den Täter} $\rightarrow$ \all{Der Täter}.
\item \textbf{Conjuguer \all{werden} au même temps} que le verbe actif, en l'accordant avec le nouveau sujet : \all{findet} (présent) $\rightarrow$ \all{wird}.
\item \textbf{Placer le participe II à la fin} et, si nécessaire, ajouter l'agent avec \all{von} + datif : \all{Der Täter wird (von der Polizei) gefunden.}
\end{enumerate}
Vérification : le verbe conjugué est en 2\textsuperscript{e} position, le participe II ferme la phrase.
\end{methode}

\begin{exoresolu}[Transformation actif $\rightarrow$ passif]
Énoncé : mettez au passif (même temps) : \all{Der Täter verbreitete die Fotos.} / \all{Die Lehrer schützen die Schüler.} / \all{Man half dem Opfer.}

\tcblower

Solution détaillée :

\textbf{1.} \all{Der Täter verbreitete die Fotos.} — COD : \all{die Fotos} (pluriel) $\rightarrow$ nominatif \all{die Fotos} ; \all{verbreitete} est au prétérit $\rightarrow$ \all{wurden} ; participe II \all{verbreitet} en fin. Résultat : \all{Die Fotos wurden (von dem Täter) verbreitet.} « Les photos ont été diffusées (par l'auteur). »

\textbf{2.} \all{Die Lehrer schützen die Schüler.} — COD : \all{die Schüler} ; présent $\rightarrow$ \all{werden} ; participe II \all{geschützt}. Résultat : \all{Die Schüler werden (von den Lehrern) geschützt.} « Les élèves sont protégés (par les professeurs). »

\textbf{3.} \all{Man half dem Opfer.} — \all{helfen} est un verbe à \textbf{datif} : pas de COD, donc pas de sujet nominatif ; \all{dem Opfer} reste au datif ; prétérit $\rightarrow$ \all{wurde} (singulier). Résultat : \all{Dem Opfer wurde geholfen.} « On a aidé la victime. »
\end{exoresolu}

\courssub{Entraînement sur un texte : le passif dans la presse}

Lisons ce court article original, dans lequel le passif est omniprésent — comme dans tout article de presse allemand sur la cybercriminalité.

\begin{textebox}[Cybermobbing an einer Schule in Berlin]
An einer Schule in Berlin wurde ein schwerer Fall von Cybermobbing entdeckt. Über mehrere Wochen wurden private Fotos einer Schülerin in einer Chatgruppe verbreitet. Die Fotos wurden mit beleidigenden Kommentaren versehen, und das Mädchen wurde täglich von unbekannten Nummern angerufen. Zuerst wurde der Fall von den Klassenkameraden nicht ernst genommen. Erst als die Schülerin nicht mehr zur Schule kam, wurde die Direktorin informiert. Sofort wurde eine Anzeige bei der Polizei gemacht. Das Handy der Schülerin wurde untersucht, und schnell wurde der Täter gefunden: ein Schüler aus der eigenen Klasse. Er wurde von der Schule für zwei Wochen suspendiert, und ein Gespräch mit einem Psychologen wurde angeordnet. Die Fotos wurden aus dem Internet gelöscht. In der Klasse wurde danach ein Projekt zum Thema „Respekt im Netz“ gestartet. Die Schüler wurden über die Gefahren des Cybermobbings informiert, und es wurden Regeln für den Umgang mit Handys vereinbart. „So etwas darf nie wieder passieren“, sagte die Direktorin. „Die Opfer müssen geschützt werden, und die Täter müssen bestraft werden.“
\end{textebox}

\textbf{Glossaire :} \all{entdecken} (découvrir) ; \all{die Chatgruppe, -n} (le groupe de discussion) ; \all{beleidigend} (insultant) ; \all{der Kommentar, -e} (le commentaire) ; \all{anrufen} (téléphoner à) ; \all{ernst nehmen} (prendre au sérieux) ; \all{die Direktorin, -nen} (la directrice) ; \all{untersuchen} (examiner) ; \all{suspendieren} (exclure temporairement) ; \all{anordnen} (ordonner) ; \all{löschen} (supprimer) ; \all{vereinbaren} (convenir) ; \all{der Umgang mit} (l'usage de) ; \all{bestrafen} (punir).

\textbf{Questions de compréhension :}
\begin{enumerate}
\item \all{Was wurde in einer Chatgruppe verbreitet?}
\item \all{Wann wurde die Direktorin informiert?}
\item \all{Von wem wurde der Täter gefunden?}
\item \all{Welche Maßnahmen wurden nach dem Fall getroffen?}
\item Relevez cinq formes passives du texte et identifiez leur temps.
\end{enumerate}

\begin{exercice}[À vous de jouer]
Mettez au passif, au temps indiqué :
\begin{enumerate}
\item \all{Die Polizei verhaftet den Täter.} (présent)
\item \all{Die Schüler vereinbarten neue Regeln.} (prétérit)
\item \all{Man informierte die Eltern.} (prétérit)
\item \all{Der Psychologe hilft dem Opfer.} (présent — attention au datif !)
\end{enumerate}
\end{exercice}

\begin{corrige}[Exercice « À vous de jouer »]
\begin{enumerate}
\item \all{Der Täter wird (von der Polizei) verhaftet.} « L'auteur est arrêté (par la police). »
\item \all{Neue Regeln wurden (von den Schülern) vereinbart.} « De nouvelles règles ont été convenues. »
\item \all{Die Eltern wurden informiert.} « Les parents ont été informés. »
\item \all{Dem Opfer wird (vom Psychologen) geholfen.} « On aide la victime. » — \all{helfen} exige le datif : aucun nominatif n'apparaît, le verbe reste au singulier.
\end{enumerate}
\end{corrige}

\begin{aretenir}[Bilan du passif]
\begin{itemize}
\item Forme : \all{werden} conjugué + participe II en fin de phrase ; agent = \all{von} + datif.
\item Présent : \all{wird geschützt} ; prétérit : \all{wurde geschützt} ; passé composé : \all{ist geschützt worden}.
\item Verbes à datif : pas de sujet nominatif (\all{Ihm wurde geholfen.}) ; passif impersonnel : \all{Es wurde gelacht.}
\item Action = \all{werden} ; état = \all{sein}. Au Bac, privilégiez le \all{Vorgangspassiv}.
\end{itemize}
\end{aretenir}

\coursec{Grammaire 2 : concession et opposition}

Dans la section précédente, nous avons étudié le passif (\all{Vorgangspassiv}), qui permet de mettre l'accent sur l'action et non sur son auteur : \all{Die Nachrichten wurden gesendet} (« les messages ont été envoyés »). Mais un texte sur le harcèlement numérique ne se contente pas de décrire des faits : il oppose des idées, il exprime des contrastes. Comment dire en allemand « bien que », « pourtant », « mais », « malgré », « alors que » ? C'est l'objet de cette section : les \cle{connecteurs de concession et d'opposition}. Attention : chaque connecteur allemand impose une \cle{construction différente} (place du verbe, cas). C'est l'un des points les plus piégeux pour les francophones.

\courssub{Observation : découvrir la concession dans le texte}

Reprenons une phrase du texte support :

\begin{exemplebox}[Phrase d'observation]
\all{Obwohl er die Nachrichten löschte, wurden sie immer wieder gesendet.}

« Bien qu'il ait effacé les messages, ils ont été renvoyés sans cesse. »
\end{exemplebox}

Observons la structure de cette phrase :

\begin{itemize}
\item La première partie, \all{Obwohl er die Nachrichten löschte}, exprime une \cle{concession} : on admet un fait (il a effacé les messages), mais ce fait n'empêche pas le résultat contraire (ils ont été renvoyés).
\item Après \all{obwohl}, le verbe conjugué \all{löschte} se trouve \textbf{à la fin} de la subordonnée.
\item La subordonnée occupe la première place de la phrase ; la principale commence donc par le verbe \all{wurden} (inversion sujet-verbe : \all{wurden sie}).
\end{itemize}

En français, on pourrait dire : « Il a effacé les messages, \emph{pourtant} ils ont été renvoyés » ou « \emph{Bien qu'}il ait effacé les messages, ils ont été renvoyés ». L'allemand, lui, choisit \textbf{un seul} connecteur et adapte toute la construction de la phrase. Étudions maintenant chaque outil en détail.

\courssub{obwohl : « bien que » (subordonnée, verbe en fin)}

\begin{propriete}[La subordonnée de concession avec \all{obwohl}]
\all{Obwohl} (« bien que », « quoique ») introduit une \cle{subordonnée circonstancielle de concession}. Règle absolue : le verbe conjugué est rejeté \textbf{à la dernière place} de la subordonnée.

\all{Obwohl er Angst hatte, ging er zur Polizei.}

« Bien qu'il ait eu peur, il est allé à la police. »

Deux constructions possibles :
\begin{enumerate}
\item Subordonnée d'abord : \all{Obwohl er Angst hatte, ging er zur Polizei.} (la principale suit avec inversion : verbe puis sujet).
\item Principale d'abord : \all{Er ging zur Polizei, obwohl er Angst hatte.} (ordre normal dans la principale).
\end{enumerate}
\end{propriete}

\begin{exemplebox}[Exemples avec \all{obwohl}]
\begin{itemize}
\item \all{Obwohl das Passwort geheim war, wurde das Konto gehackt.} « Bien que le mot de passe fût secret, le compte a été piraté. »
\item \all{Obwohl das Opfer die Nachrichten blockierte, schrieb der Täter weiter.} « Bien que la victime ait bloqué les messages, l'auteur a continué à écrire. »
\item \all{Sie erstattete Anzeige, obwohl sie sich schämte.} « Elle a porté plainte, bien qu'elle ait eu honte. »
\item \all{Obwohl viele Schüler das Mobbing sahen, sagte niemand etwas.} « Bien que beaucoup d'élèves aient vu le harcèlement, personne n'a rien dit. »
\end{itemize}
\end{exemplebox}

\begin{attention}[Pas de subjonctif en allemand !]
En français, « bien que » exige le subjonctif (« bien qu'il \emph{ait} eu peur »). En allemand, \all{obwohl} est suivi de l'\textbf{indicatif} : \all{obwohl er Angst \emph{hatte}}. Ne cherchez jamais à calquer le subjonctif français.
\end{attention}

\courssub{trotzdem : « pourtant » (adverbe, inversion)}

\begin{propriete}[L'adverbe de concession \all{trotzdem}]
\all{Trotzdem} (« pourtant », « néanmoins ») est un \cle{adverbe}, non une conjonction. Il relie deux phrases (ou deux propositions) dont la seconde contredit ce qu'on attendait. Placé en \textbf{première position}, il provoque l'\textbf{inversion} : le verbe conjugué vient en deuxième place, avant le sujet (règle V2).

\all{Er hatte Angst. Trotzdem ging er zur Polizei.}

« Il avait peur. Pourtant, il est allé à la police. »

On peut aussi le placer après le verbe : \all{Er ging trotzdem zur Polizei.}
\end{propriete}

\begin{exemplebox}[Exemples avec \all{trotzdem}]
\begin{itemize}
\item \all{Das Opfer hatte große Angst. Trotzdem zeigte es den Täter an.} « La victime avait très peur. Pourtant, elle a dénoncé l'auteur. »
\item \all{Die Schule warnte die Schüler. Trotzdem ging das Mobbing weiter.} « L'école a averti les élèves. Pourtant, le harcèlement a continué. »
\item \all{Er löschte sein Konto, trotzdem fanden ihn die Täter wieder.} « Il a supprimé son compte ; pourtant, les harceleurs l'ont retrouvé. »
\end{itemize}
\end{exemplebox}

\courssub{aber : « mais » (coordination, ordre normal)}

\begin{propriete}[La conjonction de coordination \all{aber}]
\all{Aber} (« mais ») est une \cle{conjonction de coordination} : elle relie deux propositions de même niveau et \textbf{ne change pas l'ordre des mots}. La proposition qui suit garde l'ordre normal : sujet puis verbe (V2 classique).

\all{Er hatte Angst, aber er ging zur Polizei.}

« Il avait peur, mais il est allé à la police. »

Notez la virgule devant \all{aber} quand la seconde proposition a son propre sujet.
\end{propriete}

\begin{exemplebox}[Exemples avec \all{aber}]
\begin{itemize}
\item \all{Das Gesetz existiert, aber viele Täter werden nicht bestraft.} « La loi existe, mais beaucoup d'auteurs ne sont pas punis. »
\item \all{Sie wollte schweigen, aber ihre Freundin überzeugte sie.} « Elle voulait se taire, mais son amie l'a convaincue. »
\item \all{Der Täter entschuldigte sich, aber das Opfer verzieh ihm nicht.} « L'auteur s'est excusé, mais la victime ne lui a pas pardonné. »
\end{itemize}
\end{exemplebox}

\courssub{trotz + génitif : « malgré » (préposition)}

\begin{propriete}[La préposition \all{trotz} + génitif]
\all{Trotz} (« malgré ») est une \cle{préposition} : elle introduit un \textbf{groupe nominal}, pas une proposition avec verbe conjugué. À l'écrit et dans la langue soignée, elle exige le \textbf{génitif} :

\all{Trotz der Beweise wagte er keine Anzeige.}

« Malgré les preuves, il n'a pas osé porter plainte. »

Formes du génitif après \all{trotz} : \all{trotz des Beweises} (masculin/neutre singulier), \all{trotz der Beweise} (pluriel ou féminin), \all{trotz eines Verbrechens}.

\textbf{Exception tolérée} : à l'oral et dans la langue familière (surtout dans le sud de l'Allemagne, en Autriche et en Suisse), on entend \all{trotz} + datif : \all{trotz dem Regen}. À l'écrit et à l'examen, exigez toujours le génitif.
\end{propriete}

\begin{exemplebox}[Exemples avec \all{trotz}]
\begin{itemize}
\item \all{Trotz des Gesetzes gibt es noch Cybermobbing.} « Malgré la loi, le harcèlement numérique existe encore. »
\item \all{Trotz der Angst sprach das Opfer mit einem Lehrer.} « Malgré la peur, la victime a parlé avec un professeur. »
\item \all{Trotz des Passworts konnte der Täter das Konto öffnen.} « Malgré le mot de passe, l'auteur a pu ouvrir le compte. »
\end{itemize}
\end{exemplebox}

\begin{attention}[Génitif obligatoire à l'écrit]
Ne dites pas \all{*trotz den Beweisen} à l'écrit : c'est un datif pluriel, fautif en langue soignée. Écrivez \all{trotz der Beweise}. Astuce : si le nom n'a pas d'article et ne porte pas de marque de génitif, le datif est admis (\all{trotz Regen}), mais avec article, le génitif s'impose.
\end{attention}

\courssub{während : « alors que » (opposition)}

\begin{propriete}[\all{während} exprimant l'opposition]
Outre son sens temporel (« pendant que »), \all{während} peut exprimer une \cle{opposition} : « alors que », « tandis que ». C'est une conjonction de subordination : le verbe conjugué va \textbf{en fin de subordonnée}, comme avec \all{obwohl}.

\all{Während das Opfer litt, lachte der Täter.}

« Alors que la victime souffrait, l'auteur riait. »
\end{propriete}

\begin{exemplebox}[Exemples avec \all{während} (opposition)]
\begin{itemize}
\item \all{Während die einen schweigen, helfen die anderen dem Opfer.} « Alors que les uns se taisent, les autres aident la victime. »
\item \all{Während der Täter anonym blieb, musste das Opfer alles öffentlich ertragen.} « Alors que l'auteur restait anonyme, la victime devait tout subir publiquement. »
\end{itemize}
\end{exemplebox}

\courssub{Tableau récapitulatif et schéma}

\begin{center}
\begin{tabularx}{\linewidth}{|l|l|X|X|}
\hline
\rowcolor{popblueL} \textbf{Connecteur} & \textbf{Sens} & \textbf{Construction} & \textbf{Exemple traduit} \\
\hline
\all{obwohl} & bien que & subordonnée, verbe en fin & \all{Obwohl er Angst hatte, ging er zur Polizei.} « Bien qu'il ait eu peur, il est allé à la police. » \\
\hline
\all{trotzdem} & pourtant & adverbe en 1re place, inversion (V2) & \all{Er hatte Angst. Trotzdem ging er zur Polizei.} « Il avait peur. Pourtant, il est allé à la police. » \\
\hline
\all{aber} & mais & coordination, ordre normal & \all{Er hatte Angst, aber er ging zur Polizei.} « Il avait peur, mais il est allé à la police. » \\
\hline
\all{trotz} + gén. & malgré & préposition + groupe nominal (génitif) & \all{Trotz der Beweise wagte er keine Anzeige.} « Malgré les preuves, il n'a pas osé porter plainte. » \\
\hline
\all{während} & alors que & subordonnée, verbe en fin & \all{Während das Opfer litt, lachte der Täter.} « Alors que la victime souffrait, l'auteur riait. » \\
\hline
\end{tabularx}
\end{center}

\begin{popfigure}
\begin{tikzpicture}[
  conn/.style={rectangle, rounded corners=3pt, draw=popblue, fill=popblueL, align=center, font=\small, minimum height=0.9cm},
  sens/.style={rectangle, rounded corners=3pt, draw=poporange, fill=poporangeL, align=center, font=\small, minimum height=0.9cm},
  verb/.style={rectangle, rounded corners=3pt, draw=popgreen, fill=popgreenL, align=center, font=\small, minimum height=0.9cm},
  head/.style={font=\bfseries\small, text=popdark}
]
\node[head] at (0,0) {Connecteur};
\node[head] at (4.2,0) {Sens};
\node[head] at (9.2,0) {Place du verbe};
\node[conn] at (0,-1.2) {\all{obwohl}};
\node[sens] at (4.2,-1.2) {bien que};
\node[verb, align=center] at (9.2,-1.2){verbe conjugué\\ en fin de subordonnée};
\node[conn] at (0,-2.6) {\all{trotzdem}};
\node[sens] at (4.2,-2.6) {pourtant};
\node[verb, align=center] at (9.2,-2.6){1re place :\\ inversion sujet--verbe};
\node[conn] at (0,-4.0) {\all{aber}};
\node[sens] at (4.2,-4.0) {mais};
\node[verb, align=center] at (9.2,-4.0){ordre normal\\ (sujet -- verbe)};
\node[conn] at (0,-5.4) {\all{trotz}};
\node[sens] at (4.2,-5.4) {malgré};
\node[verb, align=center] at (9.2,-5.4){pas de verbe :\\ nom au génitif};
\end{tikzpicture}
\legende{Les quatre connecteurs de concession : sens et place du verbe}
\end{popfigure}

\courssub{Comparaison français--allemand}

\begin{center}
\begin{tabularx}{\linewidth}{|X|X|X|}
\hline
\rowcolor{popblueL} \textbf{Français} & \textbf{Deutsch} & \textbf{Remarque} \\
\hline
Bien qu'il ait peur, il parle. & \all{Obwohl er Angst hat, spricht er.} & indicatif en allemand, pas de subjonctif \\
\hline
Il a peur, pourtant il parle. & \all{Er hat Angst. Trotzdem spricht er.} & inversion après \all{trotzdem} \\
\hline
Il a peur, mais il parle. & \all{Er hat Angst, aber er spricht.} & ordre inchangé après \all{aber} \\
\hline
Malgré sa peur, il parle. & \all{Trotz seiner Angst spricht er.} & génitif : \all{seiner Angst} \\
\hline
\end{tabularx}
\end{center}

\begin{attention}[Un seul connecteur par idée !]
En français on entend parfois « bien qu'il ait peur, \emph{pourtant} il est allé à la police ». En allemand, cette double marque est une faute de style (pléonasme) : ne combinez \textbf{jamais} \all{obwohl} et \all{trotzdem} dans la même phrase. Choisissez :

\all{Obwohl er Angst hatte, ging er zur Polizei.} \emph{ou} \all{Er hatte Angst. Trotzdem ging er zur Polizei.}

Jamais : \all{*Obwohl er Angst hatte, trotzdem ging er zur Polizei.}
\end{attention}

\courssub{Texte d'application : le courage de Lea}

\begin{textebox}[Mut gegen Cybermobbing]
Lea ist 17 Jahre alt und besucht ein Gymnasium in Berlin. Obwohl sie eine gute Schülerin war, verlor sie plötzlich ihre Freude an der Schule. Der Grund: Ein Mitschüler hatte ein peinliches Foto von ihr ins Internet gestellt. Obwohl Lea ihn bat, das Foto zu löschen, weigerte er sich. Trotz ihrer Angst vor den Reaktionen der Klasse sprach sie zuerst mit ihrer besten Freundin. Diese riet ihr: „Du musst mit jemandem sprechen. Trotzdem hast du Angst? Das verstehe ich, aber schweigen hilft nicht.“ Lea hatte große Angst, aber sie ging zur Vertrauenslehrerin. Während Lea unter den Nachrichten litt, machte der Täter immer neue Kommentare. Trotz des Drucks blieb Lea stark. Obwohl der Täter das Foto endlich löschte, war es schon hundertmal geteilt worden. Die Schule reagierte sofort: Der Täter wurde bestraft, und die Klasse diskutierte eine Woche lang über Cybermobbing. Trotz der schwierigen Zeit bereut Lea ihre Entscheidung nicht. „Ich habe gelernt“, sagt sie, „dass man sich wehren muss. Obwohl es schwer war, bin ich heute stärker als vorher.“ Ihre Geschichte zeigt: Trotz aller Gefahren des Internets kann man sich schützen — aber nur, wenn man nicht schweigt.
\end{textebox}

\textbf{Glossaire} : der Mitschüler, -- (le camarade de classe) ; peinlich (embarrassant) ; löschen (effacer) ; sich weigern (refuser) ; die Vertrauenslehrerin, --nen (la conseillère d'élèves) ; der Druck, --e (la pression) ; teilen (partager) ; bereuen (regretter) ; sich wehren (se défendre) ; schweigen (se taire).

\textbf{Questions de compréhension} :
\begin{enumerate}
\item Warum verlor Lea ihre Freude an der Schule ?
\item Mit wem sprach Lea zuerst, und was riet ihr diese Person ?
\item Wie reagierte die Schule auf den Fall ?
\end{enumerate}

\courssub{Exercice résolu et entraînement}

\begin{exoresolu}[Choisir le bon connecteur]
Complétez avec \all{obwohl}, \all{trotzdem}, \all{aber} ou \all{trotz} (+ génitif), puis traduisez :
\begin{enumerate}
\item \all{\ldots\ das Opfer Angst hatte, zeigte es den Täter an.}
\item \all{Der Täter wurde bestraft, \ldots\ das Mobbing hörte nicht auf.}
\item \all{\ldots\ der Beweise leugnete der Täter alles.}
\item \all{Lea hatte Angst, \ldots\ sie ging zur Vertrauenslehrerin.}
\end{enumerate}
\tcblower
\textbf{Solution détaillée} :
\begin{enumerate}
\item \all{Obwohl das Opfer Angst hatte, zeigte es den Täter an.} — Verbe conjugué \all{hatte} en fin de subordonnée : il faut \all{obwohl}. « Bien que la victime ait eu peur, elle a dénoncé l'auteur. »
\item \all{Der Täter wurde bestraft, trotzdem hörte das Mobbing nicht auf.} — Après \all{trotzdem}, inversion : \all{hörte} avant le sujet \all{das Mobbing}. « L'auteur a été puni ; pourtant, le harcèlement n'a pas cessé. »
\item \all{Trotz der Beweise leugnete der Täter alles.} — Groupe nominal au génitif pluriel : \all{der Beweise}. « Malgré les preuves, l'auteur a tout nié. »
\item \all{Lea hatte Angst, aber sie ging zur Vertrauenslehrerin.} — Coordination, ordre normal après \all{aber}. « Lea avait peur, mais elle est allée voir la conseillère. »
\end{enumerate}
\end{exoresolu}

\begin{exercice}[À vous de jouer]
\begin{enumerate}
\item Traduisez en allemand : a) « Bien que le mot de passe soit secret, le compte a été piraté. » b) « Il avait peur ; pourtant, il a porté plainte. » c) « Malgré la loi, le harcèlement numérique existe encore. » d) « Alors que la victime souffrait, l'auteur riait. »
\item Transformez : \all{Obwohl er sich schämte, sprach er mit der Polizei.} $\rightarrow$ phrase avec \all{trotzdem}, puis avec \all{aber}, puis avec \all{trotz} + génitif.
\item Corrigez la faute : \all{*Obwohl sie Angst hatte, trotzdem ging sie zur Schule.}
\end{enumerate}
\end{exercice}

\begin{corrige}[Exercice « À vous de jouer »]
\begin{enumerate}
\item a) \all{Obwohl das Passwort geheim ist, wurde das Konto gehackt.} b) \all{Er hatte Angst. Trotzdem erstattete er Anzeige.} c) \all{Trotz des Gesetzes gibt es noch Cybermobbing.} d) \all{Während das Opfer litt, lachte der Täter.}
\item \all{Er schämte sich. Trotzdem sprach er mit der Polizei.} / \all{Er schämte sich, aber er sprach mit der Polizei.} / \all{Trotz seiner Scham sprach er mit der Polizei.}
\item Il faut un seul connecteur : \all{Obwohl sie Angst hatte, ging sie zur Schule.} ou \all{Sie hatte Angst. Trotzdem ging sie zur Schule.}
\end{enumerate}
\end{corrige}

\begin{aretenir}[L'essentiel de la section]
\begin{itemize}
\item \all{obwohl} (bien que) : subordonnée, verbe conjugué \textbf{en fin} ; indicatif, jamais de subjonctif.
\item \all{trotzdem} (pourtant) : adverbe en première place, \textbf{inversion} sujet--verbe.
\item \all{aber} (mais) : coordination, \textbf{ordre normal} inchangé.
\item \all{trotz} + \textbf{génitif} (malgré) : groupe nominal sans verbe ; datif seulement toléré à l'oral.
\item \all{während} (alors que) : opposition, verbe en fin de subordonnée.
\item Un seul connecteur par idée : jamais \all{obwohl} + \all{trotzdem} ensemble.
\end{itemize}
\end{aretenir}

\coursec{Méthode du commentaire et commentaire modèle}

Vous savez désormais analyser la concession et l'opposition (\all{obwohl}, \all{trotzdem}, \all{während}) et rapporter des faits au passif. Ces deux outils grammaticaux ne sont pas de simples exercices : ce sont précisément les instruments dont vous aurez besoin pour rédiger un \cle{commentaire de texte} réussi au baccalauréat. Un commentaire, c'est rapporter des faits (passif) et nuancer son jugement (concession). Cette section vous donne la méthode complète, des phrases-outils prêtes à l'emploi et un modèle rédigé à imiter.

\courssub{Qu'est-ce qu'un commentaire de texte au Bac ?}

Au baccalauréat, le commentaire (\all{der Kommentar}) est une production écrite courte (8 à 12 lignes) dans laquelle vous montrez que vous avez compris le texte et que vous savez prendre position. Ce n'est ni une traduction, ni une simple paraphrase : c'est un \cle{texte argumentatif miniature}.

\begin{methode}[Les 5 étapes du commentaire de texte]
\begin{enumerate}
\item \cle{Présenter} le texte : source, type de texte (article de presse, récit, lettre...), thème général. Une seule phrase suffit.
\item \cle{Résumer} le contenu en 3 ou 4 phrases, avec vos propres mots, en rapportant les faits au passif.
\item \cle{Dégager la structure} : comment le texte est-il organisé ? (\all{Zuerst... dann... schließlich...})
\item \cle{Analyser la thèse de l'auteur} : que veut-il démontrer ou dénoncer ? Quel est son message ?
\item \cle{Donner son avis argumenté} : \all{Meiner Meinung nach...} suivi d'un argument (\all{weil}, \all{denn}) et d'une nuance (\all{obwohl}, \all{trotzdem}).
\end{enumerate}
\end{methode}

\begin{popfigure}
\begin{tikzpicture}[
box/.style={draw=popblue, fill=popblueL, rounded corners=3pt, text width=7.5cm, align=center, minimum height=0.9cm, font=\small},
arr/.style={-{Stealth[length=2.5mm]}, thick, popdark}]
\node[box] (a) {\textbf{1. Présentation} : source, type de texte, thème};
\node[box, below=0.35cm of a] (b) {\textbf{2. Résumé} : 3--4 phrases, au passif};
\node[box, below=0.35cm of b] (c) {\textbf{3. Structure} : zuerst -- dann -- schließlich};
\node[box, below=0.35cm of c] (d) {\textbf{4. Analyse} : thèse et intention de l'auteur};
\node[box, below=0.35cm of d, fill=poporangeL, draw=poporange] (e) {\textbf{5. Avis personnel} : Meiner Meinung nach... + argument + nuance};
\draw[arr] (a) -- (b);
\draw[arr] (b) -- (c);
\draw[arr] (c) -- (d);
\draw[arr] (d) -- (e);
\end{tikzpicture}
\legende{Organigramme des 5 étapes du commentaire de texte}
\end{popfigure}

\courssub{Les phrases-outils du commentaire}

Un bon candidat ne cherche pas ses mots : il utilise des \cle{phrases-outils} (\all{Satzbausteine}) mémorisées à l'avance. Voici le répertoire indispensable, classé par étape.

\begin{center}
\begin{tabularx}{\linewidth}{|l|X|X|}
\hline
\rowcolor{popblueL} \textbf{Étape} & \textbf{Deutsch} & \textbf{Français} \\
\hline
Présentation & \all{Der Text handelt von...} & Le texte traite de... \\
\hline
Présentation & \all{Der Artikel erschien in...} & L'article a paru dans... \\
\hline
Résumé & \all{Der Autor berichtet, dass...} & L'auteur rapporte que... \\
\hline
Résumé & \all{Es wird beschrieben, wie...} & Il est décrit comment... \\
\hline
Structure & \all{Zuerst..., dann..., schließlich...} & D'abord..., ensuite..., enfin... \\
\hline
Analyse & \all{Der Autor will zeigen, dass...} & L'auteur veut montrer que... \\
\hline
Analyse & \all{Das Ziel des Textes ist...} & Le but du texte est... \\
\hline
Avis & \all{Meiner Meinung nach...} & À mon avis... \\
\hline
Avis & \all{Ich bin der Meinung, dass...} & Je suis d'avis que... \\
\hline
Argument & \all{..., weil / denn...} & ..., parce que / car... \\
\hline
Nuance & \all{Obwohl..., ... / Trotzdem...} & Bien que..., ... / Pourtant... \\
\hline
\end{tabularx}
\end{center}

\begin{attention}[La place du verbe dans les phrases-outils]
Deux pièges de construction attendent les francophones :
\begin{itemize}
\item Après \all{Meiner Meinung nach}, le verbe conjugué reste en \cle{deuxième position} : \all{Meiner Meinung nach \textbf{müssen} Opfer besser geschützt \textbf{werden}.} (et non « Meiner Meinung nach Opfer müssen... »).
\item Après \all{Ich bin der Meinung, dass...}, la subordonnée envoie le verbe à la \cle{fin} : \all{Ich bin der Meinung, dass die Schulen mehr präventiv arbeiten \textbf{sollten}.}
\end{itemize}
\end{attention}

\begin{exemplebox}[Exemples traduits]
\all{Meiner Meinung nach müssen Opfer geschützt werden.} « À mon avis, les victimes doivent être protégées. »

\all{Ich bin der Meinung, dass das Passwort geheim bleiben muss.} « Je suis d'avis que le mot de passe doit rester secret. »

\all{Der Autor berichtet, dass ein Schüler im Internet gemobbt wurde.} « L'auteur rapporte qu'un élève a été harcelé sur Internet. »

\all{Obwohl der Täter bestraft wurde, bleibt das Problem aktuell.} « Bien que l'auteur (des faits) ait été puni, le problème reste d'actualité. »
\end{exemplebox}

\courssub{Un commentaire modèle sur le texte de Berlin}

Lisons maintenant un commentaire complet, rédigé selon les 5 étapes, qui réemploie le passif et les connecteurs de concession étudiés dans cette leçon.

\begin{textebox}[Modellkommentar — commentaire modèle]
Der Text „Cybermobbing -- ein Fall aus Berlin" berichtet von einem Schüler, der im Internet gemobbt wurde. Zuerst wurden Fotos verbreitet, dann wurde das Opfer bedroht. Der Autor will zeigen, dass Cybermobbing ein ernstes Problem an den Schulen ist. Obwohl der Täter bestraft wurde, bleibt das Problem aktuell. Meiner Meinung nach müssen Opfer besser geschützt werden, denn Cybermobbing kann Leben zerstören. Trotzdem sollte man nicht nur bestrafen, sondern auch präventiv in den Schulen arbeiten.

\emph{Glossaire :} der Schüler, die Schüler (l'élève) ; mobben (harceler) ; verbreiten (diffuser, répandre) ; bedrohen (menacer) ; bestrafen (punir) ; zerstören (détruire) ; präventiv arbeiten (travailler en prévention).
\end{textebox}

\textbf{Lecture commentée du modèle.} Analysons phrase par phrase comment ce modèle applique la méthode :

\begin{itemize}
\item \all{Der Text ... berichtet von einem Schüler, der im Internet gemobbt wurde.} → \cle{Étape 1} (présentation : titre, thème) avec une première phrase-outil et un \cle{passif prétérit} : \all{wurde gemobbt}.
\item \all{Zuerst wurden Fotos verbreitet, dann wurde das Opfer bedroht.} → \cle{Étapes 2 et 3} (résumé + structure) : les connecteurs chronologiques \all{zuerst / dann} organisent les faits, rapportés au \cle{passif} (\all{wurden verbreitet}, \all{wurde bedroht}) — on rapporte ce qui a été subi, sans recopier le texte.
\item \all{Der Autor will zeigen, dass...} → \cle{Étape 4} (analyse de l'intention de l'auteur).
\item \all{Obwohl der Täter bestraft wurde, bleibt das Problem aktuell.} → nuance grammaticale de la leçon : \cle{concession} avec \all{obwohl} (verbe en fin de subordonnée) + passif prétérit.
\item \all{Meiner Meinung nach müssen Opfer besser geschützt werden, denn...} → \cle{Étape 5} (avis) : phrase-outil d'opinion + argument introduit par \all{denn} + \cle{passif présent} avec modal (\all{müssen ... geschützt werden}).
\item \all{Trotzdem sollte man nicht nur bestrafen, sondern auch präventiv arbeiten.} → ouverture nuancée avec \all{trotzdem} et la structure \all{nicht nur..., sondern auch...}.
\end{itemize}

\begin{aretenir}[La règle d'or du commentaire]
Un bon commentaire = \cle{30 \% de résumé} et \cle{70 \% d'analyse et d'avis personnel}. L'examinateur veut lire \emph{votre} pensée, pas une paraphrase du texte. Justifiez toujours votre avis par un argument (\all{weil}, \all{denn}) et apportez une nuance (\all{obwohl}, \all{trotzdem}) : c'est ce qui distingue une copie excellente d'une copie moyenne.
\end{aretenir}

\begin{attention}[Les erreurs fatales au résumé]
\begin{itemize}
\item \cle{Ne recopiez jamais} des phrases entières du texte : reformulez avec vos propres mots. Le correcteur repère immédiatement le copier-coller.
\item Employez le \cle{passif} pour rapporter les faits : \all{Fotos wurden verbreitet} est plus objectif et plus élégant que \all{Der Täter verbreitete Fotos} quand l'agent importe peu.
\item Attention à la concession : après \all{obwohl}, le verbe va à la fin (\all{Obwohl der Täter bestraft \textbf{wurde}, ...}) ; \all{trotzdem} occupe la première position et fait basculer le sujet après le verbe (\all{\textbf{Trotzdem bleibt} das Problem aktuell}).
\item Faux ami : \all{aktualisieren} signifie « mettre à jour » ; pour « d'actualité », dites \all{aktuell} ou \all{immer noch ein Problem}.
\end{itemize}
\end{attention}

\courssub{Et au Cameroun ? Le harcèlement numérique dans nos lycées}

Le commentaire gagne toujours à s'ouvrir sur le vécu du candidat. Le cyberharcèlement n'est pas un problème réservé à Berlin : dans les lycées de Yaoundé, de Douala ou de Bamenda, des photos humiliantes circulent sur WhatsApp, des groupes de classe servent parfois à exclure un camarade, des rumeurs se répandent en quelques minutes. Relier le texte à cette réalité camerounaise montre à l'examinateur que vous savez \cle{transférer} vos connaissances.

\begin{exemplebox}[Phrases pour parler du contexte camerounais]
\all{Auch in Kamerun werden Schüler oft über WhatsApp gemobbt.} « Au Cameroun aussi, des élèves sont souvent harcelés via WhatsApp. »

\all{In unseren Gymnasien wird dieses Problem selten diskutiert.} « Dans nos lycées, ce problème est rarement débattu. »

\all{Obwohl viele Jugendliche ein Handy besitzen, werden sie nicht über die Gefahren informiert.} « Bien que beaucoup de jeunes possèdent un téléphone portable, ils ne sont pas informés des dangers. »

\all{Ich bin der Meinung, dass auch unsere Schulen Präventionsarbeit leisten müssen.} « Je suis d'avis que nos écoles aussi doivent faire un travail de prévention. »

\emph{Glossaire :} das Gymnasium, die Gymnasien (le lycée) ; das Handy, die Handys (le portable) ; die Gefahr, die Gefahren (le danger) ; die Präventionsarbeit (le travail de prévention).
\end{exemplebox}

\begin{exoresolu}[Rédiger un mini-commentaire]
\textbf{Énoncé.} À partir du texte « Cybermobbing -- ein Fall aus Berlin », rédigez les deux premières phrases d'un commentaire (présentation + début du résumé), en utilisant \all{handeln von} et un passif prétérit.

\tcblower
\textbf{Solution détaillée.} Étape 1, présentation : \all{Der Text „Cybermobbing -- ein Fall aus Berlin" handelt von einem Schüler, der zum Opfer von Cybermobbing wurde.} (Le texte traite d'un élève devenu victime de cyberharcèlement — passif prétérit : \all{wurde}.) Étape 2, résumé : \all{Zuerst wurden private Fotos im Internet verbreitet, dann wurde das Opfer bedroht.} (D'abord des photos privées ont été diffusées sur Internet, ensuite la victime a été menacée.) On remarque : le verbe en deuxième position, le participe II en fin de phrase, et aucune phrase recopiée du texte.
\end{exoresolu}

\courssub{Production guidée et entraînement}

\begin{experience}[À vous de commenter !]
En groupes de quatre, rédigez un commentaire de 8 à 10 lignes sur le thème : \all{Cybermobbing an kamerunischen Schulen}. Chaque membre du groupe rédige une étape (présentation, résumé, structure, analyse, avis), puis le groupe assemble et relit le texte complet. Contraintes obligatoires : au moins \cle{deux passifs} (présent et prétérit), \cle{une concession} avec \all{obwohl}, \cle{une nuance} avec \all{trotzdem} et \cle{un avis} avec \all{Meiner Meinung nach}. Chaque groupe lit ensuite son commentaire à voix haute ; la classe repère les phrases-outils utilisées.
\end{experience}

\begin{exercice}[Entraînement à la maison]
\begin{enumerate}
\item Remettez dans l'ordre les cinq étapes d'un commentaire : avis personnel / résumé / présentation / analyse / structure.
\item Traduisez en allemand : « À mon avis, les victimes doivent être mieux protégées, car le cyberharcèlement peut détruire des vies. »
\item Complétez avec \all{obwohl} ou \all{trotzdem} : \all{... der Täter bestraft wurde, bleibt das Opfer traumatisiert.} / \all{Der Täter wurde bestraft. ... bleibt das Opfer traumatisiert.}
\item Rédigez la phrase d'analyse (étape 4) d'un commentaire sur un article traitant des réseaux sociaux au Cameroun, avec \all{Der Autor will zeigen, dass...}.
\end{enumerate}
\end{exercice}

\begin{corrige}[Exercice : Entraînement à la maison]
\begin{enumerate}
\item Ordre correct : présentation → résumé → structure → analyse → avis personnel.
\item \all{Meiner Meinung nach müssen die Opfer besser geschützt werden, denn Cybermobbing kann Leben zerstören.}
\item \all{\textbf{Obwohl} der Täter bestraft wurde, bleibt das Opfer traumatisiert.} (subordonnée, verbe à la fin) ; \all{Der Täter wurde bestraft. \textbf{Trotzdem} bleibt das Opfer traumatisiert.} (adverbe en position 1, inversion du sujet).
\item Exemple : \all{Der Autor will zeigen, dass soziale Netzwerke in Kamerun sowohl Chancen als auch Gefahren für Jugendliche darstellen.} (L'auteur veut montrer que les réseaux sociaux représentent au Cameroun à la fois des chances et des dangers pour les jeunes.)
\end{enumerate}
\end{corrige}

\begin{aretenir}[Ce qu'il faut retenir]
Le commentaire de texte suit un plan en cinq étapes : présenter, résumer, structurer, analyser, donner son avis. Les phrases-outils (\all{Der Text handelt von...}, \all{Meiner Meinung nach...}) sécurisent votre expression. Le passif rapporte les faits avec objectivité ; les connecteurs de concession (\all{obwohl}, \all{trotzdem}) apportent la nuance qui fait la différence au baccalauréat. Retenez la proportion d'or : 30 \% de résumé, 70 \% d'analyse et d'avis argumenté.
\end{aretenir}

\coursec{Production guidée et entraînement}

Après avoir étudié la méthode du commentaire de texte et analysé un commentaire modèle sur le cas de Cybermobbing à Berlin, il est temps de passer à la pratique. Cette dernière étape de la leçon vous fait produire vous-mêmes : d'abord à l'écrit, avec une courte rédaction guidée et fortement contrainte, puis à l'oral, avec un débat en binômes. Des exercices de grammaire et de conjugaison consolident enfin le passif, avant une correction collective fondée sur une grille d'auto-évaluation. C'est exactement le type de tâche attendu au baccalauréat A : un contenu pertinent, une langue correcte, des connecteurs maîtrisés.

\courssub{Production écrite guidée : Ratschläge gegen Cybermobbing}

\textbf{La tâche.} Vous rédigez un court texte de 6 à 8 phrases intitulé \all{Ratschläge gegen Cybermobbing} (Des conseils contre le harcèlement numérique). Le plan est imposé :

\begin{enumerate}
\item \cle{Problème} (1 à 2 phrases) : présenter brièvement le phénomène du Cybermobbing.
\item \cle{Conseils} (3 à 4 phrases) : donner des recommandations concrètes aux victimes et aux témoins.
\item \cle{Conclusion} (1 à 2 phrases) : ouvrir sur l'importance de la solidarité ou de la prévention.
\end{enumerate}

\textbf{Les contraintes linguistiques obligatoires :}

\begin{itemize}
\item au moins \cle{3 verbes au passif} (présent ou prétérit) ;
\item au moins \cle{1 phrase avec \all{obwohl}} (bien que, subordonnée, verbe à la fin) ;
\item au moins \cle{1 phrase avec \all{trotz} + génitif} (malgré) ;
\item au moins \cle{5 mots du lexique de la leçon} : \all{das Cybermobbing}, \all{das Opfer}, \all{der Täter}, \all{die Anzeige}, \all{schützen}, \all{das Passwort}, \all{der Beweis}, \all{melden}, \all{speichern}\ldots
\end{itemize}

\begin{methode}[Rédiger en trois temps : plan, écriture, relecture]
\begin{enumerate}
\item \textbf{Rédiger d'abord un plan en 3 lignes.} Une ligne par partie : notez deux ou trois mots-clés allemands, pas des phrases entières. Exemple : \all{Problem: viele Jugendliche, Opfer schweigen} / \all{Ratschläge: Passwort, Beweise speichern, Anzeige} / \all{Schluss: gemeinsam handeln}.
\item \textbf{Écrire} en suivant le plan, une phrase simple à la fois. Vérifiez au fur et à mesure que chaque contrainte est remplie (cochez : 3 passifs ? \all{obwohl} ? \all{trotz} ? 5 mots du lexique ?).
\item \textbf{Relire} en vérifiant trois points critiques : la \cle{place du verbe} (2\textsuperscript{e} position dans la principale, dernière position après \all{obwohl}), les \cle{terminaisons} de l'auxiliaire \all{werden}, et le \cle{participe II à la fin} au passif (\all{wird geschützt}, \all{wurde gemeldet}).
\end{enumerate}
\end{methode}

\begin{exemplebox}[Phrases-modèles au passif pour la rédaction]
\all{Passwörter müssen geheim gehalten werden.} « Les mots de passe doivent être gardés secrets. »

\all{Beweise sollten gespeichert werden.} « Les preuves devraient être sauvegardées. »

\all{Das Opfer muss ernst genommen werden.} « La victime doit être prise au sérieux. »

\all{Der Täter wurde von der Schule gemeldet.} « Le harceleur a été signalé par l'école. »
\end{exemplebox}

\begin{attention}[Passif + verbe modal : l'ordre exact]
Avec un verbe modal, la structure est : \textbf{modal conjugué (2\textsuperscript{e} position) + participe II + \all{werden} à l'infinitif en toute fin de phrase}.

\all{Opfer müssen geschützt werden.} « Les victimes doivent être protégées. »

Ne dites jamais \all{*Opfer müssen werden geschützt} : l'infinitif \all{werden} se place toujours \emph{après} le participe II, à la toute fin. Astuce : mémorisez le schéma \textbf{modal \ldots Partizip II + werden}.
\end{attention}

\begin{textebox}[Modèle de rédaction : Ratschläge gegen Cybermobbing]
\all{Cybermobbing ist ein ernstes Problem: Viele Jugendliche werden im Internet beleidigt oder bedroht, und oft schweigen die Opfer aus Angst. Obwohl das Cybermobbing sehr verbreitet ist, wird es selten gemeldet.}

\all{Was kann man tun? Zuerst müssen die Passwörter geheim gehalten werden, damit niemand das Konto des Opfers benutzen kann. Zweitens sollten alle Beweise, zum Beispiel Screenshots, gespeichert werden. Drittens muss der Täter bei der Plattform und in der Schule gemeldet werden. Trotz der Angst vor Reaktionen sollte eine Anzeige bei der Polizei erstattet werden, wenn die Drohungen ernst sind.}

\all{Zum Schluss: Niemand darf allein gelassen werden. Wenn Opfer geschützt und unterstützt werden, verliert der Täter seine Macht. Gemeinsam kann das Cybermobbing bekämpft werden!}

\textbf{Glossaire} : \all{beleidigen} = insulter ; \all{bedrohen} = menacer ; \all{schweigen} = se taire ; \all{verbreitet} = répandu ; \all{das Konto, die Konten} = le compte ; \all{der Screenshot, die Screenshots} = la capture d'écran ; \all{die Plattform, die Plattformen} = la plateforme ; \all{eine Anzeige erstatten} = porter plainte ; \all{die Drohung, die Drohungen} = la menace ; \all{unterstützen} = soutenir ; \all{die Macht} = le pouvoir ; \all{bekämpfen} = combattre.

\textbf{Questions} : 1) Combien de formes passives comptez-vous dans ce texte ? 2) Où se trouve le verbe conjugué dans la phrase avec \all{obwohl} ? 3) Quel mot suit \all{trotz} et à quel cas est-il ?
\end{textebox}

\begin{corrige}[Questions sur le modèle]
1) Onze formes passives : \all{werden beleidigt}, \all{bedroht [werden]}, \all{wird gemeldet}, \all{müssen geheim gehalten werden}, \all{sollten gespeichert werden}, \all{muss gemeldet werden}, \all{sollte erstattet werden}, \all{darf allein gelassen werden}, \all{werden geschützt und unterstützt}, \all{kann bekämpft werden}. 2) À la fin de la subordonnée : \all{Obwohl das Cybermobbing sehr verbreitet \textbf{ist}}. 3) \all{der Angst} : génitif féminin singulier (\all{die Angst} $\rightarrow$ \all{der Angst}).
\end{corrige}

\courssub{Expression orale : débat en binômes}

\begin{experience}[Soll man den Täter anzeigen ?]
Par groupes de deux : l'un défend la plainte (\all{die Anzeige}), l'autre préfère d'autres solutions (dialogue avec le harceleur, médiation scolaire). Durée : 5 minutes de préparation, 3 minutes de débat. Utilisez les phrases-outils du commentaire de texte :

\begin{itemize}
\item Donner son avis : \all{Meiner Meinung nach \ldots} / \all{Ich bin der Ansicht, dass \ldots} (verbe à la fin !)
\item Argumenter : \all{Ein wichtiges Argument ist, dass \ldots} / \all{Erstens \ldots, zweitens \ldots}
\item Concéder : \all{Obwohl du recht hast, \ldots} / \all{Trotz deiner Argumente glaube ich, dass \ldots}
\item Conclure : \all{Zusammenfassend kann man sagen, dass \ldots}
\end{itemize}

Exemple d'échange : \all{— Ich bin der Ansicht, dass der Täter angezeigt werden muss. — Obwohl das stimmt, sollte zuerst das Opfer geschützt werden.} « — Je suis d'avis que le harceleur doit être dénoncé. — Bien que ce soit vrai, la victime devrait d'abord être protégée. »
\end{experience}

\courssub{Exercices de grammaire et de conjugaison}

\begin{exoresolu}[De l'actif au passif]
\textbf{Énoncé.} Mettez au passif, d'abord au présent puis au prétérit : \all{Die Schule schützt die Opfer.}

\tcblower
\textbf{Solution détaillée.} Étape 1 : repérer l'objet accusatif (\all{die Opfer}) qui devient le sujet nominatif du passif. Étape 2 : conjuguer \all{werden} à la personne du nouveau sujet ; le participe II (\all{geschützt}) va à la fin ; l'ancien sujet devient \all{von + datif}.

Présent : \all{Die Opfer werden (von der Schule) geschützt.} « Les victimes sont protégées (par l'école). »

Prétérit : \all{Die Opfer wurden (von der Schule) geschützt.} « Les victimes ont été / furent protégées. »
\end{exoresolu}

\begin{exercice}[Transformations au passif]
Mettez les phrases suivantes au passif, au présent puis au prétérit.

\begin{enumerate}
\item \all{Der Täter beleidigt das Opfer.}
\item \all{Die Polizei vernimmt den Täter.}
\item \all{Man speichert die Beweise.}
\item \all{Die Lehrer melden den Fall.}
\item \all{Die Plattform löscht die Nachrichten.}
\end{enumerate}
\end{exercice}

\begin{corrige}[Exercice de transformations]
1) \all{Das Opfer wird / wurde (vom Täter) beleidigt.} 2) \all{Der Täter wird / wurde (von der Polizei) vernommen.} 3) \all{Die Beweise werden / wurden gespeichert.} (avec \all{man}, l'agent disparaît au passif). 4) \all{Der Fall wird / wurde (von den Lehrern) gemeldet.} 5) \all{Die Nachrichten werden / wurden (von der Plattform) gelöscht.}

Attention aux participes II irréguliers : \all{vernehmen $\rightarrow$ vernommen}, et à \all{von + dem = vom}.
\end{corrige}

\begin{exercice}[Conjugaison de werden : mini-texte à trous]
Complétez avec la bonne forme de \all{werden} au présent (phrases 1 à 3) ou au prétérit (phrases 4 à 5).

\all{1) Viele Jugendliche \ldots\ im Internet beleidigt. 2) Das Passwort \ldots\ oft nicht geheim gehalten. 3) Die Opfer \ldots\ von ihren Freunden unterstützt. 4) Gestern \ldots\ der Täter in der Schule gemeldet. 5) Die Beweise \ldots\ von der Polizei gespeichert.}
\end{exercice}

\begin{corrige}[Exercice de conjugaison]
1) \all{werden} (pluriel, présent). 2) \all{wird} (3\textsuperscript{e} pers. sg., présent). 3) \all{werden}. 4) \all{wurde} (3\textsuperscript{e} pers. sg., prétérit). 5) \all{wurden} (pluriel, prétérit).

Rappel : présent \all{ich werde, du wirst, er/sie/es wird, wir werden, ihr werdet, sie werden} ; prétérit \all{ich wurde, du wurdest, er wurde, wir wurden, ihr wurdet, sie wurden}.
\end{corrige}

\begin{center}
\begin{tabular}{|l|l|l|}
\hline
\rowcolor{popblueL} Personne & Präsens & Präteritum \\
\hline
ich & werde & wurde \\
\hline
du & wirst & wurdest \\
\hline
er/sie/es & wird & wurde \\
\hline
wir & werden & wurden \\
\hline
ihr & werdet & wurdet \\
\hline
sie/Sie & werden & wurden \\
\hline
\end{tabular}
\end{center}

\begin{attention}[Pièges fréquents des francophones au passif]
\begin{itemize}
\item Ne confondez pas \all{werden} (auxiliaire du passif) et \all{sein} : \all{Die Tür wird geschlossen} = on ferme la porte (action, \all{Vorgangspassiv}) ; \all{Die Tür ist geschlossen} = la porte est fermée (état, \all{Zustandspassiv}).
\item Le participe II reste invariable : \all{Die Opfer werden geschützt}, jamais d'accord comme en français (« protégées »).
\item \all{Man speichert die Beweise} se transforme sans agent : \all{Die Beweise werden gespeichert} — on ne peut pas traduire « on » par \all{von man} !
\end{itemize}
\end{attention}

\courssub{Correction collective et auto-évaluation}

La correction se fait au tableau : un élève écrit une phrase de sa production, la classe la corrige ensemble (place du verbe, terminaison de \all{werden}, participe II final, connecteurs). Chaque élève évalue ensuite sa propre copie avec la grille ci-dessous, sur 10 points.

\begin{popfigure}
\begin{tikzpicture}
\node[draw=popblue, fill=popblueL, rounded corners, minimum width=4.2cm, minimum height=0.9cm, align=center, font=\bfseries] (c) at (0,0) {Contenu\\ (2 pts)};
\node[draw=popgreen, fill=popgreenL, rounded corners, minimum width=4.2cm, minimum height=0.9cm, align=center, font=\bfseries] (l) at (5,0) {Lexique\\ (2 pts)};
\node[draw=poporange, fill=poporangeL, rounded corners, minimum width=4.2cm, minimum height=0.9cm, align=center, font=\bfseries] (p) at (10,0) {Passif\\ (2 pts)};
\node[draw=poppurple, fill=poppurpleL, rounded corners, minimum width=4.2cm, minimum height=0.9cm, align=center, font=\bfseries] (k) at (2.5,-1.6) {Connecteurs\\ (2 pts)};
\node[draw=popgold, fill=popgoldL, rounded corners, minimum width=4.2cm, minimum height=0.9cm, align=center, font=\bfseries] (o) at (7.5,-1.6) {Orthographe\\ (2 pts)};
\node[draw=popdark, fill=popcream, rounded corners, minimum width=4.2cm, minimum height=0.9cm, align=center, font=\bfseries] (t) at (5,-3.2) {Total : 10 pts};
\draw[-{Stealth}, thick, popline] (c) -- (t);
\draw[-{Stealth}, thick, popline] (l) -- (t);
\draw[-{Stealth}, thick, popline] (p) -- (t);
\draw[-{Stealth}, thick, popline] (k) -- (t);
\draw[-{Stealth}, thick, popline] (o) -- (t);
\end{tikzpicture}
\legende{Grille d'auto-évaluation de la production écrite : cinq critères de 2 points chacun.}
\end{popfigure}

\textbf{Critères détaillés} : \emph{Contenu} = plan respecté (problème, conseils, conclusion) ; \emph{Lexique} = au moins 5 mots du thème correctement employés (article, orthographe) ; \emph{Passif} = au moins 3 formes correctes ; \emph{Connecteurs} = \all{obwohl} avec verbe final et \all{trotz} + génitif ; \emph{Orthographe} = majuscules aux noms, Umlaute, terminaisons.

\begin{savaistu}[Le passif valorise votre copie au Bac A]
Au baccalauréat A, l'expression écrite est notée sur la pertinence du contenu \emph{et} la correction grammaticale. Un passif bien construit — auxiliaire \all{werden} correctement conjugué, participe II à la fin, agent introduit par \all{von} + datif — montre à l'examinateur une maîtrise réelle de la syntaxe allemande et fait monter la note. À l'inverse, un passif « français » calqué mot à mot (\all{*Die Opfer sind geschützt} pour dire « les victimes sont protégées » au sens d'action) est sanctionné. Entraînez-vous donc à placer au moins deux ou trois passifs dans chaque rédaction !
\end{savaistu}

\begin{aretenir}[L'essentiel de la production guidée]
\begin{itemize}
\item Toujours rédiger un plan en 3 lignes avant d'écrire, puis relire : place du verbe, terminaisons, participe II final.
\item Passif : \all{werden} conjugué + participe II à la fin ; avec modal : modal + participe II + \all{werden}.
\item Concession : \all{obwohl} + verbe à la fin ; \all{trotz} + génitif.
\item À l'oral, utiliser les phrases-outils du commentaire pour donner son avis, concéder et conclure.
\end{itemize}
\end{aretenir}

\newpage

\coursec{Activité d'intégration}

\courssub{Objectifs de l'activité}
\par
Cette activité de synthèse vise à mobiliser l'ensemble des savoirs et savoir-faire travaillés lors de la leçon AL18 dans une tâche finale authentique, collaborative et créative. Elle permet à l'élève de :
\begin{itemize}
    \item \textbf{Réinvestir le lexique thématique} du harcèlement numérique (\cle{das Cybermobbing}, \cle{das Opfer}, \cle{der Täter}, \cle{die Anzeige}, \cle{schützen}, \cle{beweisen}, \cle{melden}, \cle{blockieren}) dans une production écrite cohérente.
    \item \textbf{Maîtriser le passif de procès (\all{Vorgangspassiv})} aux temps du présent et du prétérit, notamment avec les verbes modaux (\all{müssen, sollen, können, dürfen}) pour formuler des obligations, des recommandations et des interdictions claires.
    \item \textbf{Utiliser les connecteurs de concession et d'opposition} (\all{obwohl, trotzdem, dennoch, aber}) pour nuancer le discours, montrer la réalité du problème malgré l'existence de lois.
    \item \textbf{Développer des compétences transversales} : travail en équipe, conception graphique, expression orale continue (présentation de 2 minutes), esprit critique et citoyenneté numérique.
\end{itemize}

\courssub{Situation}
\par
Vous êtes élèves en classe de Terminale A au \textbf{Lycée Bilingue d'Etoug-Ebe} à Yaoundé. L'administration a lancé une campagne « \all{Eine Schule ohne Gewalt – Une école sans violence} » pour le premier trimestre. Le proviseur, M. \all{Atangana}, a confié aux classes d'allemand LV2 la réalisation d'affiches de sensibilisation en langue allemande qui seront exposées dans les couloirs du bâtiment principal, au CDI et à la salle informatique. L'objectif est d'alerter les camarades (germanophones et non-germanophones) sur les dangers du \all{Cybermobbing} et de leur donner des conduites à tenir concrètes.

\begin{textebox}[Contexte : Lancement du projet par le professeur d'allemand]
\all{
Liebe Schülerinnen und Schüler,
das Cybermobbing ist kein Randphänomen mehr – es trifft Schüler in Yaoundé genauso wie in Berlin oder München. WhatsApp-Gruppen, TikTok, Instagram: Die Täter sind oft anonym, die Opfer hilflos. Aber: Das Gesetz schützt euch! § 238 StGB (Nachstellung) und § 185-187 StGB (Beleidigung, Verleumdung) gelten auch im Netz.
Eure Aufgabe: Gestaltet eine Plakat-Aktion „Stoppt Cybermobbing!“ für unseren Schulflur.
Das Plakat muss enthalten:
\begin{enumerate}
    \item Einen starken Slogan (kurz, einprägsam, im Passiv).
    \item Fünf konkrete Handlungsempfehlungen im Passiv mit Modalverben (Präsens).
    \item Einen Satz der Konzession (Obwohl / Trotzdem), der zeigt: Gesetze allein reichen nicht.
    \item Wichtige Anlaufstellen mit Namen und Nummern (Schulpsychologe, Vertrauenslehrer, Polizei, Grüne Nummer).
\end{enumerate}
Arbeitet in Gruppen zu viert. Präsentationszeit: 2 Minuten pro Gruppe. Das beste Plakat wird laminiert und ausgehängt.
Viel Erfolg! \\
\emph{Ihr Deutschlehrer, Herr Ngono}
}
\end{textebox}

\courssub{Consignes}
\par
Travaillez en groupes de 4 élèves. Répartissez les rôles : \textbf{Rédacteur principal} (allemand), \textbf{Correcteur grammatical} (passif, modaux, connecteurs), \textbf{Graphiste/Mise en page} (visuel, lisibilité), \textbf{Orateur} (présentation orale). Suivez les étapes ci-dessous :

\begin{enumerate}
    \item \textbf{Analyse du cahier des charges (10 min) :} Relisez le message de M. Ngono (texte ci-dessus). Identifiez les 4 éléments obligatoires. Soulignez les verbes modaux au passif attendus (\all{müssen, sollen, können}). Notez le connecteur de concession imposé (\all{obwohl}).
    \item \textbf{Brainstorming lexical \& grammatical (15 min) :} Sur une feuille de brouillon, listez 10 mots-clés du thème (noms, verbes). Conjuguez au passif présent avec modaux : \all{schützen $\rightarrow$ müssen geschützt werden}, \all{melden $\rightarrow$ sollen gemeldet werden}, \all{speichern $\rightarrow$ müssen gespeichert werden}, \all{blockieren $\rightarrow$ können blockiert werden}, \all{strafen $\rightarrow$ müssen bestraft werden}. Formulez 2 phrases de concession avec \all{obwohl} + verbe à la fin et \all{trotzdem} en position 1 ou 3.
    \item \textbf{Rédaction de l'affiche (25 min) :} Rédigez la version finale sur papier A3 (fourni) ou sur tablette (Canva/PowerPoint). Respectez la structure imposée :
    \begin{itemize}
        \item \textbf{Slogan} (très grand, centré, passif + modal). Ex. : \all{Opfer müssen gehört und geschützt werden!}
        \item \textbf{5 Conseils} (liste numérotée, phrases complètes au passif + modal). Variez les modaux (\all{müssen, sollen, können}).
        \item \textbf{Phrase de concession} (encadrée ou en couleur). Structure : \all{Obwohl [subordonnée], [principale avec trotzdem/dennoch].}
        \item \textbf{Contacts} : \all{Ansprechpartner: Herr Ngono (Raum 12), Frau Mvondo (Schulpsychologin), Polizei Yaoundé (117), Grüne Nummer: 1500.}
    \end{itemize}
    \item \textbf{Révision croisée (10 min) :} Échangez l'affiche avec un autre groupe. Vérifiez : \textbf{Passif correct ?} (\all{werden} + Partizip II à la fin), \textbf{Modaux au bon endroit ?} (position 2), \textbf{Concession bien construite ?} (\all{obwohl} renvoie le verbe à la fin), \textbf{Orthographe/Umlauts/majuscules aux noms ?} Corrigez au stylo vert.
    \item \textbf{Préparation orale (10 min) :} L'orateur répète la présentation (2 min max) : lire le slogan, expliquer le choix des 5 conseils, lire la phrase de concession, citer les contacts. Les 3 autres membres gèrent l'affichage au tableau.
    \item \textbf{Présentation \& Vote (15 min) :} Chaque groupe présente. La classe note sur une grille (Contenu / Grammaire / Visuel / Oral). Vote à main levée pour la meilleure affiche. Remise du prix symbolique (bonbon \all{Haribo} / note de participation).
\end{enumerate}

\courssub{Ressources à mobiliser}
\par
Voici un mémo synthétique des outils linguistiques vus en classe pour réussir l'activité.

\begin{propriete}[Le Vorgangspassiv (Passif de procès) avec verbes modaux — Présent]
\par
\emph{Structure :} \textbf{Modal (conjugué, pos. 2) + Sujet + ... + Partizip II + \all{werden} (infinitif, pos. finale)}.
\par
\emph{Formation du Partizip II :} \all{ge- + radical + -t} (verbes faibles) / \all{ge- + radical modifié + -en} (verbes forts).
\par
\begin{center}
\begin{tabularx}{\linewidth}{|X|X|X|}
\hline
\rowcolor{popblueL} \textbf{Actif (Aktiv)} & \textbf{Passif avec modal (Vorgangspassiv)} & \textbf{Traduction} \\
\hline
\all{Man muss die Opfer schützen.} & \all{Die Opfer \textbf{müssen} geschützt \textbf{werden}.} & Les victimes doivent être protégées. \\
\hline
\all{Man soll Beweise speichern.} & \all{Beweise \textbf{sollen} gespeichert \textbf{werden}.} & Les preuves doivent être conservées. \\
\hline
\all{Man kann den Täter blockieren.} & \all{Der Täter \textbf{kann} blockiert \textbf{werden}.} & L'auteur peut être bloqué. \\
\hline
\all{Man darf Fotos nicht teilen.} & \all{Fotos \textbf{dürfen} nicht geteilt \textbf{werden}.} & Les photos ne doivent pas être partagées. \\
\hline
\end{tabularx}
\end{center}
\cle{Attention :} L'agent (\all{von wem}) est souvent omis dans les consignes générales. Le verbe \all{werden} reste à l'infinitif à la fin.
\end{propriete}

\begin{propriete}[Concession et Opposition : \all{obwohl} vs \all{trotzdem / dennoch}]
\par
\begin{center}
\begin{tabularx}{\linewidth}{|X|X|X|}
\hline
\rowcolor{poporangeL} \textbf{Connecteur} & \textbf{Structure syntaxique} & \textbf{Exemple (Cybermobbing)} \\
\hline
\all{obwohl} (bien que) & Subordonnée : verbe \textbf{rejeté à la fin}. & \all{Obwohl das Gesetz streng \textbf{ist}, gibt es viele Täter.} \\
\hline
\all{trotzdem} / \all{dennoch} (néanmoins) & Principale : verbe en \textbf{position 2} (sujet ou \all{trotzdem} en pos. 1). & \all{Das Gesetz ist streng. \textbf{Trotzdem} gibt es viele Täter.} \\
& & \all{Das Gesetz ist streng. Es gibt \textbf{trotzdem} viele Täter.} \\
\hline
\all{aber} (mais) & Coordination : verbe en position 2. Pas de rejet. & \all{Das Gesetz ist streng, \textbf{aber} es gibt viele Täter.} \\
\hline
\end{tabularx}
\end{center}
\cle{Piège francophone :} Ne pas écrire \all{*Obwohl das Gesetz streng ist, trotzdem gibt es...} (doublon incorrect). Choisir \textbf{soit} \all{obwohl} \textbf{soit} \all{trotzdem/dennoch/aber}.
\end{propriete}

\begin{definition}[Lexique clé — \all{Das Cybermobbing} (Nom, article, pluriel)]
\par
\begin{itemize}
    \item \all{das Cybermobbing} (pas de pluriel usuel) — le cyberharcèlement
    \item \all{das Opfer, die Opfer} — la victime
    \item \all{der Täter, die Täter} — l'auteur, le harceleur
    \item \all{die Anzeige, die Anzeigen} — la plainte, le signalement
    \item \all{der Beweis, die Beweise} — la preuve
    \item \all{die Belästigung, die Belästigungen} — le harcèlement
    \item \all{die Beleidigung, die Beleidigungen} — l'injure
    \item \all{verletzen} (verletzt, hat verletzt) — blesser (psychiquement)
    \item \all{melden} — signaler, déclarer
    \item \all{blockieren} — bloquer (un contact)
    \item \all{speichern} — sauvegarder, conserver
    \item \all{schützen} — protéger
    \item \all{anonym} — anonyme
    \item \all{der Vertrauenslehrer, die Vertrauenslehrer} — le professeur de confiance
    \item \all{die Schulpsychologin, die Schulpsychologinnen} — la psychologue scolaire
    \item \all{die Polizei} — la police
    \item \all{die Grüne Nummer, die Grünen Nummern} — le numéro vert (1500 au Cameroun)
\end{itemize}
\end{definition}

\courssub{Proposition de production et corrigé commenté}
\par
Voici une production modèle réalisée par un groupe type (niveau Terminale A LV2), suivie de l'analyse linguistique justifiant les choix.

\begin{exoresolu}[Affiche modèle : « Stoppt Cybermobbing! »]
\textbf{Énoncé :} Produire l'affiche finale prête à être affichée.
\tcblower
\textbf{Solution détaillée (Texte de l'affiche) :}

\begin{textebox}[Affiche finale — Groupe 3 : Marie, Jean, Amina, Paul]
\vspace{0.2cm}
\centering
\textbf{\Huge \all{STOPPT CYBERMOBBING!}} \\[0.3cm]
\textbf{\Large \all{Opfer müssen gehört und geschützt werden!}} \\[0.5cm]
\noindent\rule{\linewidth}{0.6pt}\par
\textbf{\large \all{5 REGELN FÜR DEINE SICHERHEIT:}} \\[0.2cm]
\begin{enumerate}
    \item \all{Beleidigende Nachrichten \textbf{müssen} sofort \textbf{gespeichert werden}.}
    \item \all{Der Täter \textbf{soll} bei der Plattform \textbf{gemeldet werden}.}
    \item \all{Das Profil des Täters \textbf{kann} einfach \textbf{blockiert werden}.}
    \item \all{Keine privaten Daten \textbf{dürfen} im Netz \textbf{veröffentlicht werden}.}
    \item \all{Die Eltern oder Vertrauenslehrer \textbf{müssen} informiert \textbf{werden}.}
\end{enumerate}
\vspace{0.2cm}
\noindent\rule{\linewidth}{0.6pt}\par
\centering
\textbf{\all{OBWOHL CYBERMOBBING IN DEUTSCHLAND UND KAMERUN VERBOTEN IST,}}\\
\textbf{\all{GIBT ES TROTZDEM NOCH VIELE OPFER. SCHWEIG NICHT!}} \\[0.5cm]
\noindent\rule{\linewidth}{0.6pt}\par
\textbf{\large \all{ANSPRECHPARTNER / HILFE:}} \\
\all{• Herr Ngono (Deutschlehrer, Raum 12)} \\
\all{• Frau Mvondo (Schulpsychologin, Büro B3)} \\
\all{• Polizei Yaoundé: \textbf{117}} \\
\all{• Grüne Nummer (MINJEC): \textbf{1500}} \\
\all{• \textbf{Nummer gegen Kummer: 116 111 (Deutschland)}}
\end{textebox}

\textbf{Commentaire linguistique du corrigé :}
\begin{enumerate}
    \item \textbf{Slogan :} \all{Opfer müssen gehört und geschützt werden.} $\rightarrow$ Passif présent avec modal \all{müssen} (obligation forte). Deux Partizip II coordonnés (\all{gehört} / \all{geschützt}) partageant l'auxiliaire \all{werden}. Efficace, percutant.
    \item \textbf{Conseil 1 :} \all{müssen ... gespeichert werden} (obligation de preuve). \all{sofort} (adverbe) placé après le modal, avant le Partizip II. Correct.
    \item \textbf{Conseil 2 :} \all{sollen ... gemeldet werden} (recommandation morale/institutionnelle). Préposition \all{bei} + datif (\all{der Plattform}). Correct.
    \item \textbf{Conseil 3 :} \all{können ... blockiert werden} (possibilité technique). \all{einfach} (adverbe de manière). Correct.
    \item \textbf{Conseil 4 :} \all{Keine privaten Daten dürfen im Netz veröffentlicht werden.} (interdiction). La négation est portée par l'article indéfini négatif \all{Keine} (nominatif pluriel), ce qui évite l'usage de \all{nicht} et place la focalisation sur l'objet. Structure élégante et correcte.
    \item \textbf{Conseil 5 :} \all{müssen ... informiert werden} (obligation de communication). Objet \all{die Eltern oder Vertrauenslehrer} au nominatif (sujet passif). Correct.
    \item \textbf{Phrase de concession :} \all{Obwohl Cybermobbing in Deutschland und Kamerun verboten ist, gibt es trotzdem noch viele Opfer.}
    \begin{itemize}
        \item Subordonnée \all{Obwohl ... ist} : verbe \all{ist} correctement rejeté à la fin. Sujet \all{Cybermobbing}. Compléments de lieu \all{in Deutschland und Kamerun}.
        \item Principale \all{gibt es trotzdem noch viele Opfer} : inversion sujet-verbe (\all{gibt es}) car la subordonnée est en position 1 (topicalisation). \all{trotzdem} en position 3 (après \all{es}) pour insister sur la concession. \all{noch} (encore) bien placé. \textbf{Excellente maîtrise de la topologie de la phrase allemande.}
    \end{itemize}
    \item \textbf{Contacts :} Vocabulaire précis (\all{Schulpsychologin}, \all{Vertrauenslehrer}), numéros locaux (Cameroun) et allemand (\all{Nummer gegen Kummer}) pour l'authenticité culturelle.
\end{enumerate}
\end{exoresolu}

\courssub{Bilan de l'activité}
\par
Cette activité d'intégration a permis de vérifier l'appropriation des trois piliers grammaticaux de la leçon AL18 dans un contexte de communication réel et motivant :
\begin{itemize}
    \item \textbf{Le passif modalisé} s'est révélé l'outil grammatical \emph{par excellence} pour formuler des consignes de sécurité impersonnelles, objectives et autoritaires (\all{müssen, sollen, dürfen}). Les élèves ont compris que le choix du modal nuance la force de l'injonction (obligation légale vs conseil technique vs interdiction).
    \item \textbf{La concession (\all{obwohl} / \all{trotzdem})} a permis d'élever le discours au-delà de la simple liste de conseils : elle introduit une réflexion critique sur l'écart entre la norme juridique et la réalité sociale, compétence attendue en Terminale (niveau B1+/B2).
    \item \textbf{L'ancrage camerounais} (numéros d'urgence locaux, noms de professeurs, contexte WhatsApp/TikTok) a garanti l'authenticité de la tâche et l'engagement des élèves.
\end{itemize}

\begin{attention}[Erreurs fréquentes observées lors de la correction croisée — Points de vigilance pour l'évaluation]
\begin{itemize}
    \item \textbf{Oubli du \all{werden} final :} \all{*Opfer müssen geschützt.} $\rightarrow$ Rappel : \cle{Vorgangspassiv = Modal + Partizip II + werden}.
    \item \textbf{Mauvaise place du verbe dans la subordonnée \all{obwohl} :} \all{*Obwohl Cybermobbing ist verboten...} $\rightarrow$ Verbe \textbf{obligatoirement à la fin}.
    \item \textbf{Confusion \all{trotzdem} / \all{obwohl} :} Doublon \all{*Obwohl ..., trotzdem ...}. Choix unique.
    \item \textbf{Majuscules manquantes sur les noms :} \all{*opfer, *Täter, *anzeige}. $\rightarrow$ Règle d'or : \cle{Tout nom = Majuscule}.
    \item \textbf{Accord du participe passé (français) vs Partizip II (allemand) :} L'allemand ne s'accorde pas en genre/nombre au passif (\all{die Opfer werden geschützt}, \all{der Täter wird bestraft}). Ne pas ajouter de \all{-e} ou \all{-en} d'accord.
\end{itemize}
\end{attention}

\begin{savaistu}[Citoyenneté numérique : Le cadre légal au Cameroun et en Allemagne]
\par
Saviez-vous que le Cameroun a adopté la \textbf{Loi n°2010/012 du 21 décembre 2010 relative à la cybercriminalité} ? Les articles relatifs au harcèlement par voie électronique (\all{Cybermobbing}) prévoient des peines d'emprisonnement et d'amende. En Allemagne, le \all{Strafgesetzbuch (StGB)} sanctionne la \all{Nachstellung (§ 238 StGB)} — le « stalking » numérique — jusqu'à 3 ans de prison, et les injures (\all{Beleidigung § 185}), calomnies (\all{Verleumdung § 186}) et diffamation (\all{üble Nachrede § 187}) sont poursuivies d'office si la victime est mineure. La plateforme \all{jugendschutz.net} et le \all{Nummer gegen Kummer (116 111)} offrent une aide gratuite et anonyme. Au Cameroun, le \all{MINJEC} (Ministère de la Jeunesse et de l'Éducation Civique) gère le \textbf{1500} (Grüne Nummer). \textbf{La loi protège, mais la prévention par l'éducation (vos affiches !) reste la première barrière.}
\end{savaistu}

\newpage

\coursec{Méthodes et exercices résolus}

\coursec{AL18 : La cybercriminalité 2 : le harcèlement numérique}

\courssub{Texte support : Cybermobbing – ein Fall aus Berlin}

\begin{textebox}[Article original — Contexte : Berlin, lycée professionnel]
\all{Die 16-jährige Lena besucht ein Berufskolleg in Berlin-Mitte. Seit drei Monaten wird sie über WhatsApp und Instagram systematisch beleidigt. Unbekannte Täter haben gefälschte Fotos von ihr erstellt und in Klassenchats veröffentlicht. „Ich hatte Angst, zur Schule zu gehen“, sagt Lena. Obwohl sie die Nachrichten löschte und die Nummern blockierte, gingen die Angriffe weiter. Die Täter nutzten immer neue Accounts.}

\all{Lena zog sich zurück, ihre Noten verschlechterten sich. Erst als ihr Vater die Veränderungen bemerkte, vertraute sie sich ihm an. Zusammen erstatteten sie Anzeige bei der Polizei. Der Beamte erklärte: „Bei Cybermobbing ist es wichtig, Beweise zu sichern: Screenshots machen, nicht einfach löschen.“ Die Polizei ermittelt nun wegen Beleidigung und Verleumdung nach § 185 ff. StGB.}

\all{Die Schulleitung reagierte: Es gab eine Projektwoche „Respekt im Netz“. Die Schüler lernten, wie man Passwörter schützt, Privatsphäre-Einstellungen prüft und wo man Hilfe findet (Nummer gegen Kummer: 116 111). Lena geht wieder zur Schule. „Das Mobbing ist nicht ganz weg“, sagt sie, „aber ich weiß jetzt: Ich bin nicht schuld, und ich habe Rechte.“}
\end{textebox}

\begin{savaistu}[Contexte Allemagne — Aide aux victimes]
En Allemagne, le \all{Cybermobbing} n'est pas un délit autonome, mais est poursuivi via les infractions classiques : \all{Beleidigung} (insulte, § 185 StGB), \all{Üble Nachrede} (calomnie, § 186), \all{Verleumdung} (dénonciation calomnieuse, § 187) ou \all{Nachstellung} (harcèlement/stalking, § 238). Les victimes peuvent porter plainte (\all{Anzeige erstatten}) en ligne (\all{Onlinewache}) ou au commissariat. Numéros d'urgence : \all{Nummer gegen Kummer} (116 111, anonyme et gratuit), \all{Telefonseelsorge} (0800 111 0 111). L'initiative \all{„Schule ohne Rassismus – Schule mit Courage„} mobilise les lycées.
\end{savaistu}

\courssub{Compréhension du texte}

\begin{methode}[Lire et exploiter un texte de presse sur le Cybermobbing]
\begin{enumerate}
\item \textbf{Première lecture globale.} Lis le titre, la source (ici : article de presse simulé) et le premier paragraphe. Repère le \cle{thème} (\all{das Cybermobbing}), le \cle{lieu} (\all{Berlin}), les \cle{personnages} (\all{das Opfer} Lena, \all{der Täter} / \all{die Täter}, \all{der Vater}, \all{die Polizei}, \all{die Schulleitung}) et le \cle{type de texte} (\all{Zeitungsartikel}, \all{Bericht}).
\item \textbf{Deuxième lecture détaillée.} Souligne les mots-clés du champ lexical : \all{die Anzeige}, \all{das Passwort}, \all{schützen}, \all{beleidigen}, \all{veröffentlichen}, \all{erstatten}, \all{ermitteln}, \all{die Verleumdung}, \all{der Account}. Entoure les mots transparents (\all{die Cyberkriminalität}, \all{das Internet}, \all{die Polizei}, \all{der StGB} = Code pénal).
\item \textbf{Repérage grammatical.} Note les phrases au \cle{passif} (\all{wird ... beleidigt}, \all{wurden ... veröffentlicht}, \all{wird ... erklärt}) et les connecteurs de \cle{concession} (\all{obwohl}, \all{Trotz}) : ils portent souvent la réponse aux questions de compréhension fine.
\item \textbf{Répondre aux questions.} Reformule avec tes propres mots, en phrases complètes, en citant un élément du texte. Une réponse commence par le sujet, jamais par \all{weil} seul.
\item \textbf{Vérification.} Relis : verbe conjugué en 2\up{e} position, majuscules à tous les noms, cas corrects après les prépositions (\all{bei der Polizei}, \all{mit den Eltern}, \all{von den Tätern}).
\end{enumerate}
\end{methode}

\begin{exoresolu}[Compréhension de l'écrit — Type Bac]
\textbf{Texte :} voir le texte support ci-dessus.\par
\textbf{Questions :}
\begin{enumerate}
\item Wer ist das Opfer und wie alt ist sie?
\item Was haben die Täter mit Fotos von Lena gemacht?
\item Warum half das Blockieren der Nummern nicht?
\item Was rät der Polizist den Opfern?
\item Welche Maßnahme hat die Schule ergriffen?
\end{enumerate}
\tcblower
\textbf{Raisonnement :}
\begin{enumerate}
\item Information explicite premier paragraphe : \all{Die 16-jährige Lena}.
\item Action des auteurs : \all{gefälschte Fotos ... erstellt und ... veröffentlicht} (passif possible dans la réponse : \all{wurden erstellt und veröffentlicht}).
\item Concession \all{Obwohl sie ... blockierte, gingen die Angriffe weiter} -> nouvelle comptes (\all{neue Accounts}).
\item Conseil direct (discours indirect) : \all{Beweise sichern, Screenshots machen}.
\item Réaction institutionnelle : \all{Projektwoche „Respekt im Netz“}.
\end{enumerate}
\textbf{Réponses rédigées :}\par
1. \all{Das Opfer ist Lena, eine 16-jährige Schülerin aus Berlin.} « La victime est Lena, une élève de 16 ans de Berlin. »\par
2. \all{Die Täter haben gefälschte Fotos von Lena erstellt und in Klassenchats veröffentlicht.} (Ou passif : \all{Gefälschte Fotos von Lena wurden erstellt und veröffentlicht.}) « Les auteurs ont créé de fausses photos de Lena et les ont publiées dans des chats de classe. »\par
3. \all{Das Blockieren half nicht, weil die Täter immer neue Accounts nutzten.} « Le blocage n'a pas aidé car les auteurs utilisaient toujours de nouveaux comptes. »\par
4. \all{Der Polizist rät, Beweise zu sichern und Screenshots zu machen, statt nur zu löschen.} « Le policier conseille de sécuriser des preuves et de faire des captures d'écran, au lieu de simplement supprimer. »\par
5. \all{Die Schulleitung hat eine Projektwoche „Respekt im Netz“ organisiert.} « La direction a organisé une semaine-projet "Respect sur le Net". »\par
\textbf{Conclusion :} Chaque réponse est une phrase complète, verbe en 2\up{e} position, vocabulaire du texte repris correctement (cas, temps).
\end{exoresolu}

\begin{exercice}[Compréhension — Entraînement personnel]
Répondez en allemand, en phrases complètes, d'après le texte support.
\begin{enumerate}
\item Wie hat sich Lenas Verhalten verändert? (Deux éléments)
\item Welche Straftaten liegen laut Polizei vor?
\item Was haben die Schüler in der Projektwoche gelernt? (Trois éléments)
\item Wie fühlt sich Lena am Ende? Citez un passage du texte.
\end{enumerate}
\end{exercice}

\courssub{Lexique thématique : das Cybermobbing}

\begin{propriete}[Vocabulaire de base — Noms (article, pluriel, genre) et verbes (construction)]
\begin{center}
\begin{tabularx}{\linewidth}{|X|X|X|}
\hline
\rowcolor{popblueL} \textbf{Deutsch (Artikel, Plural)} & \textbf{Français} & \textbf{Construction / Famille / Exemple} \\
\hline
\all{das Opfer, die Opfer} & la victime & \all{Opfer sein} (+ génitif : \all{des Mobbing}) ; \all{jemanden zum Opfer fallen} \\
\all{der Täter, die Täter} & l'auteur, le harceleur & \all{der Täter} (masc.) ; fém. \all{die Täterin} ; verbe \all{begehen} (une infraction) \\
\all{die Anzeige, die Anzeigen} & la plainte, la dénonciation & \all{Anzeige erstatten} (porter plainte) ; \all{Anzeige gegen jdn. erstatten} (+ accusatif) \\
\all{das Passwort, die Passwörter} & le mot de passe & \all{Passwort ändern / weitergeben / vergessen} ; \all{sicheres Passwort} \\
\all{die Nachricht, die Nachrichten} & le message & \all{Nachricht löschen / blockieren / speichern} ; \all{Sprachnachricht} \\
\all{der Account, die Accounts} & le compte & \all{Account erstellen / hacken / löschen} ; \all{fake Account} \\
\all{der Screenshot, die Screenshots} & la capture d'écran & \all{Screenshot machen / speichern / als Beweis sichern} \\
\all{die Verleumdung, die Verleumdungen} & la calomnie, la diffamation & \all{wegen Verleumdung anzeigen} ; verbe \all{verleumden} \\
\all{die Beleidigung, die Beleidigungen} & l'insulte & \all{wegen Beleidigung anzeigen} ; verbe \all{beleidigen} \\
\all{der Beweis, die Beweise} & la preuve & \all{Beweise sichern / vorlegen / vernichten} \\
\hline
\end{tabularx}
\end{center}
\end{propriete}

\begin{propriete}[Verbes clés et constructions]
\begin{center}
\begin{tabularx}{\linewidth}{|X|X|X|}
\hline
\rowcolor{popblueL} \textbf{Verbe} & \textbf{Sens / Construction} & \textbf{Exemple (Présent / Prétérit)} \\
\hline
\all{beleidigen} & insulter qqn (+ Akk) & \all{Er beleidigt sie. / Er beleidigte sie.} \\
\all{veröffentlichen} & publier, rendre public (+ Akk) & \all{Sie veröffentlicht das Foto. / Sie veröffentlichte es.} \\
\all{schützen} & protéger qqn/qqch (+ Akk) & \all{Man muss Kinder schützen. / Man schützte sie.} \\
\all{erstatten} & déposer (plainte) — \all{Anzeige erstatten} & \all{Sie erstattet Anzeige. / Sie erstattete Anzeige.} \\
\all{ermitteln} & enquêter (+ \all{wegen} + génitif) & \all{Die Polizei ermittelt. / Die Polizei ermittelte.} \\
\all{blockieren} & bloquer (+ Akk) & \all{Er blockiert die Nummer. / Er blockierte sie.} \\
\all{weitersagen / weitergeben} & divulguer, transmettre (part. sép.) & \all{Gib das Passwort nicht weiter!} \\
\all{sich vertrauen} & se confier à qqn (+ Dat) & \all{Sie vertraute sich ihrem Vater an.} \\
\hline
\end{tabularx}
\end{center}
\end{propriete}

\begin{methode}[Mémoriser et utiliser le lexique thématique]
\begin{enumerate}
\item \textbf{Apprends chaque nom avec article et pluriel} : \all{das Opfer, die Opfer} ; \all{der Täter, die Täter} ; \all{die Anzeige, die Anzeigen} ; \all{das Passwort, die Passwörter} ; \all{die Nachricht, die Nachrichten} ; \all{der Account, die Accounts}.
\item \textbf{Apprends les verbes avec leur construction casuelle} : \all{jemanden beleidigen} (Akk) ; \all{jemanden schützen} (Akk) ; \all{Anzeige erstatten} (Akk pour \all{Anzeige}) ; \all{sich jemandem anvertrauen} (Dat) ; \all{ermitteln wegen + Genitiv}.
\item \textbf{Construis des familles de mots (dérivation)} :
\begin{itemize}
\item \all{mobben} (v) $\rightarrow$ \all{das Mobbing} (n) $\rightarrow$ \all{das Cybermobbing} (n) $\rightarrow$ \all{der Mobber} / \all{die Mobberin}
\item \all{die Tat} $\rightarrow$ \all{der Täter} / \all{die Täterin} $\rightarrow$ \all{täterlos} (adj)
\item \all{anzeigen} $\rightarrow$ \all{die Anzeige} $\rightarrow$ \all{anzeigepflichtig}
\item \all{schützen} $\rightarrow$ \all{der Schutz} $\rightarrow$ \all{schützenswert} $\rightarrow$ \all{der Datenschutz}
\end{itemize}
\item \textbf{Réutilise le lexique dans des phrases personnelles} liées à ton contexte (école, réseaux sociaux, téléphone portable, Yaoundé/Douala).
\end{enumerate}
\end{methode}

\begin{exoresolu}[Vocabulaire — Phrases à compléter (Type Bac)]
Complète avec le mot qui convient \textbf{à la forme correcte} (cas, nombre) : \all{das Passwort, der Täter, die Anzeige, schützen, das Opfer, die Nachricht, blockieren, Beweis}.\par
1. Lena ist \ldots\ des Cybermobbings. \quad 2. Die Polizei sucht \ldots\ . \quad 3. Die Familie erstattet \ldots\ . \quad 4. Man muss Kinder im Internet \ldots\ . \quad 5. Gib niemals dein \ldots\ weiter! \quad 6. Sie hat die beleidigenden \ldots\ gelöscht. \quad 7. Als \ldots\ hat er die Nummern der Täter. \quad 8. Für die Polizei sind Screenshots wichtige \ldots\ .
\tcblower
\textbf{Raisonnement :} Identifier la fonction (sujet, COD, prépositionnel) et le nombre.\par
1. \all{Lena ist \textbf{das Opfer} des Cybermobbings.} (Prédicat après \all{sein} $\rightarrow$ Nominatif ; \all{des Cybermobbings} = Génitif).\par
2. \all{Die Polizei sucht \textbf{den Täter}.} (\all{suchen} + Accusatif ; masculin $\rightarrow$ \all{den}).\par
3. \all{Die Familie erstattet \textbf{Anzeige}.} (Expression figée \all{Anzeige erstatten}, \all{Anzeige} = Accusatif singulier sans article).\par
4. \all{Man muss Kinder im Internet \textbf{schützen}.} (Infinitif en fin de phrase après modal \all{muss}).\par
5. \all{Gib niemals dein \textbf{Passwort} weiter!} (\all{weitergeben} séparable, \all{Passwort} = Accusatif neutre ; impératif).\par
6. \all{Sie hat die beleidigenden \textbf{Nachrichten} gelöscht.} (Accusatif pluriel, adjectif \all{beleidigend} + \all{-en}).\par
7. \all{Als \textbf{Beweis} hat er die Nummern der Täter.} (Prédicat après \all{als} $\rightarrow$ Nominatif ; ou Accusatif si \all{als} = comme preuve. Ici : \all{Als Beweis} [ Nominatif/Accusatif ] \all{dienen}... mais phrase elliptique. Mieux : \all{Er hat die Nummern als \textbf{Beweis}.} $\rightarrow$ Accusatif).\par
8. \all{Für die Polizei sind Screenshots wichtige \textbf{Beweise}.} (\all{Für} + Accusatif ; \all{wichtige Beweise} = Nominatif pluriel attribut du sujet).\par
\textbf{Conclusion :} Le bon mot ne suffit pas : il faut accorder le cas (\all{den Täter}), le nombre (\all{Nachrichten}, \all{Beweise}) et placer correctement la particule séparable (\all{weiter}).
\end{exoresolu}

\begin{exercice}[Vocabulaire — Entraînement]
Complétez les phrases avec le mot approprié (forme correcte) : \all{die Täterin, das Mobbing, verleumden, der Datenschutz, sichern, die Schulleitung, anonym, das Opfer}.
\begin{enumerate}
\item \all{Nicht nur der Täter, auch \ldots\ kann strafrechtlich verfolgt werden.}
\item \all{Das \ldots\ fand nicht nur in der Schule statt, sondern auch im Internet.}
\item \all{Man darf niemanden \ldots\, das ist eine Straftat.}
\item \all{Opfer sollten Beweise \ldots\, bevor sie Nachrichten löschen.}
\item \all{Die \ldots\ hat eine Präventionsveranstaltung organisiert.}
\item \all{Viele Täter agieren \ldots\ hinter gefälschten Profilen.}
\item \all{Ein starkes \ldots\ schützt vor Identitätsdiebstahl.}
\item \all{Das \ldots\ leidet oft unter Angst und Schlafstörungen.}
\end{enumerate}
\end{exercice}

\courssub{Grammaire 1 : Le passif (Vorgangspassiv) — Présent et Prétérit}

\begin{propriete}[Formation du Vorgangspassiv (Passif d'action / de procès)]
Le \cle{Vorgangspassiv} décrit une \textbf{action} ou un \textbf{processus} en cours ou accompli. Il s'oppose au \all{Zustandspassiv} (état résultant : \all{das Foto ist veröffentlicht}).
\begin{itemize}
\item \textbf{Forme :} Auxiliaire \all{werden} (conjugué) + \cle{Participe II} du verbe lexical (à la fin de la phrase).
\item \textbf{Sujet passif} = COD de la phrase active (devient Nominatif).
\item \textbf{Complément d'agent} (optionnel) = Sujet actif $\rightarrow$ \all{von} + \textbf{Datif} (personne) ou \all{durch} + \textbf{Accusatif} (moyen/instrument, moins fréquent au Bac).
\item \textbf{Condition :} Seuls les verbes \textbf{transitifs} (avec COD / Akkusativobjekt) forment un passif personnel. Verbes intransitifs (\all{schlafen, gehen, sein}) $\rightarrow$ passif impersonnel seulement (\all{es wird getanzt}).
\end{itemize}
\end{propriete}

\begin{exemplebox}[Tableau de conjugaison — Verbe modèle : \all{beleidigen} (insulter)]
\begin{center}
\begin{tabular}{|l|l|l|l|}
\hline
\rowcolor{popblueL} \textbf{Personne} & \textbf{Présent (Gegenwart)} & \textbf{Prétérit (Präteritum)} & \textbf{Participe II} \\
\hline
ich & \all{werde beleidigt} & \all{wurde beleidigt} & \cellcolor{popblueL!20} \all{beleidigt} \\
du & \all{wirst beleidigt} & \all{wurdest beleidigt} & \\
er / sie / es & \all{wird beleidigt} & \all{wurde beleidigt} & \\
wir & \all{werden beleidigt} & \all{wurden beleidigt} & \\
ihr & \all{werdet beleidigt} & \all{wurdet beleidigt} & \\
sie / Sie & \all{werden beleidigt} & \all{wurden beleidigt} & \\
\hline
\end{tabular}
\end{center}
\emph{Remarque :} Le participe II \all{beleidigt} ne change jamais. Seul \all{werden} se conjugue. Pour les verbes forts/irréguliers : \all{veröffentlichen $\rightarrow$ veröffentlicht} (faible) ; \all{schreiben $\rightarrow$ geschrieben} (fort) ; \all{sehen $\rightarrow$ gesehen} (irrégulier).
\end{exemplebox}

\begin{methode}[Transformer une phrase active en passif (Personnel)]
\begin{enumerate}
\item \textbf{Repérer le temps} de la phrase active (Présent ou Prétérit).
\item \textbf{Conjuguer \all{werden}} à ce même temps, à la personne du nouveau sujet (l'ancien COD).
\item \textbf{Nouveau sujet} = Ancien COD (passer à l'article Nominatif : \all{den Täter} $\rightarrow$ \all{der Täter} ; \all{das Mädchen} $\rightarrow$ \all{das Mädchen}).
\item \textbf{Participe II} du verbe lexical à la \textbf{fin de la phrase}.
\item \textbf{Ancien sujet} $\rightarrow$ Complément d'agent \all{von} + \textbf{Datif} (\all{der Täter} $\rightarrow$ \all{von dem Täter} / \all{von den Tätern} ; \all{die Polizei} $\rightarrow$ \all{von der Polizei}). Si l'agent est inconnu/indifférent, on l'omet.
\end{enumerate}
\end{methode}

\begin{exoresolu}[Transformation Actif $\rightarrow$ Passif — Type Bac]
Mets les phrases au passif (Vorgangspassiv), \textbf{au même temps}. Omet l'agent si non pertinent.\par
1. \all{Die Täter veröffentlichen das Foto.} (Présent) \quad 2. \all{Ein Schüler beleidigte das Mädchen.} (Prétérit) \quad 3. \all{Die Polizei verhaftete den Täter.} (Prétérit) \quad 4. \all{Jemand hat das Passwort gestohlen.} (Perfekt — révision) \quad 5. \all{Man sollte das Problem melden.} (Modal + Infinitif — passif modal)
\tcblower
\textbf{Raisonnement \& Solutions :}\par
1. Présent. COD : \all{das Foto} (neutre). $\rightarrow$ \all{\textbf{Das Foto wird von den Tätern veröffentlicht.}} (\all{von} + Datif pluriel \all{den Tätern}).\par
2. Prétérit. COD : \all{das Mädchen} (neutre). $\rightarrow$ \all{\textbf{Das Mädchen wurde von einem Schüler beleidigt.}} (\all{von} + Datif \all{einem Schüler}).\par
3. Prétérit. COD : \all{den Täter} (masc.) $\rightarrow$ Sujet \all{der Täter}. $\rightarrow$ \all{\textbf{Der Täter wurde von der Polizei verhaftet.}} (\all{von} + Datif fém. \all{der Polizei}).\par
4. Perfekt. $\rightarrow$ \all{\textbf{Das Passwort ist gestohlen worden.}} (Passif Perfekt : \all{sein} + Part. II + \all{worden} ; agent omis).\par
5. Modal \all{sollen}. $\rightarrow$ \all{\textbf{Das Problem sollte gemeldet werden.}} (Passif modal : \all{modal} + \all{Part. II} + \all{werden}).\par
\textbf{Conclusion :} Trois réflexes : bon temps de \all{werden}, participe II final, complément d'agent en \all{von} + datif.
\end{exoresolu}

\begin{exercice}[Passif — Entraînement]
Transformez au passif (Vorgangspassiv), même temps. Omettez l'agent si "man" ou "jemand".
\begin{enumerate}
\item \all{Man mobbt viele Jugendliche im Netz.} (Présent)
\item \all{Die Schule organisierte eine Projektwoche.} (Prétérit)
\item \all{Jemand hat das Video geteilt.} (Perfekt)
\item \all{Die Eltern müssen das Kind schützen.} (Modal présent)
\item \all{Der Richter verurteilte den Täter zu einer Geldstrafe.} (Prétérit)
\end{enumerate}
\end{exercice}

\courssub{Grammaire 2 : La concession et l'opposition}

\begin{propriete}[Connecteurs de concession et d'opposition — Syntaxe et Cas]
\begin{center}
\begin{tabularx}{\linewidth}{|X|X|X|X|}
\hline
\rowcolor{popblueL} \textbf{Connecteur} & \textbf{Structure syntaxique} & \textbf{Place du verbe} & \textbf{Exemple} \\
\hline
\all{\cle{obwohl}} & Subordonnée (Sujet + ... + Verbe conjugué) & Verbe \textbf{à la fin} de la subordonnée. Si subordonnée en tête : verbe principal en \textbf{1\up{re} position} (juste après la virgule). & \all{Obwohl sie Angst \textbf{hatte}, \textbf{ging} sie zur Polizei.} \\
\all{\cle{trotz}} + Génitif & Groupe nominal (Article/Adj/Nom au \textbf{Génitif}) & Pas de verbe. La principale suit ordre normal (V2). & \all{Trotz \textbf{der} Angst \textbf{ging} sie zur Polizei.} \\
\all{\cle{aber}} & Coordination (Phrase principale + \all{aber} + Phrase principale) & Verbe en \textbf{2\up{e} position} dans chaque principale. & \all{Sie hatte Angst, \textbf{aber} sie \textbf{ging} zur Polizei.} \\
\all{\cle{dennoch}} / \all{dennoch} & Adverbe conjonctif (début phrase principale) & Verbe en \textbf{1\up{re} position} (inversion). & \all{Sie hatte Angst. \textbf{Dennoch} \textbf{ging} sie zur Polizei.} \\
\hline
\end{tabularx}
\end{center}
\emph{Attention :} \all{trotz} régit le \textbf{Génitif} (écrit standard). À l'oral : souvent Datif (\all{trotz dem Regen}), mais \textbf{écrit Bac = Génitif}. \all{obwohl} ne se combine jamais avec \all{aber} dans la même phrase (\all{Obwohl ... , aber ...} = FAUTE).
\end{propriete}

\begin{methode}[Choisir et construire la concession / opposition]
\begin{enumerate}
\item \textbf{Analyse le français :} Y a-t-il un verbe conjugué après "bien que / quoique" ? $\rightarrow$ \all{obwohl} + subordonnée. Y a-t-il seulement un nom/groupe nominal après "malgré / en dépit de" ? $\rightarrow$ \all{trotz} + Génitif. Est-ce une simple opposition "mais" ? $\rightarrow$ \all{aber} (ou \all{dennoch} en début de 2\up{e} phrase).
\item \textbf{Vérifie le cas après \all{trotz} :} \all{die Angst} $\rightarrow$ \all{der Angst} ; \all{die Drohungen} $\rightarrow$ \all{der Drohungen} ; \all{das Verbot} $\rightarrow$ \all{des Verbots} ; \all{seine Wut} $\rightarrow$ \all{seiner Wut}.
\item \textbf{Place le verbe :} Subordonnée \all{obwohl} en tête $\rightarrow$ Verbe principal juste après la virgule. \all{Dennoch} en tête $\rightarrow$ Inversion (Verbe - Sujet).
\item \textbf{Particules séparables :} En fin de subordonnée (\all{obwohl er aufgibt}) ou de principale (\all{er gibt auf}), jamais au milieu.
\end{enumerate}
\end{methode}

\begin{exoresolu}[Thème — Traduction concession / opposition (Type Bac)]
Traduis en allemand :\par
1. « Bien que Lena ait supprimé les messages, le harcèlement a continué. »\par
2. « Malgré la peur, elle a parlé avec ses parents. »\par
3. « Le harceleur était un camarade de classe, mais personne ne le savait. »\par
4. « Malgré les menaces, elle a porté plainte. »\par
5. « Elle avait honte, néanmoins elle en a parlé à un prof. »
\tcblower
\textbf{Raisonnement \& Solutions :}\par
1. Deux verbes (\all{supprimé}, \all{a continué}) $\rightarrow$ \all{obwohl} + Prétérit. \all{Obwohl Lena die Nachrichten \textbf{löschte}, \textbf{ging} das Mobbing \textbf{weiter.}} (\all{weitergehen} séparable).\par
2. Nom seul (\all{la peur}) $\rightarrow$ \all{trotz} + Génitif. \all{Trotz \textbf{der} Angst \textbf{sprach} sie mit ihren Eltern.} (\all{mit} + Datif pluriel \all{ihren Eltern}).\par
3. Opposition simple $\rightarrow$ \all{aber}. \all{Der Täter war ein Klassenkamerad, \textbf{aber} niemand \textbf{wusste} das.} (\all{wissen} Prétérit \all{wusste}).\par
4. Nom pluriel (\all{les menaces}) $\rightarrow$ \all{trotz} + Génitif pluriel. \all{Trotz \textbf{der} Drohungen \textbf{erstattete} sie Anzeige.} (\all{erstatten} Prétérit).\par
5. \all{dennoch} (adverbe) pour "néanmoins". \all{Sie schämte sich. \textbf{Dennoch} \textbf{sprach} sie mit einem Lehrer.}\par
\textbf{Conclusion :} Le choix dépend de la nature grammaticale (verbe vs nom). La place du verbe est le critère n°1 noté au Bac.
\end{exoresolu}

\begin{exercice}[Concession / Opposition — Entraînement]
Complétez avec \all{obwohl}, \all{trotz} (+ génitif), \all{aber} ou \all{dennoch}. Conjuguez les verbes entre parenthèses au prétérit.
\begin{enumerate}
\item \all{\ldots\ sie Angst \textbf{hatte}, \textbf{ging} sie zur Polizei.}
\item \all{\ldots\ \textbf{der} Angst \textbf{ging} sie zur Polizei.}
\item \all{Sie hatte Angst, \ldots\ sie \textbf{ging} zur Polizei.}
\item \all{Sie hatte Angst. \ldots\ \textbf{ging} sie zur Polizei.}
\item \all{\ldots\ \textbf{des} Verbots \textbf{postete} er das Foto.}
\item \all{Er wusste, dass es verboten war, \ldots\ \textbf{postete} er es.}
\end{enumerate}
\end{exercice}

\courssub{Méthode du commentaire et commentaire modèle}

\begin{methode}[Rédiger un commentaire structuré (Type Bac — 10-12 lignes)]
\begin{enumerate}
\item \textbf{Introduction (2 phrases) :} Présenter le document (type, source, thème, date/lieu si donnés). Annoncer le problème (\all{Problematik}) : \all{Der Text thematisiert das Cybermobbing am Beispiel von Lena\ldots}
\item \textbf{Analyse du contenu (4-5 phrases) :} Résumer les faits clés dans l'ordre chronologique/logique. Utiliser le \textbf{passif} pour les faits subis (\all{Sie wurde beleidigt}), la \textbf{concession} pour les obstacles (\all{Obwohl sie blockierte\ldots}). Citer des détails (durée, plateformes, réaction police/école).
\item \textbf{Analyse de la forme / langue (2-3 phrases) :} Repérer 2-3 procédés linguistiques : passif (effet d'objectivité, focus victime), concession (résilience), vocabulaire juridique (\all{Anzeige, Verleumdung, StGB}), discours direct/indirect (\all{sagt Lena}, \all{erklärte der Beamte}).
\item \textbf{Conclusion / Ouverture personnelle (2 phrases) :} Donner ton avis (\all{Meiner Meinung nach\ldots}), lier au contexte camerounais ou général, proposer une solution (\all{Man sollte\ldots}, \all{Es ist wichtig, dass\ldots}).
\end{enumerate}
\end{methode}

\begin{exoresolu}[Commentaire modèle — Type Bac]
\textbf{Consigne :} Commentez le texte \all{Cybermobbing – ein Fall aus Berlin} en montrant comment la victime sort de l'isolement. (Environ 12 lignes).\par
\tcblower
\textbf{Commentaire :}\par
\all{Der vorliegende Zeitungsbericht thematisiert das Cybermobbing am konkreten Fall der 16-jährigen Lena aus Berlin. Das Problem: Jugendliche werden systematisch im Netz beleidigt und bloßgestellt, oft anonym und über lange Zeit.}

\all{Der Text schildert chronologisch das Leid der Opfer: Lena wird monatelang über WhatsApp und Instagram beleidigt, gefälschte Fotos werden veröffentlicht. Obwohl sie technische Gegenmaßnahmen ergreift (Löschen, Blockieren), gelingt es ihr nicht, die Angriffe zu stoppen, da die Täter neue Accounts nutzen. Erst die Intervention des Vaters bricht den Teufelskreis: Zusammen erstatten sie Anzeige. Die Polizei betont die Wichtigkeit der Beweissicherung (Screenshots). Die Schule reagiert präventiv mit einer Projektwoche „Respekt im Netz“, wo Schüler lernen, Passwörter zu schützen und Hilfe zu suchen (Nummer gegen Kummer).}

\all{Sprachlich nutzt der Text das Vorgangspassiv („wird beleidigt“, „wurden veröffentlicht“), um den Fokus auf das Opfer und das Geschehen zu lenken, nicht auf die Täter. Der Konnektor „obwohl“ markiert die Ohnmacht der technischen Selbsthilfe. Fachbegriffe wie „Anzeige“, „Verleumdung“, „StGB“ verleihen juristische Autorität.}

\all{Meiner Meinung nach zeigt der Fall exemplarisch: Cybermobbing ist kein Kavaliersdelikt, sondern Straftat. Wichtig ist, dass Opfer nicht schweigen, sondern Beweise sichern und Vertrauenspersonen (Eltern, Lehrer, Polizei) einbeziehen. Auch in Kamerun sollten Schulen Präventionsworkshops anbieten, damit Jugendliche ihre Rechte im digitalen Raum kennen.}
\end{exoresolu}

\courssub{Production guidée et entraînement}

\begin{methode}[Rédiger une production guidée (court paragraphe / article de journal scolaire)]
\begin{enumerate}
\item \textbf{Analyse la consigne :} Sujet imposé ? Nombre de phrases ? Contraintes grammaticales (passif, concession, vocabulaire précis) ?
\item \textbf{Plan en 4 étapes :} 1. Constat / Accroche (\all{Cybermobbing ist\ldots}). 2. Conseils aux victimes (\all{Man sollte\ldots}, \all{Opfer müssen\ldots}). 3. Rôle de l'institution (école, police, loi). 4. Opinion personnelle / Appel (\all{Ich finde\ldots}, \all{Wir alle müssen\ldots}).
\item \textbf{Écriture :} Phrases courtes, correctes. \textbf{Check-list obligatoire avant de rendre :}
\begin{itemize}
\item Verbe en 2\up{e} position (principale) / fin (subordonnée).
\item Majuscules à \textbf{tous} les noms.
\item Particules séparables en fin de phrase (\all{weitergeben, anrufen, mitmachen}).
\item Cas après prépositions : \all{mit} (Dat), \all{von} (Dat), \all{bei} (Dat), \all{wegen} (Gen), \all{trotz} (Gen).
\item Passif : \all{werden} + Part. II (pas \all{sein} pour l'action !).
\item Concession : \all{obwohl} (verbe fin) / \all{trotz} (Génitif) / \all{aber} (V2).
\end{itemize}
\end{enumerate}
\end{methode}

\begin{exoresolu}[Production guidée modèle — Type Bac]
\textbf{Consigne :} Rédige un court paragraphe (6-7 phrases) : \all{„Was kann man gegen Cybermobbing tun?“} Utilise \textbf{au moins une phrase au passif} et \textbf{un connecteur de concession} (\all{obwohl}, \all{trotz}, \all{aber}, \all{dennoch}).\par
\tcblower
\textbf{Plan :} 1. Constat général. 2. Conseil technique (passif). 3. Concession (peur vs action). 4. Rôle école/parents. 5. Appel à la loi. 6. Opinion personnelle.\par
\textbf{Production :}\par
\all{Cybermobbing ist ein großes Problem für viele Jugendliche in Deutschland und weltweit. Oft werden beleidigende Nachrichten und Fotos im Internet veröffentlicht, obwohl die Opfer das nicht wollen. Meiner Meinung nach sollte man sein Passwort niemals weitergeben und die Privatsphäre-Einstellungen prüfen. Wenn jemand gemobbt wird, muss er sofort mit den Eltern oder Lehrern sprechen. Trotz der Angst sollte das Opfer Anzeige bei der Polizei erstatten. Ich finde es wichtig, dass Schulen und Familien die Jugendlichen besser schützen und über Gesetze aufklären.}\par
\textbf{Justifications grammaticales :}
\begin{itemize}
\item \all{werden ... veröffentlicht} = Passif présent pluriel (sujet \all{Nachrichten und Fotos}).
\item \all{obwohl die Opfer das nicht wollen} = Subordonnée concession, verbe \all{wollen} à la fin.
\item \all{weitergeben} = Particule séparable en fin de proposition infinitive (\all{sollte ... weitergeben}).
\item \all{Wenn jemand gemobbt wird} = Subordonnée temporelle, passif présent, verbe \all{wird} à la fin.
\item \all{muss er sofort ... sprechen} = Principale après subordonnée : verbe \all{muss} en 1\up{re} position (inversion).
\item \all{Trotz der Angst} = \all{trotz} + Génitif (\all{die Angst} $\rightarrow$ \all{der Angst}) ; verbe \all{sollte} en 2\up{e} position.
\item \all{dass Schulen ... schützen} = Subordonnée \all{dass}, verbe \all{schützen} à la fin.
\end{itemize}
\textbf{Conclusion :} Le paragraphe respecte toutes les contraintes (passif, concession, vocabulaire thème), reste simple et sans faute : stratégie gagnante au Bac.
\end{exoresolu}

\begin{exercice}[Production écrite — Entraînement Bac]
\textbf{Sujet :} \all{„Schutz vor Cybermobbing – Was kann die Schule tun?“} Rédigez un court article pour le journal du lycée (8-10 phrases).\par
\textbf{Contraintes obligatoires :}
\begin{itemize}
\item Utilisez \textbf{3 noms du vocabulaire thème} (avec articles corrects).
\item Intégrez \textbf{2 phrases au passif} (présent ou prétérit).
\item Utilisez \textbf{2 connecteurs de concession différents} (\all{obwohl}, \all{trotz}, \all{aber}, \all{dennoch}).
\item Donnez votre avis avec \all{Meiner Meinung nach} ou \all{Ich finde, dass\ldots}.
\end{itemize}
\end{exercice}

\begin{corrige}[Exercice Production écrite]
\textbf{Proposition de correction (une possibilité parmi d'autres) :}\par
\all{Schutz vor Cybermobbing – Was kann die Schule tun? Cybermobbing betrifft viele Schüler an unserer Schule. Oft werden die Opfer im Klassenchat beleidigt, obwohl die Lehrer das nicht sehen. Meiner Meinung nach sollte die Schule eine AG „Medienkompetenz“ anbieten. Dabei lernen die Jugendlichen, wie man ein sicheres Passwort erstellt und Beweise sichert. Trotz guter Prävention gibt es manchmal Täter. Dann muss die Schulleitung konsequent handeln und die Eltern informieren. Ich finde es wichtig, dass jede Schule einen Vertrauenslehrer für digitale Gewalt benennt. Nur so können wir das Netz sicherer machen.}
\end{corrige}

\begin{attention}[Les 5 erreurs qui coûtent des points au Bac — Récapitulatif]
\begin{enumerate}
\item \textbf{Confondre \all{werden} et \all{sein} au passif.} Le \cle{Vorgangspassiv} (action) exige \all{werden} : \all{Sie \textbf{wurde} beleidigt}. \all{Sie war beleidigt} = \cle{Zustandspassiv} (état : « elle était insultée / offensée »), ne raconte pas l'action. \textbf{Test :} Peut-on ajouter \all{von ...} ? Oui $\rightarrow$ \all{werden}.
\item \textbf{Oublier le verbe en fin de subordonnée avec \all{obwohl}.} Faute type : \all{Obwohl sie \textbf{hatte} Angst\ldots} — Correct : \all{Obwohl sie Angst \textbf{hatte},\ldots} Et après la virgule : verbe principal en 1\up{re} position : \all{\ldots, \textbf{ging} sie zur Polizei.}
\item \textbf{Calquer le français « malgré de / malgré les ».} \all{trotz} se construit \textbf{sans préposition} et avec le \textbf{Génitif} : \all{trotz \textbf{der} Drohungen}, \all{trotz \textbf{des} Verbots}. Jamais \all{trotz von den Drohungen}.
\item \textbf{Majuscules et articles des noms.} \all{das Opfer}, \all{der Täter}, \all{die Anzeige}, \all{das Passwort}, \all{die Nachricht} prennent \textbf{toujours} la majuscule. Apprends l'article \textbf{avec} le nom, car il commande les cas (\all{von dem Täter}, \all{bei der Polizei}, \all{wegen der Verleumdung}).
\item \textbf{Faux amis et particules séparables.}
\begin{itemize}
\item \all{die Anzeige} = la plainte policière (pas « l'annonce » publicitaire $\rightarrow$ \all{die Anzeige / die Werbung}).
\item \all{aktiv} = actif (pas « actif » au sens « élève actif » $\rightarrow$ \all{engagiert}).
\item \all{das Mobbing} = le harcèlement (pas « le mobbing » au sens anglais pur, mot intégré).
\item Particules : \all{Gib dein Passwort nicht \textbf{weiter}!} (\all{weitergeben}) ; \all{Er ruft die Polizei \textbf{an}.} (\all{anrufen}) ; \all{Sie macht einen Screenshot \textbf{von} dem Foto.} (\all{von} + Dat).
\end{itemize}
\end{enumerate}
\end{attention}

\begin{savaistu}[Pour aller plus loin — Cameroun \& Allemagne]
Au Cameroun, la loi n° 2010/012 du 21 décembre 2010 relative à la cybercriminalité et la loi n° 2016/007 du 12 juillet 2016 portant code pénal (art. 302-1 ss) répriment le harcèlement par moyens électroniques. La \all{ANTIC} (Agence Nationale des Technologies de l'Information et de la Communication) gère les incidents. En Allemagne, le \all{NetzDG} (Netzwerkdurchsetzungsgesetz) oblige les réseaux sociaux à supprimer les contenus haineux sous 24h. Comparer les deux cadres juridiques est un excellent sujet d'exposé oral (\all{mündliche Prüfung}).
\end{savaistu}

\newpage

\coursec{Exercices}

\courssub{Je vérifie mes connaissances}

\begin{exercice}[QCM : le passif]
Entoure la bonne réponse.
\begin{enumerate}
\item \all{Die Fotos \dots\ im Internet verbreitet.} \quad a) werden \quad b) wurden \quad c) sind
\item Au passif, le participe II se place : \quad a) en position 2 \quad b) à la fin de la phrase \quad c) après le sujet
\item \all{Das Opfer wird \dots\ der Polizei geschützt.} \quad a) von \quad b) mit \quad c) für
\item \all{Gestern \dots\ der Täter gefunden.} \quad a) wird \quad b) wurde \quad c) ist geworden
\end{enumerate}
\end{exercice}

\begin{exercice}[Vrai ou faux ? Justifie ta réponse]
Réponds par \all{richtig} ou \all{falsch} et justifie en une phrase en allemand.
\begin{enumerate}
\item \all{Cybermobbing findet nur in der Schule statt.}
\item \all{Das Opfer kann eine Anzeige bei der Polizei erstatten.}
\item \all{Man soll sein Passwort allen Freunden geben.}
\item \all{Beim Passiv steht das Partizip II am Satzende.}
\end{enumerate}
\end{exercice}

\begin{exercice}[Vocabulaire : à chaque mot son article et son pluriel]
Complète le tableau avec l'article et le pluriel de chaque nom.
\begin{center}
\begin{tabularx}{\linewidth}{|X|X|X|}
\hline
\rowcolor{popblueL} Nom & Article & Pluriel \\ \hline
Cybermobbing & & \\ \hline
Opfer & & \\ \hline
Täter & & \\ \hline
Anzeige & & \\ \hline
Passwort & & \\ \hline
Nachricht & & \\ \hline
\end{tabularx}
\end{center}
\end{exercice}

\begin{exercice}[Conjugaison : complète avec \all{werden} au bon temps]
Conjugue l'auxiliaire \all{werden} au présent ou au prétérit selon le contexte.
\begin{enumerate}
\item Gestern \dots\ ein Schüler gemobbt. Heute \dots\ er von einer Psychologin betreut.
\item Jeden Tag \dots\ viele Nachrichten verschickt.
\item Letzte Woche \dots\ eine Anzeige erstattet.
\item Die Schüler \dots\ von der Schule geschützt.
\end{enumerate}
\end{exercice}

\courssub{Je m'entraîne}

\begin{exercice}[Mets au passif]
Transforme les phrases suivantes au passif, d'abord au présent, puis au prétérit.
\begin{enumerate}
\item \all{Der Täter verbreitet die Fotos.}
\item \all{Die Polizei findet den Täter.}
\item \all{Die Psychologin betreut das Opfer.}
\item \all{Man löscht die Nachrichten.}
\item \all{Die Schule schützt die Schüler.}
\end{enumerate}
\end{exercice}

\begin{exercice}[Concession et opposition : trois formulations]
Relie les deux phrases de trois manières différentes, avec \all{obwohl}, \all{trotzdem} et \all{trotz}. Attention à la place du verbe et à la construction du nom.

\all{Das Opfer hatte Angst. Es ging zur Polizei.}

Fais de même avec : \all{Es gab viele Beweise. Der Täter wurde nicht gefunden.}
\end{exercice}

\begin{exercice}[Texte à trous : le passif]
Complète le texte avec la forme correcte du passif (présent ou prétérit) des verbes entre parenthèses.

\begin{textebox}[Ein Fall von Cybermobbing]
In Douala \dots\ ein 15-jähriger Schüler im Internet (mobben). Jeden Abend \dots\ böse Nachrichten an ihn (schicken). Seine Fotos \dots\ ohne seine Erlaubnis (verbreiten). Zuerst \dots\ nichts gegen den Täter (unternehmen). Schließlich \dots\ eine Anzeige bei der Polizei (erstatten). Der Täter \dots\ schnell (finden) und die Nachrichten \dots\ (löschen). Heute \dots\ der Schüler von einer Psychologin (betreuen).
\end{textebox}
\end{exercice}

\begin{exercice}[Appariements : mots et définitions]
Relie chaque mot allemand à sa définition française.
\begin{center}
\begin{tabularx}{\linewidth}{|X|X|}
\hline
\rowcolor{popblueL} Deutsch & Français \\ \hline
1. das Opfer & a) plainte déposée auprès de la police \\ \hline
2. der Täter & b) code secret pour protéger un compte \\ \hline
3. die Anzeige & c) personne qui subit le harcèlement \\ \hline
4. das Passwort & d) protéger quelqu'un contre un danger \\ \hline
5. schützen & e) personne qui commet le harcèlement \\ \hline
\end{tabularx}
\end{center}
\end{exercice}

\courssub{Je me prépare au Bac}

\begin{exercice}[Type Bac — Compréhension écrite]
Lis le texte suivant, puis réponds aux questions.

\begin{textebox}[Cybermobbing in Yaoundé : der Fall Marie]
Marie ist 16 Jahre alt und besucht ein Gymnasium in Yaoundé. Seit drei Monaten wird sie im Internet gemobbt. Jeden Abend werden beleidigende Nachrichten an sie geschickt, und ihre Fotos wurden in einer Gruppe verbreitet. Obwohl Marie große Angst hatte, sprach sie zuerst mit niemandem über ihr Problem. Sie wurde immer trauriger und ihre Noten wurden schlechter.

Eines Tages wurde ihre beste Freundin über die Nachrichten informiert. Sie überzeugte Marie, mit ihren Eltern zu sprechen. Zusammen gingen sie zur Polizei, und eine Anzeige wurde erstattet. Trotz der vielen Beweise wurde der Täter nicht sofort gefunden, denn er benutzte ein falsches Profil. Schließlich wurde er von der Polizei identifiziert : es war ein Schüler aus ihrer eigenen Klasse.

Heute wird Marie von einer Psychologin betreut, und an ihrer Schule werden die Schüler über die Gefahren des Internets informiert. „Man muss darüber sprechen“, sagt Marie. „Schweigen hilft nur dem Täter.“
\end{textebox}

\textbf{Fragen zum Text} — Réponds en allemand par des phrases complètes.
\begin{enumerate}
\item Was passiert Marie seit drei Monaten ?
\item Warum sprach sie zuerst mit niemandem ?
\item Wer half Marie, mit ihren Eltern zu sprechen ?
\item Warum wurde der Täter nicht sofort gefunden ?
\item Was wird heute an der Schule getan ?
\end{enumerate}

\textbf{Richtig oder falsch ?} — Justifie ta réponse par une citation du texte.
\begin{enumerate}
\item \all{Maries Noten blieben gut.}
\item \all{Der Täter war ein unbekannter Mann.}
\item \all{Marie findet, dass man über Cybermobbing sprechen muss.}
\end{enumerate}
\end{exercice}

\begin{exercice}[Type Bac — Thème et version]
\textbf{Thème} — Traduis en allemand (attention au passif et à la concession) :
\begin{enumerate}
\item Bien que la victime ait porté plainte, le harcèlement a continué. Les messages ont été envoyés chaque soir.
\item Les élèves doivent être protégés par l'école. Malgré la peur, la victime a parlé à ses parents.
\end{enumerate}

\textbf{Version} — Traduis en français :

\all{Trotz der vielen Beweise wurde lange keine Anzeige erstattet. Schließlich wurde der Täter von der Polizei gefunden. Obwohl das Opfer noch Angst hat, geht es wieder gern zur Schule.}
\end{exercice}

\begin{exercice}[Type Bac — Expression écrite]
Sujet : \all{Cybermobbing an unserer Schule : Was muss getan werden ?}

Rédige un texte d'environ 80 mots en allemand. Consignes obligatoires :
\begin{itemize}
\item utilise au moins deux phrases au passif ;
\item utilise au moins une phrase avec \all{obwohl}, \all{trotz} ou \all{trotzdem} ;
\item emploie au moins quatre mots du vocabulaire du thème (\all{das Opfer}, \all{der Täter}, \all{die Anzeige}, \all{schützen}, \all{das Passwort}, \all{das Cybermobbing}) ;
\item donne ton opinion personnelle dans la conclusion.
\end{itemize}
\end{exercice}

\newpage

\coursec{Corrigés des exercices}

\begin{corrige}[Exercice 1 — QCM : le passif]
\begin{enumerate}
\item \textbf{b) wurden} — \all{Die Fotos wurden im Internet verbreitet.} Le sujet \all{die Fotos} est pluriel ; le contexte évoque une action passée, d'où le prétérit de \all{werden} : \all{wurden}. « Les photos ont été diffusées sur Internet. »
\item \textbf{b) à la fin de la phrase} — Au passif, le participe II occupe la dernière place de la phrase (cadre verbal : \all{werden} conjugué en position 2, Partizip II en fin).
\item \textbf{a) von} — \all{Das Opfer wird von der Polizei geschützt.} L'agent (celui qui fait l'action) est introduit par \all{von} + datif. « La victime est protégée par la police. »
\item \textbf{b) wurde} — \all{Gestern wurde der Täter gefunden.} L'indicateur temporel \all{gestern} impose le prétérit ; le sujet \all{der Täter} est au singulier : \all{wurde}. « Hier, l'auteur a été retrouvé. »
\end{enumerate}
\textbf{Point de vigilance :} ne pas confondre le passif (\all{werden} + Partizip II) avec le Perfekt actif (\all{sein/haben} + Partizip II). \all{ist geworden} signifie « est devenu », ce qui n'a aucun sens ici. De même, \all{Die Fotos haben verbreitet} est impossible : \all{verbreiten} exige un COD (\all{Die Fotos hat der Täter verbreitet} à l'actif).
\end{corrige}

\begin{corrige}[Exercice 2 — Vrai ou faux ?]
\begin{enumerate}
\item \textbf{falsch} — \all{Cybermobbing findet überall statt, wo das Internet benutzt wird : zu Hause, in der Schule, unterwegs.} « Le harcèlement numérique a lieu partout où Internet est utilisé : à la maison, à l'école, en déplacement. »
\item \textbf{richtig} — \all{Das Opfer kann mit seinen Eltern zur Polizei gehen und eine Anzeige erstatten.} « La victime peut aller au commissariat avec ses parents et porter plainte. »
\item \textbf{falsch} — \all{Das Passwort muss geheim bleiben, sonst kann das Konto missbraucht werden.} « Le mot de passe doit rester secret, sinon le compte peut être utilisé abusivement. »
\item \textbf{richtig} — \all{Beim Passiv steht das Partizip II immer am Satzende, z. B. : „Der Täter wurde gefunden.“} « Au passif, le participe II se place toujours à la fin de la phrase. »
\end{enumerate}
\textbf{Point de vigilance :} la justification doit être une phrase complète avec le verbe conjugué en position 2 ; dans la subordonnée en \all{wo} (phrase 1), le verbe conjugué \all{wird} est rejeté à la fin.
\end{corrige}

\begin{corrige}[Exercice 3 — Vocabulaire : article et pluriel]
\begin{center}
\begin{tabularx}{\linewidth}{|X|X|X|}
\hline
\rowcolor{popblueL} Nom & Article & Pluriel \\ \hline
Cybermobbing & das & (ohne Plural) \\ \hline
Opfer & das & die Opfer \\ \hline
Täter & der & die Täter \\ \hline
Anzeige & die & die Anzeigen \\ \hline
Passwort & das & die Passwörter \\ \hline
Nachricht & die & die Nachrichten \\ \hline
\end{tabularx}
\end{center}
\textbf{Remarques :}
\begin{itemize}
\item \all{das Cybermobbing} est un nom abstrait : il n'a pas de pluriel en usage courant.
\item \all{das Opfer} et \all{der Täter} ont un pluriel identique au singulier : seul l'article change (\all{die Opfer}, \all{die Täter}).
\item \all{das Passwort} prend un Umlaut au pluriel : \all{die Passwörter} (comme \all{das Wort} $\rightarrow$ \all{die Wörter}).
\end{itemize}
\textbf{Point de vigilance :} apprends toujours un nom allemand avec son article et son pluriel ; le genre ne se devine pas à partir du français (\emph{la} victime mais \all{das Opfer} !). Conséquence grammaticale : \all{das Opfer} se reprend par le pronom neutre \all{es}, même si la victime est une fille.
\end{corrige}

\begin{corrige}[Exercice 4 — Conjugaison : \all{werden} au bon temps]
\begin{enumerate}
\item Gestern \textbf{wurde} ein Schüler gemobbt. Heute \textbf{wird} er von einer Psychologin betreut. (\all{gestern} $\rightarrow$ prétérit singulier ; \all{heute} $\rightarrow$ présent singulier.)
\item Jeden Tag \textbf{werden} viele Nachrichten verschickt. (habitude au présent ; sujet pluriel \all{viele Nachrichten}.)
\item Letzte Woche \textbf{wurde} eine Anzeige erstattet. (passé $\rightarrow$ prétérit ; sujet singulier \all{eine Anzeige}.)
\item Die Schüler \textbf{werden} von der Schule geschützt. (présent ; sujet pluriel \all{die Schüler}.)
\end{enumerate}
\textbf{Rappel de conjugaison :}
\begin{center}
\begin{tabular}{|l|l|l|}
\hline
\rowcolor{popblueL} Personne & Präsens & Präteritum \\ \hline
ich & werde & wurde \\ \hline
du & wirst & wurdest \\ \hline
er/sie/es & wird & wurde \\ \hline
wir & werden & wurden \\ \hline
ihr & werdet & wurdet \\ \hline
sie/Sie & werden & wurden \\ \hline
\end{tabular}
\end{center}
\textbf{Point de vigilance :} au prétérit, \all{werden} ne prend jamais d'Umlaut : on écrit \all{wurde}, \all{wurden} (prononcé « vour-de », « vour-den »), et non une forme avec \emph{ü}.
\end{corrige}

\begin{corrige}[Exercice 5 — Mets au passif]
\begin{enumerate}
\item \all{Die Fotos werden (von dem Täter) verbreitet.} / \all{Die Fotos wurden (von dem Täter) verbreitet.} « Les photos sont / ont été diffusées par l'auteur. »
\item \all{Der Täter wird von der Polizei gefunden.} / \all{Der Täter wurde von der Polizei gefunden.} « L'auteur est / a été retrouvé par la police. »
\item \all{Das Opfer wird von der Psychologin betreut.} / \all{Das Opfer wurde von der Psychologin betreut.} « La victime est / a été suivie par la psychologue. »
\item \all{Die Nachrichten werden gelöscht.} / \all{Die Nachrichten wurden gelöscht.} « Les messages sont / ont été effacés. » (Avec \all{man}, le complément d'agent disparaît au passif.)
\item \all{Die Schüler werden von der Schule geschützt.} / \all{Die Schüler wurden von der Schule geschützt.} « Les élèves sont / ont été protégés par l'école. »
\end{enumerate}
\textbf{Méthode :} (1) le COD de la phrase active devient le sujet au nominatif ; (2) le verbe se met à la forme \all{werden} + Partizip II, \all{werden} s'accordant avec le nouveau sujet ; (3) l'ancien sujet devient complément d'agent avec \all{von} + datif (ou disparaît si c'est \all{man}).
\textbf{Point de vigilance :} \all{der Täter} au datif devient \all{von dem Täter} ; \all{den Täter} (accusatif) de la phrase 2 redevient \all{der Täter} (nominatif) au passif. Attention aux participes sans \all{ge-} : \all{verbreiten} $\rightarrow$ \all{verbreitet}, \all{betreuen} $\rightarrow$ \all{betreut}.
\end{corrige}

\begin{corrige}[Exercice 6 — Concession et opposition : trois formulations]
\textbf{Premier couple de phrases :}
\begin{enumerate}
\item \all{Obwohl das Opfer Angst hatte, ging es zur Polizei.} (subordonnée avec \all{obwohl} : verbe conjugué rejeté à la fin ; la subordonnée occupant la position 1, le verbe de la principale vient juste après la virgule, avant le sujet \all{es}.) « Bien que la victime ait eu peur, elle est allée à la police. »
\item \all{Das Opfer hatte Angst. Trotzdem ging es zur Polizei.} (\all{trotzdem} est un connecteur en position 1 : inversion sujet-verbe.) « La victime avait peur. Pourtant, elle est allée à la police. »
\item \all{Trotz seiner Angst ging das Opfer zur Polizei.} (\all{trotz} + génitif du nom : \all{die Angst} $\rightarrow$ \all{trotz seiner Angst}.) « Malgré sa peur, la victime est allée à la police. »
\end{enumerate}
\textbf{Deuxième couple de phrases :}
\begin{enumerate}
\item \all{Obwohl es viele Beweise gab, wurde der Täter nicht gefunden.} « Bien qu'il y ait eu de nombreuses preuves, l'auteur n'a pas été retrouvé. »
\item \all{Es gab viele Beweise. Trotzdem wurde der Täter nicht gefunden.} « Il y avait de nombreuses preuves. Pourtant, l'auteur n'a pas été retrouvé. »
\item \all{Trotz der vielen Beweise wurde der Täter nicht gefunden.} (\all{trotz} + génitif : \all{die vielen Beweise} $\rightarrow$ \all{trotz der vielen Beweise}.) « Malgré les nombreuses preuves, l'auteur n'a pas été retrouvé. »
\end{enumerate}
\textbf{Point de vigilance :} avec \all{obwohl}, le verbe va à la fin de la subordonnée ; avec \all{trotzdem}, on inverse sujet et verbe ; avec \all{trotz}, il faut un \textbf{nom} au génitif, jamais une phrase verbale. Ne pas écrire \all{trotz die Angst} : la préposition exige le génitif (\all{trotz der Angst}).
\end{corrige}

\begin{corrige}[Exercice 7 — Texte à trous : le passif]
Texte complété :

In Douala \textbf{wurde} ein 15-jähriger Schüler im Internet \textbf{gemobbt}. Jeden Abend \textbf{wurden} böse Nachrichten an ihn \textbf{geschickt}. Seine Fotos \textbf{wurden} ohne seine Erlaubnis \textbf{verbreitet}. Zuerst \textbf{wurde} nichts gegen den Täter \textbf{unternommen}. Schließlich \textbf{wurde} eine Anzeige bei der Polizei \textbf{erstattet}. Der Täter \textbf{wurde} schnell \textbf{gefunden} und die Nachrichten \textbf{wurden gelöscht}. Heute \textbf{wird} der Schüler von einer Psychologin \textbf{betreut}.

Traduction : « À Douala, un élève de 15 ans a été harcelé sur Internet. Chaque soir, des messages méchants lui étaient envoyés. Ses photos ont été diffusées sans son autorisation. D'abord, rien n'a été entrepris contre l'auteur. Finalement, une plainte a été déposée auprès de la police. L'auteur a été vite retrouvé et les messages ont été effacés. Aujourd'hui, l'élève est suivi par une psychologue. »

\textbf{Justifications :} tout le récit est au passé $\rightarrow$ prétérit du passif (\all{wurde(n)} + Partizip II) ; seule la dernière phrase (\all{heute}) passe au présent (\all{wird betreut}). Attention aux participes des verbes à préfixe inséparable : \all{verbreiten} $\rightarrow$ \all{verbreitet}, \all{erstatten} $\rightarrow$ \all{erstattet}, \all{betreuen} $\rightarrow$ \all{betreut} (pas de \all{ge-} !). \all{unternehmen} est un verbe fort : \all{unternommen}. Enfin, \all{ein 15-jähriger Schüler} est bien au nominatif singulier : c'est le sujet du passif.
\end{corrige}

\begin{corrige}[Exercice 8 — Appariements : mots et définitions]
\begin{center}
\begin{tabularx}{\linewidth}{|X|X|}
\hline
\rowcolor{popblueL} Deutsch & Français \\ \hline
1. das Opfer & c) personne qui subit le harcèlement \\ \hline
2. der Täter & e) personne qui commet le harcèlement \\ \hline
3. die Anzeige & a) plainte déposée auprès de la police \\ \hline
4. das Passwort & b) code secret pour protéger un compte \\ \hline
5. schützen & d) protéger quelqu'un contre un danger \\ \hline
\end{tabularx}
\end{center}
\textbf{Expressions utiles à retenir :} \all{eine Anzeige erstatten} « porter plainte » ; \all{jemanden vor etwas schützen} « protéger quelqu'un de quelque chose » (avec \all{vor} + datif) ; \all{Opfer von Cybermobbing werden} « devenir victime de harcèlement numérique ».
\textbf{Point de vigilance :} \all{der Täter} est un faux ami partiel : il désigne l'\emph{auteur} d'un acte (délit, harcèlement), pas un « tuteur » ; la victime est \all{das Opfer}, toujours neutre.
\end{corrige}

\begin{corrige}[Exercice 9 — Type Bac : compréhension écrite]
\textbf{Fragen zum Text} (réponses modèles, en phrases complètes) :
\begin{enumerate}
\item \all{Seit drei Monaten wird Marie im Internet gemobbt : sie bekommt beleidigende Nachrichten, und ihre Fotos wurden verbreitet.} « Depuis trois mois, Marie est harcelée sur Internet : elle reçoit des messages insultants et ses photos ont été diffusées. »
\item \all{Sie sprach zuerst mit niemandem, weil sie große Angst hatte.} « Elle n'a d'abord parlé à personne parce qu'elle avait très peur. »
\item \all{Ihre beste Freundin half ihr, mit ihren Eltern zu sprechen.} « Sa meilleure amie l'a aidée à parler à ses parents. »
\item \all{Der Täter wurde nicht sofort gefunden, weil er ein falsches Profil benutzte.} « L'auteur n'a pas été retrouvé immédiatement parce qu'il utilisait un faux profil. »
\item \all{Heute werden die Schüler an ihrer Schule über die Gefahren des Internets informiert.} « Aujourd'hui, les élèves de son école sont informés des dangers d'Internet. »
\end{enumerate}

\textbf{Richtig oder falsch ?}
\begin{enumerate}
\item \textbf{falsch} — Le texte dit : \all{„ihre Noten wurden schlechter“} (ses notes sont devenues moins bonnes).
\item \textbf{falsch} — Le texte dit : \all{„es war ein Schüler aus ihrer eigenen Klasse“} (c'était un élève de sa propre classe).
\item \textbf{richtig} — Le texte dit : \all{„Man muss darüber sprechen“, sagt Marie. „Schweigen hilft nur dem Täter.“} (« Il faut en parler, dit Marie. Le silence n'aide que l'auteur. »)
\end{enumerate}

\textbf{Barème indicatif (sur 10) :} questions de compréhension : 5 $\times$ 1 pt (contenu 0,5 + correction linguistique 0,5) ; richtig/falsch : 3 $\times$ 1 pt (réponse 0,5 + citation justificative 0,5) ; qualité globale de la langue : 2 pts.
\textbf{Points de vigilance :} répondre par des phrases complètes (verbe en position 2) ; reprendre le temps de la question (\all{sprach} au prétérit car le texte raconte au passé) ; dans les subordonnées en \all{weil}, verbe conjugué à la fin (\all{weil sie große Angst hatte}) ; pour la justification, citer le texte entre guillemets allemands „…“.
\end{corrige}

\begin{corrige}[Exercice 10 — Type Bac : thème et version]
\textbf{Thème} (traductions proposées) :
\begin{enumerate}
\item \all{Obwohl das Opfer eine Anzeige erstattet hat, ging das Mobbing weiter. Die Nachrichten wurden jeden Abend geschickt.} (Variantes acceptées : \all{... erstattete, ...} au prétérit ; \all{... setzte sich das Mobbing fort.})
\item \all{Die Schüler müssen von der Schule geschützt werden. Trotz der Angst sprach das Opfer mit seinen Eltern.} (Variante : \all{Obwohl es Angst hatte, sprach das Opfer mit seinen Eltern.})
\end{enumerate}
\textbf{Points de vigilance (thème) :}
\begin{itemize}
\item « porter plainte » = \all{eine Anzeige erstatten} (et non \emph{machen}) ;
\item « les messages ont été envoyés » = passif au prétérit : \all{wurden geschickt} (sujet pluriel !) ;
\item « doivent être protégés » = \all{müssen} + passif à l'infinitif : \all{geschützt werden} (double infinitif en fin de phrase, \all{werden} en tout dernier) ;
\item « malgré la peur » = \all{trotz der Angst} (génitif) ;
\item après \all{obwohl}, le verbe conjugué va en fin de subordonnée, et la principale qui suit commence par le verbe (\all{..., ging das Mobbing weiter}).
\end{itemize}

\textbf{Version} (traduction proposée) : « Malgré les nombreuses preuves, aucune plainte n'a été déposée pendant longtemps. Finalement, l'auteur a été retrouvé par la police. Bien que la victime ait encore peur, elle retourne volontiers à l'école. »

\textbf{Points de vigilance (version) :} \all{es geht wieder gern zur Schule} : \all{es} reprend \all{das Opfer} (neutre !) — traduire par « elle » en français ; \all{lange keine Anzeige erstattet} = « aucune plainte ... pendant longtemps » ; \all{wurde ... gefunden} = passif au prétérit, à rendre par « a été retrouvé ».

\textbf{Barème indicatif (sur 10) :} thème 5 pts (2,5 pts par phrase : passif 1 pt, concession 0,5 pt, lexique et ordre des mots 1 pt) ; version 5 pts (fidélité 3 pts, qualité du français 2 pts).
\end{corrige}

\begin{corrige}[Exercice 11 — Type Bac : expression écrite]
\textbf{Proposition de rédaction modèle (environ 90 mots) :}

\all{Cybermobbing an unserer Schule : Was muss getan werden ?}

\all{Cybermobbing ist auch an unserer Schule ein Problem. Viele Schüler werden im Internet beleidigt, und ihre Fotos werden ohne Erlaubnis verbreitet. Obwohl die Opfer große Angst haben, schweigen sie oft. Das muss geändert werden : die Schüler müssen über die Gefahren des Internets informiert werden, und die Opfer müssen von der Schule geschützt werden. Wenn jemand gemobbt wird, soll eine Anzeige bei der Polizei erstattet werden, damit der Täter gefunden wird. Meiner Meinung nach hilft Schweigen nur dem Täter : man muss sprechen und handeln.}

Traduction : « Le harcèlement numérique est aussi un problème dans notre école. Beaucoup d'élèves sont insultés sur Internet et leurs photos sont diffusées sans autorisation. Bien que les victimes aient très peur, elles se taisent souvent. Cela doit être changé : les élèves doivent être informés des dangers d'Internet et les victimes doivent être protégées par l'école. Quand quelqu'un est harcelé, une plainte devrait être déposée auprès de la police afin que l'auteur soit retrouvé. À mon avis, le silence n'aide que l'auteur : il faut parler et agir. »

\textbf{Vérification des consignes :}
\begin{itemize}
\item passif (au moins 2 phrases) : \all{werden beleidigt}, \all{werden verbreitet}, \all{muss geändert werden}, \all{müssen informiert werden}, \all{müssen geschützt werden}, \all{soll erstattet werden}, \all{wird gefunden} ;
\item concession : \all{Obwohl die Opfer große Angst haben, schweigen sie oft.} ;
\item vocabulaire du thème (au moins 4 mots) : \all{das Cybermobbing}, \all{die Opfer}, \all{der Täter}, \all{die Anzeige}, \all{schützen} ;
\item opinion personnelle en conclusion : \all{Meiner Meinung nach ...}.
\end{itemize}

\textbf{Barème indicatif (sur 10) :} respect des consignes de contenu 4 pts (passif 1, concession 1, vocabulaire 1, opinion 1) ; correction grammaticale (ordre des mots, conjugaisons, cas) 3 pts ; richesse du lexique et orthographe (majuscules, Umlaute) 2 pts ; cohérence et longueur (environ 80 mots) 1 pt.

\textbf{Points de vigilance :} verbe conjugué toujours en position 2 ; après \all{obwohl}, verbe à la fin ; avec modal + passif, double infinitif \all{geschützt werden} en fin de phrase ; majuscule à tous les noms (\all{Schweigen}, \all{Erlaubnis}, \all{Gefahren}) ; ne pas écrire \all{die Opfer haben Angst gehabt} si le contexte est au présent.
\end{corrige}

\newpage

\coursec{Fiche bilan}

\begin{popfigure}
\begin{tikzpicture}[mindmap, grow cyclic, every node/.style={concept, text=white, font=\small},
  root concept/.append style={concept color=popblue, font=\bfseries},
  level 1/.append style={level distance=3.4cm, sibling angle=60},
  level 1 concept/.append style={font=\small},
  level 2/.append style={level distance=2.2cm, sibling angle=45, font=\scriptsize}]
\node[root concept, align=center]{Cybermobbing\\ (AL18)}
  child[concept color=poporange]{ node {Texte : ein Fall aus Berlin}
    child { node {Opfer, Täter, Schule} }
    child { node {Anzeige, Polizei} }
  }
  child[concept color=popgreen]{ node {Lexique}
    child { node {das Opfer, der Täter} }
    child { node {die Anzeige, das Passwort} }
    child { node {schützen, beleidigen} }
  }
  child[concept color=popblue]{ node {Passif}
    child { node {Präsens : wird + Part. II} }
    child { node {Präteritum : wurde + Part. II} }
  }
  child[concept color=popgold]{ node {Concession}
    child { node {obwohl + verbe final} }
    child { node {trotz + Génitif} }
  }
  child[concept color=poppurple]{ node {Opposition}
    child { node {aber, sondern} }
    child { node {während + verbe final} }
  }
  child[concept color=poppink]{ node {Méthode}
    child { node {lire -- repérer -- résumer} }
    child { node {citer le texte} }
  };
\end{tikzpicture}
\legende{Carte mentale de la leçon : le harcèlement numérique (das Cybermobbing)}
\end{popfigure}

\begin{bilan}
\textbf{1. Lexique clé (à connaître avec article et pluriel)}
\begin{itemize}
  \item \all{das Cybermobbing} (Sg.) : le cyberharcèlement ; \all{das Mobbing} : le harcèlement.
  \item \all{das Opfer} (--) : la victime ; \all{der Täter} (--) : l'auteur, le coupable.
  \item \all{die Anzeige} (-n) : la plainte ; \all{Anzeige erstatten} : porter plainte.
  \item \all{das Passwort} (\"-er) : le mot de passe ; \all{das Profil} (-e) : le profil.
  \item \all{die Nachricht} (-en) : le message ; \all{das soziale Netzwerk} (-e) : le réseau social.
  \item \all{schützen} : protéger ; \all{beleidigen} : insulter ; \all{veröffentlichen} : publier.
  \item \all{jmdm. helfen} (+ datif) : aider qqn ; \all{jmdn. anzeigen} : dénoncer qqn.
\end{itemize}

\textbf{2. Grammaire à connaître par c\oe ur}
\begin{center}
\begin{tabularx}{\linewidth}{|X|X|X|}
\hline
\rowcolor{popblueL} Point & Règle & Exemple \\
\hline
Passif Präsens & \all{werden} (conjugué) + Partizip II final & \all{Er wird beleidigt.} « Il est insulté. » \\
\hline
Passif Präteritum & \all{wurden} + Partizip II final & \all{Sie wurde bedroht.} « Elle a été menacée. » \\
\hline
Agent au passif & \all{von} + datif & \all{... von einem Mitschüler} \\
\hline
Concession & \all{obwohl} + verbe en fin ; \all{trotz} + génitif & \all{Obwohl sie Angst hatte, ging sie zur Polizei.} \\
\hline
Opposition & \all{aber} (verbe 2\up{e}) ; \all{sondern} (après négation) ; \all{während} + verbe final & \all{Er schwieg, aber sie handelte.} \\
\hline
\end{tabularx}
\end{center}

\textbf{3. Méthodes express}
\begin{itemize}
  \item Compréhension : lire une première fois pour le sens global, repérer \all{wer -- was -- wo -- warum}, puis répondre en phrases complètes en citant le texte.
  \item Actif $\rightarrow$ passif : l'objet accusatif devient sujet ; le sujet devient \all{von} + datif (ou disparaît) ; \all{werden} se conjugue, le participe II va en fin de phrase.
  \item Production courte : 5 à 8 phrases, au moins une au passif et une avec \all{obwohl} ou \all{trotz}.
\end{itemize}
\end{bilan}

\begin{aretenir}[Checklist avant le Bac]
\begin{itemize}
  \item $\square$ Je connais le lexique du cyberharcèlement avec articles et pluriels (\all{das Opfer}, \all{der Täter}, \all{die Anzeige}, \all{das Passwort}).
  \item $\square$ Je sais former le passif au présent : \all{wird/werden} + Partizip II.
  \item $\square$ Je sais former le passif au prétérit : \all{wurde/wurden} + Partizip II.
  \item $\square$ Je place correctement le participe II en fin de phrase au passif.
  \item $\square$ J'introduis l'agent avec \all{von} + datif (et non \all{durch}, réservé au moyen).
  \item $\square$ Je mets le verbe conjugué en fin de proposition après \all{obwohl} et \all{während}.
  \item $\square$ Je distingue \all{aber} (opposition simple) et \all{sondern} (après une négation).
  \item $\square$ J'emploie \all{trotz} + génitif : \all{trotz der Drohungen}.
  \item $\square$ Je réponds aux questions par des phrases complètes en allemand, avec majuscules aux noms.
  \item $\square$ Je sais résumer le texte en 3 à 5 phrases au Präsens ou au Perfekt.
  \item $\square$ Je vérifie l'ordre des mots : verbe conjugué en 2\up{e} position dans la phrase principale.
  \item $\square$ Je relis ma production : accords, cas après prépositions, orthographe des Umlaute.
\end{itemize}
\end{aretenir}

\findecours

\end{document}
